\input{preamble}

% OK, commencer ici.
%
\begin{document}

\title{Algèbre différentielle graduée}


\maketitle

\phantomsection
\label{section-phantom}

\tableofcontents

\section{Introduction}
\label{section-introduction}

\noindent
Dans ce chapitre, nous traitons des algèbres différentielles graduées, des modules,
des catégories, etc. Une référence de base est \cite{Keller-Deriving}.
Un article de synthèse est \cite{Keller-survey}.

\medskip\noindent
Comme la longueur des développements ne nous préoccupe pas dans le projet Champs,
nous développons d'abord la matière dans le cadre des catégories de modules
différentiels gradués. Nous reprenons ensuite les constructions dans le cadre des
modules différentiels gradués sur des catégories différentielles graduées.



\section{Conventions}
\label{section-conventions}

\noindent
Dans ce chapitre, nous conservons la convention selon laquelle {\it anneau} signifie
anneau commutatif avec unité $1$. Si $R$ est un anneau, une {\it $R$-algèbre $A$}
sera un $R$-module $A$ muni d'une application $R$-bilinéaire $A \times A \to A$
(la multiplication), telle que celle-ci soit associative et possède une unité.
Autrement dit, ce sont des $R$-algèbres associatives avec unité
dont le morphisme structural $R \to A$ prend ses valeurs dans le centre de $A$.

\medskip\noindent
{\bf Règles de signe.} Dans ce chapitre, nous travaillerons avec des algèbres graduées
et des modules gradués, souvent munis de différentielles. Les règles de signe des
complexes sous-jacents seront toujours celles introduites dans
Compléments d'algèbre, section \ref{more-algebra-section-sign-rules} (ou des règles compatibles avec celles-ci).
Il en résultera parfois que la structure multiplicative se trouve tordue d'une manière
inattendue, surtout lorsque l'on considère des modules à gauche
ou le lien entre modules à gauche et modules à droite.





\section{Algèbres différentielles graduées}
\label{section-dga}


\noindent
Nous donnons seulement les définitions.

\begin{definition}
\label{definition-dga}
Soit $R$ un anneau commutatif. Une {\it algèbre différentielle graduée sur $R$}
est l'une ou l'autre des données suivantes :
\begin{enumerate}
\item un complexe de chaînes $A_\bullet$ de $R$-modules muni d'applications
$R$-bilinéaires $A_n \times A_m \to A_{n + m}$,
$(a, b) \mapsto ab$ telles que
$$
\text{d}_{n + m}(ab) = \text{d}_n(a)b + (-1)^n a\text{d}_m(b)
$$
et tel que $\bigoplus A_n$ soit une $R$-algèbre associative
avec unité, ou
\item un complexe de cochaînes $A^\bullet$ de $R$-modules muni d'applications
$R$-bilinéaires $A^n \times A^m \to A^{n + m}$, $(a, b) \mapsto ab$
telles que
$$
\text{d}^{n + m}(ab) = \text{d}^n(a)b + (-1)^n a\text{d}^m(b)
$$
et tel que $\bigoplus A^n$ soit une $R$-algèbre associative avec unité.
\end{enumerate}
\end{definition}

\noindent
Nous écrivons souvent simplement $A = \bigoplus A_n$ ou $A = \bigoplus A^n$ et
voyons cette somme comme une $R$-algèbre associative avec unité, munie d'une
graduation par $\mathbf{Z}$ et d'un opérateur $R$-linéaire $\text{d}$ dont le carré
est nul et qui satisfait la règle de Leibniz décrite ci-dessus. Dans ce cas,
nous disons souvent : « Soit $(A, \text{d})$ une algèbre différentielle graduée. »

\medskip\noindent
La règle de Leibniz reliant les différentielles et la multiplication dans une
$R$-algèbre différentielle graduée $A$ signifie exactement que l'application de multiplication
définit un morphisme de complexes de cochaînes
$$
\text{Tot}(A^\bullet \otimes_R A^\bullet) \to A^\bullet
$$
Ici, $A^\bullet$ désigne le complexe de cochaînes sous-jacent à $A$.

\begin{definition}
\label{definition-homomorphism-dga}
Un {\it homomorphisme d'algèbres différentielles graduées}
$f : (A, \text{d}) \to (B, \text{d})$ est un homomorphisme d'algèbres $f : A \to B$
compatible avec les graduations et avec $\text{d}$.
\end{definition}

\begin{definition}
\label{definition-cdga}
Une algèbre différentielle graduée $(A, \text{d})$ est {\it commutative} si
$ab = (-1)^{nm}ba$ pour $a$ de degré $n$ et $b$ de degré $m$.
Nous disons que $A$ est {\it strictement commutative} si, en outre, $a^2 = 0$
lorsque $\deg(a)$ est impair.
\end{definition}

\noindent
La définition suivante a un sens en général, mais n'est peut-être
« correcte » que pour le produit tensoriel d'algèbres différentielles graduées
commutatives.

\begin{definition}
\label{definition-tensor-product}
Soit $R$ un anneau.
Soient $(A, \text{d})$, $(B, \text{d})$ des algèbres différentielles graduées sur $R$.
Le {\it produit tensoriel} de $A$ et $B$ comme algèbres différentielles graduées
est l'algèbre $A \otimes_R B$ dont la multiplication est définie par
$$
(a \otimes b)(a' \otimes b') = (-1)^{\deg(a')\deg(b)} aa' \otimes bb'
$$
munie de la différentielle $\text{d}$ définie par la règle
$\text{d}(a \otimes b) = \text{d}(a) \otimes b + (-1)^m a \otimes \text{d}(b)$
où $m = \deg(a)$.
\end{definition}

\begin{lemma}
\label{lemma-total-complex-tensor-product}
Soit $R$ un anneau.
Soient $(A, \text{d})$, $(B, \text{d})$ des algèbres différentielles graduées sur $R$.
Désignons par $A^\bullet$, $B^\bullet$ les complexes de cochaînes sous-jacents.
En tant que complexes de cochaînes de $R$-modules, nous avons
$$
(A \otimes_R B)^\bullet = \text{Tot}(A^\bullet \otimes_R B^\bullet).
$$
\end{lemma}

\begin{proof}
Rappelons que la différentielle du complexe total est donnée par
$\text{d}_1^{p, q} + (-1)^p \text{d}_2^{p, q}$ sur $A^p \otimes_R B^q$.
Or c'est exactement la règle qui définit la différentielle
sur $A \otimes_R B$ dans la
définition \ref{definition-tensor-product}.
\end{proof}






\section{Modules différentiels gradués}
\label{section-modules}

\noindent
Par défaut dans ce chapitre, nous considérons les modules à
droite ; nous discutons des modules à gauche dans la section \ref{section-left-modules}.

\begin{definition}
\label{definition-dgm}
Soit $R$ un anneau.
Soit $(A, \text{d})$ une algèbre différentielle graduée sur $R$.
Un {\it module différentiel gradué} $M$ à droite sur $A$ est un $A$-module à droite
$M$ qui possède une graduation $M = \bigoplus M^n$ et une différentielle
$\text{d}$ tels que $M^n A^m \subset M^{n + m}$, tels que
$\text{d}(M^n) \subset M^{n + 1}$, et tels que
$$
\text{d}(ma) = \text{d}(m)a + (-1)^n m\text{d}(a)
$$
pour $a \in A$ et $m \in M^n$. Un
{\it homomorphisme de modules différentiels gradués} $f : M \to N$ est une
application de $A$-modules compatible avec les graduations et les
différentielles. La catégorie des $A$-modules différentiels gradués (droits) est
notée $\text{Mod}_{(A, \text{d})}$.
\end{definition}

\noindent
Notons que nous pouvons considérer $M$ comme un complexe de cochaînes
$M^\bullet$ de $R$-modules (droits). À savoir, pour $r \in R$ nous avons $\text{d}(r)
= 0$ et $r$ s'envoie sur un élément de degré $0$ de $A$,
donc $\text{d}(mr) = \text{d}(m)r$.

\medskip\noindent
La règle de Leibniz reliant les différentielles et la multiplication sur un $R$-module
différentiel gradué $M$ sur une $R$-algèbre différentielle graduée $A$ signifie
exactement que l'application de multiplication définit un morphisme de complexes de cochaînes
$$
\text{Tot}(M^\bullet \otimes_R A^\bullet) \to M^\bullet
$$
Ici $A^\bullet$ et $M^\bullet$ désignent les complexes de cochaînes
sous-jacents de $A$ et $M$.

\begin{lemma}
\label{lemma-dgm-abelian}
Soit $(A, d)$ une algèbre différentielle graduée. La catégorie
$\text{Mod}_{(A, \text{d})}$ est abélienne et possède des limites et colimites arbitraires.
\end{lemma}

\begin{proof}
Les noyaux et conoyaux commutent avec la prise des $A$-modules
sous-jacents. Il en va de même pour les sommes directes et les colimites. En d'autres
termes, ces opérations dans $\text{Mod}_{(A, \text{d})}$ commutent avec le foncteur
d'oubli vers la catégorie des $A$-modules. Ce n'est pas le cas pour les
produits et limites. À savoir, si $N_i$, $i \in I$ est
une famille de $A$-modules différentiels gradués, alors le produit
$\prod N_i$ dans $\text{Mod}_{(A, \text{d})}$ est donné en posant $(\prod N_i)^n = \prod
N_i^n$ et $\prod N_i = \bigoplus_n (\prod N_i)^n$. Ainsi, nous voyons que le produit
commute avec le foncteur d'oubli vers la catégorie des $A$-modules gradués. Une
catégorie avec produits et égalisateurs possède des limites, voir
Catégories, lemme \ref{categories-lemma-limits-products-equalizers}.
\end{proof}

\noindent
Ainsi, si $(A, \text{d})$ est une algèbre
différentielle graduée sur $R$, alors il existe un foncteur exact
$$
\text{Mod}_{(A, \text{d})} \longrightarrow \text{Comp}(R)
$$
de catégories abéliennes. Pour un module différentiel gradué $M$,
les groupes de cohomologie $H^n(M)$ sont définis comme la cohomologie
du complexe correspondant de $R$-modules. Par conséquent, une suite
exacte courte $0 \to K \to L \to M \to 0$ de modules différentiels gradués
donne lieu à une longue suite exacte
\begin{equation}
\label{equation-les}
H^n(K) \to H^n(L) \to H^n(M) \to H^{n + 1}(K)
\end{equation}
de modules de cohomologie,
voir Homologie, lemme \ref{homology-lemma-long-exact-sequence-cochain}.

\medskip\noindent
De plus, à partir de maintenant, nous empruntons toute la
terminologie utilisée pour les complexes de modules. Par exemple, nous disons
qu'un $A$-module différentiel gradué $M$ est {\it acyclique} si
$H^k(M) = 0$ pour tout $k \in \mathbf{Z}$. Nous disons qu'un homomorphisme
$M \to N$ de modules différentiels gradués sur $A$ est un {\it
quasi-isomorphisme} s'il induit des isomorphismes $H^k(M) \to H^k(N)$ pour tout $k \in
\mathbf{Z}$. Et ainsi de suite.

\begin{definition}
\label{definition-shift}
Soit $(A, \text{d})$ une algèbre différentielle graduée.
Soit $M$ un module différentiel gradué dont le complexe sous-jacent de
$R$-modules est $M^\bullet$. Pour tout $k \in \mathbf{Z}$, nous
définissons le {\it module décalé de $k$} $M[k]$ comme suit
\begin{enumerate}
\item le complexe de $R$-modules sous-jacent à $M[k]$ est $M^\bullet[k]$,
c'est-à-dire que nous avons $M[k]^n = M^{n + k}$
et $\text{d}_{M[k]} = (-1)^k\text{d}_M$ et
\item la structure de $A$-module est donnée par la multiplication
$$
(M[k])^n \times A^m \longrightarrow (M[k])^{n + m}
$$
est égale à la multiplication donnée $M^{n + k} \times A^m \to M^{n + k + m}$.
\end{enumerate}
Pour un morphisme $f : M \to N$ de $A$-modules différentiels gradués,
nous notons $f[k] : M[k] \to N[k]$ l'application égale à $f$ sur les
$A$-modules sous-jacents. Cela définit un
foncteur $[k] : \text{Mod}_{(A, \text{d})} \to \text{Mod}_{(A, \text{d})}$.
\end{definition}

\noindent
Vérifions qu'avec ce choix la règle de Leibniz est satisfaite.
Soient $x \in M[k]^n = M^{n + k}$ et $a \in A^m$. En notant
$\cdot_{M[k]}$ le produit dans $M[k]$, on obtient
\begin{align*}
\text{d}_{M[k]}(x \cdot_{M[k]} a)
& =
(-1)^k \text{d}_M(xa) \\
& =
(-1)^k \text{d}_M(x) a + (-1)^{k + n + k} x \text{d}(a) \\ & =
\text{d}_{M[k]}(x)
a + (-1)^n x \text{d}(a) \\ & =
\text{d}_{M[k]}(x)
\cdot_{M[k]} a + (-1)^n x \cdot_{M[k]} \text{d}(a)
\end{align*}
C'est ce que nous voulons puisque $x$ a un degré $n$ en tant qu'élément homogène de $M[k]$. Nous
observons également qu'avec ces choix, nous pouvons considérer
l'application de multiplication comme un morphisme de complexes
$$
\text{Tot}(M^\bullet[k] \otimes _R A^\bullet) \to
\text{Tot}(M^\bullet \otimes _R A^\bullet)[k] \to
M^\bullet[k]
$$
où la première flèche est
celle de Compléments d'algèbre, section
\ref{more-algebra-section-sign-rules} (\ref{more-algebra-item-shift-tensor}) qui dans ce
cas n'implique pas de signe. (En fait, nous aurions pu déduire que la
règle de Leibniz vaut à partir de cette observation.)

\medskip\noindent
Les remarques dans Homologie, section \ref{homology-section-homotopy-shift}
s'appliquent. En particulier, nous identifierons les groupes de cohomologie de tous les
décalages $M[k]$ sans l'intervention de signes.

\medskip\noindent
À ce stade, nous avons suffisamment de structure pour parler des {\it triangles},
voir Catégories dérivées, définition \ref{derived-definition-triangle}. En fait,
notre prochain objectif est de développer suffisamment de théorie pour pouvoir
énoncer et prouver que la catégorie homotopique des modules différentiels gradués
est une catégorie triangulée. Tout d'abord, nous définissons la catégorie homotopique.






\section{La catégorie homotopique}
\label{section-homotopy}

\noindent
Nos homotopies tiennent compte de la structure de module $A$ et de la
graduation, mais pas de la différentielle (bien sûr).

\begin{definition}
\label{definition-homotopy}
Soit $(A, \text{d})$ une algèbre différentielle graduée. Soient
$f, g : M \to N$ des homomorphismes de $A$-modules différentiels gradués. Une
{\it homotopie entre $f$ et $g$} est une application de $A$-modules $h : M \to N$
telle que
\begin{enumerate}
\item $h(M^n) \subset N^{n - 1}$ pour tout $n$, et
\item $f(x) - g(x) = \text{d}_N(h(x)) + h(\text{d}_M(x))$ pour
tout $x \in M$.
\end{enumerate}
Si une homotopie existe, alors nous disons que $f$ et $g$ sont {\it homotopes}.
\end{definition}

\noindent
Ainsi, $h$ est compatible avec la structure de $A$-module et la graduation
mais pas avec la différentielle. Si $f = g$ et si $h$ est une homotopie
au sens de la définition, alors $h$ définit un morphisme $h : M \to N[-1]$
dans $\text{Mod}_{(A, \text{d})}$.

\begin{lemma}
\label{lemma-compose-homotopy}
Soit $(A, \text{d})$ une algèbre différentielle graduée.
Soient $f, g : L \to M$ des homomorphismes de $A$-modules différentiels gradués.
Supposons donnés en outre des homomorphismes $a : K \to L$ et $c : M \to N$.
Si $h : L \to M$ est une application de $A$-modules qui définit une homotopie
entre $f$ et $g$, alors $c \circ h \circ a$ définit une homotopie
entre $c \circ f \circ a$ et $c \circ g \circ a$.
\end{lemma}

\begin{proof}
Résulte immédiatement d'Homologie, lemme \ref{homology-lemma-compose-homotopy-cochain}.
\end{proof}

\noindent
Ce lemme nous permet de définir la catégorie homotopique comme suit.

\begin{definition}
\label{definition-complexes-notation}
Soit $(A, \text{d})$ une algèbre différentielle graduée.
La {\it catégorie homotopique}, notée $K(\text{Mod}_{(A, \text{d})})$, est la
catégorie dont les objets sont les objets de
$\text{Mod}_{(A, \text{d})}$ et dont les morphismes sont les classes d'homotopie
des homomorphismes de $A$-modules différentiels gradués.
\end{definition}

\noindent
La notation $K(\text{Mod}_{(A, \text{d})})$ n'est pas standard mais elle est au
moins cohérente avec l'utilisation de $K(-)$ dans d'autres parties du projet Champs.

\begin{lemma}
\label{lemma-homotopy-direct-sums}
Soit $(A, \text{d})$ une algèbre différentielle graduée. La
catégorie homotopique $K(\text{Mod}_{(A, \text{d})})$ possède des
sommes directes et des produits.
\end{lemma}

\begin{proof}
Démonstration omise. Indication : utiliser les sommes directes et les
produits comme dans le lemme \ref{lemma-dgm-abelian}. Cela fonctionne
parce que nous avons vu que ces foncteurs commutent avec le foncteur d'oubli vers la
catégorie des $A$-modules gradués et parce que $\prod$ est un foncteur
exact sur la catégorie des familles de groupes abéliens.
\end{proof}







\section{Cônes}
\label{section-cones}

\noindent
Nous introduisons les cônes pour la catégorie des modules différentiels gradués.

\begin{definition}
\label{definition-cone}
Soit $(A, \text{d})$ une algèbre différentielle graduée.
Soit $f : K \to L$ un homomorphisme de $A$-modules différentiels gradués. Le
{\it cône} de $f$ est le $A$-module différentiel gradué $C(f)$
donné par $C(f) = L \oplus K$ avec la graduation
$C(f)^n = L^n \oplus K^{n + 1}$ et la
différentielle
$$
d_{C(f)} =
\left(
\begin{matrix}
\text{d}_L & f \\
0 & -\text{d}_K
\end{matrix}
\right)
$$
Il est muni de morphismes canoniques de complexes $i : L \to C(f)$ et $p :
C(f) \to K[1]$ induits par les applications évidentes $L \to C(f)$
et $C(f) \to K$.
\end{definition}

\noindent
La formation du triangle du cône est fonctorielle dans le sens suivant.

\begin{lemma}
\label{lemma-functorial-cone}
Soit $(A, \text{d})$ une algèbre différentielle graduée.
Supposons que
$$
\xymatrix{
K_1 \ar[r]_{f_1} \ar[d]_a & L_1 \ar[d]^b \\
K_2 \ar[r]^{f_2} & L_2
}
$$
soit un diagramme d'homomorphismes de $A$-modules différentiels gradués
qui est commutatif à homotopie
près. Alors il existe un morphisme $c : C(f_1) \to C(f_2)$ qui induit un
morphisme de triangles
$$
(a, b, c) : (K_1, L_1, C(f_1), f_1, i_1, p_1) \to
(K_1, L_1, C(f_1), f_2, i_2, p_2)
$$
dans $K(\text{Mod}_{(A, \text{d})})$.
\end{lemma}

\begin{proof}
Soit $h : K_1 \to L_2$ une homotopie entre $f_2 \circ a$ et $b \circ f_1$.
Définissons $c$ par la matrice
$$
c =
\left(
\begin{matrix}
b & h \\
0 & a
\end{matrix}
\right) :
L_1 \oplus K_1 \to L_2 \oplus K_2
$$
Un calcul matriciel montre que $c$ est un morphisme de modules
différentiels gradués. Il est trivial que $c \circ i_1 = i_2 \circ b$, et il est
également trivial de vérifier que $p_2 \circ c = a \circ p_1$.
\end{proof}











\section{Suites exactes courtes admissibles}
\label{section-admissible}

\noindent
Une suite exacte courte admissible est l'analogue des suites exactes scindées
terme à terme dans le cadre des modules différentiels gradués.

\begin{definition}
\label{definition-admissible-ses}
Soit $(A, \text{d})$ une algèbre différentielle graduée.
\begin{enumerate}
\item Un homomorphisme $K \to L$ de modules différentiels gradués sur $A$
est un {\it monomorphisme admissible} s'il existe un morphisme de $A$-modules
gradués $L \to K$ qui soit inverse à gauche de $K \to L$.
\item Un homomorphisme $L \to M$ de $A$-modules différentiels gradués
est un {\it épimorphisme admissible} s'il existe un morphisme de $A$-modules
gradués $M \to L$ qui soit un inverse à droite de $L \to M$.
\item Une suite exacte courte $0 \to K \to L \to M \to 0$ de $A$-modules
différentiels gradués est une {\it suite exacte courte admissible} si
elle se scinde en tant que suite de $A$-modules gradués.
\end{enumerate}
\end{definition}

\noindent
Ainsi, les scindages sont compatibles avec toutes les données sauf les
différentielles. Étant donnée une suite exacte courte admissible, nous
obtenons un triangle ; c'est la raison pour laquelle nous exigeons que nos
scindages soient compatibles avec la structure de $A$-module.

\begin{lemma}
\label{lemma-admissible-ses}
Soit $(A, \text{d})$ une algèbre différentielle graduée.
Soit $0 \to K \to L \to M \to 0$ une suite exacte courte admissible de
$A$-modules différentiels gradués. Soient $s : M \to L$ et $\pi : L \to K$ des
scindages tels que $\Ker(\pi) = \Im(s)$. Alors
nous obtenons un morphisme
$$
\delta = \pi \circ \text{d}_L \circ s : M \to K[1]
$$
de $\text{Mod}_{(A, \text{d})}$ qui induit les homomorphismes de
liaison dans la suite exacte longue de cohomologie (\ref{equation-les}).
\end{lemma}

\begin{proof}
L'application $\pi \circ \text{d}_L \circ s$ est compatible par construction
avec la structure de $A$-module et les graduations. Elle est compatible
avec les différentielles d'après
Homologie, lemme
\ref{homology-lemma-ses-termwise-split-cochain}. Soit $R$ l'anneau sur lequel $A$ est une algèbre différentielle
graduée. L'égalité des applications est une assertion sur les $R$-modules.
Elle découle donc d'Homologie,
lemmes \ref{homology-lemma-ses-termwise-split-cochain} et
\ref{homology-lemma-ses-termwise-split-long-cochain}.
\end{proof}

\begin{lemma}
\label{lemma-make-commute-map}
Soit $(A, \text{d})$ une algèbre différentielle graduée. Soit
$$
\xymatrix{
K \ar[r]_f \ar[d]_a & L \ar[d]^b \\ M
\ar[r]^g & N
}
$$
un diagramme d'homomorphismes de $A$-modules différentiels gradués
qui commute à homotopie près.
\begin{enumerate}
\item Si $f$ est un monomorphisme admissible, alors $b$ est homotope à un
homomorphisme qui rend le diagramme commutatif.
\item Si $g$ est un épimorphisme admissible, alors $a$ est homotope à
un morphisme qui fait commuter le diagramme.
\end{enumerate}
\end{lemma}

\begin{proof}
Soit $h : K \to N$ une homotopie entre $bf$ et $ga$, c'est-à-dire $bf
- ga = \text{d}h + h\text{d}$. Supposons que $\pi : L \to K$ soit une
application de $A$-modules gradués, inverse à gauche de $f$.
Prenons $b' = b - \text{d}h\pi - h\pi \text{d}$.
Supposons que $s : N \to M$ soit une application de $A$-modules gradués, inverse à
droite de $g$. Prenons $a' = a + \text{d}sh +
sh\text{d}$. Calculs omis.
\end{proof}

\begin{lemma}
\label{lemma-make-injective}
Soit $(A, \text{d})$ une algèbre différentielle
graduée. Soit $\alpha : K \to L$ un homomorphisme de $A$-modules
différentiels gradués. Il existe une factorisation
$$
\xymatrix{
K \ar[r]^{\tilde \alpha} \ar@/_1pc/[rr]_\alpha &
\tilde L \ar[r]^\pi & L
}
$$
dans $\text{Mod}_{(A, \text{d})}$ telle que
\begin{enumerate}
\item $\tilde \alpha$ soit un monomorphisme admissible (voir
la définition \ref{definition-admissible-ses}),
\item il existe un morphisme $s : L \to \tilde
L$ tel que $\pi \circ s = \text{id}_L$ et tel que
$s \circ \pi$ soit homotope à $\text{id}_{\tilde L}$.
\end{enumerate}
\end{lemma}

\begin{proof}
La démonstration est identique à celle de
Catégories dérivées, lemme \ref{derived-lemma-make-injective}. En effet,
nous posons $\tilde L = L \oplus C(1_K)$ et nous utilisons des propriétés
élémentaires de la construction du cône.
\end{proof}

\begin{lemma}
\label{lemma-sequence-maps-split}
Soit $(A, \text{d})$ une algèbre différentielle
graduée. Soit $L_1 \to L_2 \to \ldots \to
L_n$ une suite d'homomorphismes composables de
$A$-modules différentiels gradués.
Il existe un diagramme commutatif
$$
\xymatrix{
L_1 \ar[r] &
L_2 \ar[r] &
\ldots \ar[r] &
L_n \\ M_1
\ar[r] \ar[u] & M_2
\ar[r] \ar[u] & \ldots
\ar[r] & M_n
\ar[u]
}
$$
dans $\text{Mod}_{(A, \text{d})}$ tel que chaque $M_i \to M_{i + 1}$ soit
un monomorphisme admissible et que chaque $M_i \to L_i$ soit
une équivalence d'homotopie.
\end{lemma}

\begin{proof}
Le cas $n = 1$ est immédiat. Le lemme
\ref{lemma-make-injective} traite le cas $n = 2$.
Supposons le diagramme construit à l'exception
de $M_n$. Appliquons le lemme \ref{lemma-make-injective} à la
composition $M_{n - 1} \to L_{n - 1} \to L_n$. Le
résultat est une factorisation $M_{n - 1} \to M_n \to L_n$ comme
souhaité.
\end{proof}



\begin{lemma}
\label{lemma-nilpotent}
Soit $(A, \text{d})$ une algèbre différentielle graduée.
Soient $0 \to K_i \to L_i \to M_i \to 0$, $i = 1, 2, 3$
des suites exactes courtes admissibles de $A$-modules différentiels gradués.
Soient $b : L_1 \to L_2$ et $b' : L_2 \to L_3$
des homomorphismes de modules différentiels gradués tels que
$$
\vcenter{
\xymatrix{
K_1 \ar[d]_0 \ar[r] &
L_1 \ar[r] \ar[d]_b &
M_1 \ar[d]_0 \\ K_2
\ar[r] & L_2 \ar[r] & M_2 } }
\quad\text{et}\quad
\vcenter{
\xymatrix{
K_2
\ar[d]^0
\ar[r] & L_2 \ar[r]
\ar[d]^{b'} & M_2 \ar[d]^0
\\ K_3 \ar[r] &
L_3 \ar[r] & M_3
}
}
$$
commutent à homotopie près. Alors $b' \circ b$ est homotope à $0$.
\end{lemma}

\begin{proof}
Par le lemme \ref{lemma-make-commute-map}, nous pouvons remplacer $b$ et $b'$
par des applications homotopes telles que le carré droit du diagramme de gauche
commute et le carré gauche du diagramme de droite commute. En d'autres termes, nous avons
$\Im(b) \subset \Im(K_2 \to L_2)$ et
$\Ker(b') \supset \Im(K_2 \to L_2)$. Ensuite, $b'
\circ b = 0$ en tant qu'homomorphisme de modules.
\end{proof}
















\section{Triangles distingués}
\label{section-distinguished}

\noindent
Le lemme suivant produit nos triangles distingués.

\begin{lemma}
\label{lemma-triangle-independent-splittings}
Soit $(A, \text{d})$ une algèbre différentielle graduée.
Soit $0 \to K \to L \to M \to 0$ une suite exacte courte admissible
de $A$-modules différentiels gradués. Le triangle
\begin{equation}
\label{equation-triangle-associated-to-admissible-ses}
K \to L \to M \xrightarrow{\delta} K[1]
\end{equation}
avec $\delta$ comme dans le lemme \ref{lemma-admissible-ses} est, à
isomorphisme canonique près dans $K(\text{Mod}_{(A, \text{d})})$, indépendant des choix
effectués dans le lemme \ref{lemma-admissible-ses}.
\end{lemma}

\begin{proof}
En effet, soit $(s', \pi')$ un second choix de scindages comme
dans le lemme \ref{lemma-admissible-ses}. Alors nous affirmons que $\delta$ et
$\delta'$ sont homotopes. Écrivons $s' = s + \alpha \circ h$ et
$\pi' = \pi + g \circ \beta$ pour certains homomorphismes uniques
de $A$-modules $h : M \to K$ et $g : M \to K$ de degré $-1$. Alors
$g = -h$ et $g$ est une homotopie entre $\delta$ et $\delta'$. Les
calculs sont effectués dans la démonstration
d'Homologie, lemme \ref{homology-lemma-ses-termwise-split-homotopy-cochain}.
\end{proof}

\begin{definition}
\label{definition-distinguished-triangle}
Soit $(A, \text{d})$ une algèbre différentielle graduée.
\begin{enumerate}
\item Si $0 \to K \to L \to M \to 0$ est une suite exacte courte
admissible de $A$-modules différentiels gradués, alors le {\it triangle associé
à $0 \to K \to L \to M \to 0$} est le triangle
(\ref{equation-triangle-associated-to-admissible-ses}) de
$K(\text{Mod}_{(A, \text{d})})$.
\item Un triangle de $K(\text{Mod}_{(A, \text{d})})$ est
appelé {\it triangle distingué} s'il est isomorphe à un triangle
associé à une suite exacte courte admissible de
$A$-modules différentiels gradués.
\end{enumerate}
\end{definition}









\section{Cônes et triangles distingués}
\label{section-cones-and-triangles}

\noindent
Soit $(A, \text{d})$ une algèbre différentielle graduée.
Soit $f : K \to L$ un homomorphisme de $A$-modules différentiels gradués.
Alors $(K, L, C(f), f, i, p)$ forme un triangle :
$$
K \to L \to C(f) \to K[1]
$$
dans $\text{Mod}_{(A, \text{d})}$ et donc dans $K(\text{Mod}_{(A, \text{d})})$.
Les cônes ne sont {\bf pas} des triangles distingués en général, mais la
différence tient à un signe ou à une rotation, au choix. Voici deux énoncés précis.

\begin{lemma}
\label{lemma-rotate-cone}
Soit $(A, \text{d})$ une algèbre différentielle
graduée. Soit $f : K \to L$ un homomorphisme de modules différentiels
gradués. Le triangle $(L, C(f), K[1], i, p,
f[1])$ est le triangle associé à la suite exacte courte admissible
$$
0 \to L \to C(f) \to K[1] \to 0
$$
provenant de la définition du cône de $f$.
\end{lemma}

\begin{proof}
Immédiat d'après les définitions.
\end{proof}

\begin{lemma}
\label{lemma-rotate-triangle}
Soit $(A, \text{d})$ une algèbre différentielle graduée.
Soient $\alpha : K \to L$ et $\beta : L \to M$
définissant une suite exacte courte admissible
$$
0 \to K \to L \to M \to 0
$$
de modules différentiels gradués
sur $A$. Soit $(K, L, M, \alpha, \beta,
\delta)$ le triangle associé. Alors les triangles
$$
(M[-1], K, L, \delta[-1], \alpha, \beta)
\quad\text{et}\quad
(M[-1], K, C(\delta[-1]), \delta[-1], i, p)
$$
sont isomorphes.
\end{lemma}

\begin{proof}
En utilisant un choix de scindages, nous écrivons $L = K \oplus M$ et nous
identifions $\alpha$ et $\beta$ avec les applications d'inclusion et de projection
naturelles. Par construction de $\delta$, nous avons
$$
d_B =
\left(
\begin{matrix}
d_K & \delta \\
0 & d_M
\end{matrix}
\right)
$$
D'autre part, le cône de $\delta[-1] : M[-1] \to K$ est
donné par $C(\delta[-1]) = K \oplus M$ avec une différentielle identique
à la matrice ci-dessus ! D'où le lemme.
\end{proof}

\begin{lemma}
\label{lemma-third-isomorphism}
Soit $(A, \text{d})$ une algèbre différentielle graduée.
Soient $f_1 : K_1 \to L_1$ et $f_2 : K_2 \to L_2$ des homomorphismes de
modules différentiels gradués sur $A$. Soit
$$
(a, b, c) :
(K_1, L_1, C(f_1), f_1, i_1, p_1)
\longrightarrow
(K_1, L_1, C(f_1), f_2, i_2, p_2)
$$
un morphisme quelconque de triangles de $K(\text{Mod}_{(A, \text{d})})$. Si $a$ et
$b$ sont des équivalences par homotopie, alors il en est de même pour $c$.
\end{lemma}

\begin{proof}
Soit $a^{-1} : K_2 \to K_1$ un homomorphisme de $A$-modules différentiels gradués qui
est l'inverse de $a$ dans $K(\text{Mod}_{(A, \text{d})})$. Soit
$b^{-1} : L_2 \to L_1$ un homomorphisme de $A$-modules différentiels gradués qui est
l'inverse de $b$ dans $K(\text{Mod}_{(A, \text{d})})$. Soit $c' :
C(f_2) \to C(f_1)$ le morphisme issu du lemme
\ref{lemma-functorial-cone} appliqué à $f_1 \circ a^{-1} =
b^{-1} \circ f_2$. Si nous pouvons montrer
que $c \circ c'$ et $c' \circ c$ sont des isomorphismes dans
$K(\text{Mod}_{(A, \text{d})})$, la démonstration
sera achevée. Il suffit donc de prouver ce qui suit : étant donné un
morphisme de triangles
$(1, 1, c) : (K, L, C(f), f, i, p)$ dans
$K(\text{Mod}_{(A, \text{d})})$, le morphisme $c$ est un isomorphisme dans
$K(\text{Mod}_{(A, \text{d})})$. Par
hypothèse, les deux carrés du diagramme
$$
\xymatrix{
L \ar[r] \ar[d]_1 &
C(f) \ar[r] \ar[d]_c &
K[1] \ar[d]_1 \\ L
\ar[r] & C(f)
\ar[r] &
K[1]
}
$$
commutent à homotopie près. Par construction de $C(f)$, les lignes
forment des suites exactes courtes admissibles. Ainsi, nous voyons
que $(c - 1)^2 = 0$ dans $K(\text{Mod}_{(A, \text{d})})$ par
le lemme \ref{lemma-nilpotent}. Par
conséquent, $c$ est un isomorphisme dans $K(\text{Mod}_{(A, \text{d})})$
avec inverse $2 - c$.
\end{proof}

\noindent
Le lemme suivant montre que les triangles de la catégorie homotopique
donnés par les cônes et les triangles distingués coïncident à
isomorphisme près, du moins à un signe près !

\begin{lemma}
\label{lemma-the-same-up-to-isomorphisms}
Soit $(A, \text{d})$ une algèbre différentielle graduée.
\begin{enumerate}
\item Étant donnée une suite exacte courte
admissible $0 \to K \xrightarrow{\alpha} L \to M \to 0$
de $A$-modules différentiels gradués, il existe une équivalence d'homotopie
$C(\alpha) \to M$ telle que le diagramme
$$
\xymatrix{
K \ar[r] \ar[d] & L \ar[d] \ar[r] &
C(\alpha) \ar[r]_{-p} \ar[d] & K[1] \ar[d] \\ K
\ar[r]^\alpha & L \ar[r]^\beta & M
\ar[r]^\delta & K[1]
}
$$
définisse un isomorphisme de triangles dans $K(\text{Mod}_{(A, \text{d})})$.
\item Étant donné un morphisme de complexes $f : K
\to L$, il existe un isomorphisme de triangles
$$
\xymatrix{
K \ar[r] \ar[d] & \tilde L \ar[d] \ar[r] &
M \ar[r]_{\delta} \ar[d] & K[1] \ar[d] \\ K
\ar[r] & L \ar[r] &
C(f) \ar[r]^{-p} & K[1]
}
$$
où le triangle supérieur est le triangle associé à une
suite exacte courte admissible $K \to \tilde L \to M$.
\end{enumerate}
\end{lemma}

\begin{proof}
Démonstration de (1). Nous avons $C(\alpha) = L \oplus K$ et nous
définissons simplement $C(\alpha) \to M$ via la projection sur $L$ suivie de
$\beta$. Cela définit un morphisme de modules différentiels gradués car les
compositions $K^{n + 1} \to L^{n + 1} \to M^{n + 1}$ sont nulles.
Choisissons des scindages $s : M \to L$ et $\pi : L \to K$ avec
$\Ker(\pi) = \Im(s)$ et posons $\delta
= \pi \circ \text{d}_L \circ s$ comme d'habitude. Pour
obtenir un inverse par homotopie,
nous prenons $M \to C(\alpha)$ donné par $(s , -\delta)$. Cela est compatible
avec les différentielles car $\delta^n$ peut être caractérisé comme
l'application unique $M^n \to K^{n + 1}$
telle que $\text{d} \circ s^n - s^{n + 1} \circ \text{d} = \alpha \circ \delta^n$,
voir la
démonstration de Homologie, lemme
\ref{homology-lemma-ses-termwise-split-cochain}. La composition $M \to C(f) \to M$ est l'identité. La
composition $C(f) \to M \to C(f)$ est égale au morphisme
$$
\left(
\begin{matrix}
s \circ \beta & 0 \\
-\delta \circ \beta & 0
\end{matrix}
\right)
$$
Pour voir que cela est homotope à l'application
identité, utilisons l'homotopie $h : C(\alpha) \to C(\alpha)$
donnée par la matrice
$$
\left(
\begin{matrix}
0 & 0 \\
\pi & 0
\end{matrix}
\right) :
C(\alpha) = L \oplus K
\to
L \oplus K = C(\alpha)
$$
Il est trivial de vérifier que
$$
\left(
\begin{matrix}
1 & 0 \\
0 & 1
\end{matrix}
\right)
-
\left(
\begin{matrix}
s \\
-\delta
\end{matrix}
\right)
\left(
\begin{matrix}
\beta & 0
\end{matrix}
\right)
=
\left(
\begin{matrix}
\text{d} & \alpha \\
0 & -\text{d}
\end{matrix}
\right)
\left(
\begin{matrix}
0 & 0 \\
\pi & 0
\end{matrix}
\right)
+
\left(
\begin{matrix}
0 & 0 \\
\pi & 0
\end{matrix}
\right)
\left(
\begin{matrix}
\text{d} & \alpha \\
0 & -\text{d}
\end{matrix}
\right)
$$
Pour terminer la démonstration de (1), nous devons montrer que les
morphismes $-p : C(\alpha) \to K[1]$ (voir
définition \ref{definition-cone}) et
$C(\alpha) \to M \to K[1]$ concordent à homotopie
près. Cela est clair d'après ce qui précède. En effet, nous pouvons utiliser l'inverse
d'homotopie $(s, -\delta) : M \to C(\alpha)$ et
vérifier à la place que les deux applications
$M \to K[1]$ concordent. Notons en outre que $p
\circ (s, -\delta) = -\delta$ comme souhaité.

\medskip\noindent
Démonstration de (2). Nous notons $\tilde f : K \to
\tilde L$, $s : L \to \tilde
L$ et $\pi : L \to L$ comme dans
le lemme \ref{lemma-make-injective}.
D'après les lemmes \ref{lemma-functorial-cone} et \ref{lemma-third-isomorphism},
les triangles $(K, L, C(f), i, p)$ et $(K,
\tilde L, C(\tilde f), \tilde i, \tilde p)$ sont isomorphes.
Notons que nous pouvons composer les isomorphismes de triangles.
Ainsi, nous pouvons remplacer $L$ par
$\tilde L$ et $f$ par $\tilde f$. En d'autres termes, nous
pouvons supposer que $f$ est un monomorphisme admissible. Dans ce
cas, le résultat découle de la partie (1).
\end{proof}







\section{La catégorie homotopique est triangulée}
\label{section-homotopy-triangulated}

\noindent
Montrons d'abord qu'elle est prétriangulée.

\begin{lemma}
\label{lemma-homotopy-category-pre-triangulated}
Soit $(A, \text{d})$ une algèbre différentielle graduée.
La catégorie homotopique $K(\text{Mod}_{(A, \text{d})})$
avec ses foncteurs de translation naturels et ses triangles distingués
est une catégorie pré-triangulée.
\end{lemma}

\begin{proof}
Démonstration de TR1. Par définition, tout triangle isomorphe à un triangle
distingué est distingué. De plus, tout triangle $(K, K, 0, 1, 0,
0)$ est distingué puisque $0 \to K \to K \to 0 \to 0$ est
une suite exacte courte admissible. Enfin, étant donné tout homomorphisme
$f : K \to L$ de $A$-modules différentiels gradués, le triangle
$(K, L, C(f), f, i, -p)$ est distingué d'après le
lemme \ref{lemma-the-same-up-to-isomorphisms}.

\medskip\noindent
Démonstration de TR2. Soit $(X, Y, Z, f, g, h)$ un
triangle. Supposons que $(Y, Z, X[1], g, h, -f[1])$ soit
distingué. Alors il existe une suite exacte courte admissible
$0 \to K \to L \to M \to 0$ telle que le triangle
associé $(K, L, M, \alpha, \beta, \delta)$ soit
isomorphe à $(Y, Z, X[1], g, h, -f[1])$. En effectuant une rotation inverse,
nous voyons que $(X, Y, Z, f, g, h)$ est
isomorphe à $(M[-1], K, L, -\delta[-1], \alpha,
\beta)$. Il découle du lemme \ref{lemma-rotate-triangle} que le triangle
$(M[-1], K, L, \delta[-1], \alpha, \beta)$ est
isomorphe à $(M[-1],
K, C(\delta[-1]), \delta[-1], i, p)$. En
précomposant l'isomorphisme précédent des triangles avec $-1$ sur $Y$, il s'ensuit
que $(X, Y, Z, f, g, h)$ est isomorphe à $(M[-1], K,
C(\delta[-1]), \delta[-1], i, -p)$. Par conséquent, il
est distingué par le lemme
\ref{lemma-the-same-up-to-isomorphisms}. D'autre part,
supposons que $(X, Y, Z, f, g, h)$ soit distingué. Par le lemme
\ref{lemma-the-same-up-to-isomorphisms}, cela signifie qu'il est isomorphe à un
triangle de la forme $(K, L, C(f), f, i, -p)$
pour un certain morphisme $f$ de $\text{Mod}_{(A,
\text{d})}$. Alors, le triangle obtenu par rotation $(Y, Z,
X[1], g, h, -f[1])$ est isomorphe à
$(L, C(f), K[1], i, -p, -f[1])$, qui est isomorphe au
triangle $(L, C(f), K[1], i, p,
f[1])$. Par le lemme
\ref{lemma-rotate-cone}, ce triangle est distingué. Par conséquent, $(Y, Z, X[1], g,
h, -f[1])$ est distingué comme souhaité.

\medskip\noindent
Démonstration de TR3. Soient $(X, Y, Z, f, g, h)$ et $(X', Y', Z', f', g',
h')$ des triangles distingués de $K(\mathcal{A})$ et soient $a : X \to X'$
et $b : Y \to Y'$ des morphismes tels que $f' \circ a = b \circ f$. D'après
le lemme \ref{lemma-the-same-up-to-isomorphisms}, nous pouvons supposer
que $(X, Y, Z, f, g, h) = (X, Y, C(f), f, i, -p)$ et
$(X', Y', Z', f', g', h') = (X', Y', C(f'), f', i', -p')$. À ce
stade, nous appliquons simplement le lemme \ref{lemma-functorial-cone}
au diagramme commutatif donné par $f, f', a, b$.
\end{proof}

\noindent
Avant de prouver TR4 en général, nous le prouvons dans un cas particulier.

\begin{lemma}
\label{lemma-two-split-injections}
Soit $(A, \text{d})$ une algèbre différentielle graduée. Supposons que
$\alpha : K \to L$ et $\beta : L \to M$ soient des monomorphismes admissibles
de modules différentiels gradués sur $A$. Alors il existe des triangles distingués
$(K, L, Q_1, \alpha, p_1, d_1)$, $(K, M, Q_2, \beta \circ \alpha, p_2, d_2)$ et $(L,
M, Q_3, \beta, p_3, d_3)$ pour lesquels TR4 est vérifié.
\end{lemma}

\begin{proof}
Soient $\pi_1 : L \to K$ et $\pi_3 : M \to L$ des homomorphismes de
$A$-modules gradués qui sont des inverses à gauche de $\alpha$ et $\beta$.
Alors $K \to M$ est également un monomorphisme admissible avec un
inverse à gauche $\pi_2 = \pi_1 \circ \pi_3$. Notons $Q_1$, $Q_2$ et
$Q_3$ pour les conoyaux de $K \to L$, $K \to M$ et $L \to M$. Alors nous
obtenons les identifications (en tant que $A$-modules gradués)
$Q_1 = \Ker(\pi_1)$, $Q_3 = \Ker(\pi_3)$ et $Q_2 =
\Ker(\pi_2)$. Alors $L = K \oplus Q_1$ et $M = L \oplus
Q_3$ en tant que $A$-modules gradués. Cela implique $M = K
\oplus Q_1 \oplus Q_3$. Notons que $\pi_2 = \pi_1 \circ \pi_3$ est nul à la
fois sur $Q_1$ et $Q_3$. Par conséquent $Q_2 = Q_1 \oplus Q_3$.
Considérons le diagramme commutatif
$$
\begin{matrix}
0 & \to & K & \to & L & \to & Q_1 & \to &
0 \\ & & \downarrow & & \downarrow & & \downarrow & \\ 0 &
\to & K & \to & M & \to & Q_2 & \to & 0
\\ & & \downarrow & & \downarrow & & \downarrow & \\ 0 & \to
& L & \to & M & \to & Q_3 & \to & 0
\end{matrix}
$$
Les lignes de ce diagramme sont des suites exactes courtes admissibles, et
déterminent donc par définition des triangles distingués. De plus, les
flèches descendantes dans le diagramme ci-dessus sont compatibles avec les
scindages choisis et définissent donc des morphismes de triangles
$$
(K \to L \to Q_1 \to K[1])
\longrightarrow
(K \to M \to Q_2 \to K[1])
$$
et
$$
(K \to M \to Q_2 \to K[1])
\longrightarrow
(L \to M \to Q_3 \to L[1]).
$$
Notons que les scindages $Q_3 \to M$ de la suite inférieure dans le
diagramme fournissent un scindage pour la suite
scindée $0 \to Q_1 \to Q_2 \to Q_3 \to 0$ lorsqu'on les compose avec $M
\to Q_2$. Il en découle facilement que le morphisme $\delta : Q_3 \to
Q_1[1]$ dans le triangle distingué correspondant
$$
(Q_1 \to Q_2 \to Q_3 \to Q_1[1])
$$
est égal à la composition $Q_3 \to L[1] \to Q_1[1]$. Ainsi,
nous obtenons une structure comme dans la conclusion de l'axiome TR4.
\end{proof}

\noindent
Voici le résultat final.

\begin{proposition}
\label{proposition-homotopy-category-triangulated}
Soit $(A, \text{d})$ une algèbre différentielle graduée. La catégorie homotopique
$K(\text{Mod}_{(A, \text{d})})$ des $A$-modules différentiels gradués avec ses
foncteurs de translation naturels et ses triangles distingués est une catégorie
triangulée.
\end{proposition}

\begin{proof}
Nous savons que $K(\text{Mod}_{(A, \text{d})})$ est une catégorie pré-triangulée. Par
conséquent, il suffit de prouver TR4 et pour le prouver, nous
pouvons utiliser Catégories dérivées, lemme
\ref{derived-lemma-easier-axiom-four}. Soient $K \to L$ et $L \to M$ des morphismes composables
de $K(\text{Mod}_{(A, \text{d})})$.
D'après le lemme \ref{lemma-sequence-maps-split}, nous pouvons
supposer que $K \to L$ et $L \to M$ sont des monomorphismes
admissibles. Dans ce cas, le résultat découle du
lemme \ref{lemma-two-split-injections}.
\end{proof}











\section{Modules à gauche}
\label{section-left-modules}

\noindent
Tout ce que nous avons dit jusqu'à présent a un analogue dans le cadre
des modules différentiels gradués à gauche, sauf qu'il faut faire
attention à certaines règles de signe.

\medskip\noindent
Soit $(A, \text{d})$ une $R$-algèbre différentielle graduée.
Par analogie avec les modules à droite, nous
définissons un {\it module différentiel gradué à gauche sur $A$} $M$
comme un $A$-module à gauche $M$ muni d'une graduation $M =
\bigoplus M^n$ et d'une différentielle $\text{d}$ telles que $A^n M^m
\subset M^{n + m}$, que $\text{d}(M^n) \subset M^{n + 1}$, et que
$$
\text{d}(am) = \text{d}(a) m + (-1)^{\deg(a)}a \text{d}(m)
$$
pour les éléments homogènes $a \in A$ et $m \in M$. Comme auparavant, cette
règle de Leibniz signifie exactement que la multiplication définit un
morphisme de complexes
$$
\text{Tot}(A^\bullet \otimes_R M^\bullet) \to M^\bullet
$$
Ici, $A^\bullet$ et $M^\bullet$ désignent les complexes de $R$-modules
sous-jacents à $A$ et $M$.

\begin{definition}
\label{definition-opposite-dga}
Soit $R$ un anneau. Soit $(A, \text{d})$ une algèbre différentielle graduée
sur $R$. L'{\it algèbre différentielle graduée opposée} est l'algèbre
différentielle graduée $(A^{opp}, \text{d})$ sur $R$ où $A^{opp} = A$ en tant
que $R$-module gradué, $\text{d} = \text{d}$, et la multiplication est
donnée par
$$
a \cdot_{opp} b = (-1)^{\deg(a)\deg(b)} b a
$$
pour les éléments homogènes $a, b \in A$.
\end{definition}

\noindent
Cela a du sens parce que
\begin{align*}
\text{d}(a \cdot_{opp} b)
& =
(-1)^{\deg(a)\deg(b)} \text{d}(b a) \\ &
=
(-1)^{\deg(a)\deg(b)} \text{d}(b) a +
(-1)^{\deg(a)\deg(b) + \deg(b)}b\text{d}(a) \\ & =
(-1)^{\deg(a)}a
\cdot_{opp} \text{d}(b) + \text{d}(a) \cdot_{opp} b
\end{align*}
comme souhaité. En termes de complexes sous-jacents de $R$-modules,
cela signifie que le diagramme
$$
\xymatrix{
\text{Tot}(A^\bullet \otimes_R A^\bullet)
\ar[rrr]_-{\text{multiplication de }A^{opp}}
\ar[d]_{\text{contrainte de commutativité}} & & &
A^\bullet \ar[d]^{\text{id}} \\
\text{Tot}(A^\bullet \otimes_R A^\bullet)
\ar[rrr]^-{\text{multiplication de }A} & & &
A^\bullet
}
$$
commute. Ici, la contrainte de commutativité sur la catégorie monoïdale
symétrique des complexes de $R$-modules est donnée dans
Compléments d'algèbre, section \ref{more-algebra-section-sign-rules}.

\medskip\noindent
Soit $(A, \text{d})$ une algèbre différentielle graduée sur $R$.
Soit $M$ un $A$-module différentiel gradué à gauche. Nous noterons
$M^{opp}$ le module $M$ considéré comme un $A^{opp}$-module à droite avec
multiplication $\cdot_{opp}$ définie par la règle
$$
m \cdot_{opp} a = (-1)^{\deg(a)\deg(m)} a m
$$
pour $a$ et $m$ homogènes. Cela est compatible avec les différentielles
car nous aurions pu utiliser le diagramme
$$
\xymatrix{
\text{Tot}(M^\bullet \otimes_R A^\bullet)
\ar[rrr]_-{\text{multiplication sur }M^{opp}}
\ar[d]_{\text{contrainte de commutativité}} & & &
M^\bullet \ar[d]^{\text{id}} \\
\text{Tot}(A^\bullet \otimes_R M^\bullet)
\ar[rrr]^-{\text{multiplication sur }M} & & &
M^\bullet
}
$$
pour définir la multiplication $\cdot_{opp}$ sur $M^{opp}$.
Pour voir qu'il s'agit d'une multiplication associative, nous calculons
pour les éléments homogènes $a, b \in A$ et $m \in M$ que
\begin{align*}
m \cdot_{opp} (a \cdot_{opp} b)
& =
(-1)^{\deg(a)\deg(b)} m \cdot_{opp} (ba) \\ &
=
(-1)^{\deg(a)\deg(b) + \deg(ab)\deg(m)} bam \\ & =
(-1)^{\deg(a)\deg(b)
+ \deg(ab)\deg(m) + \deg(b)\deg(am)} (am)
\cdot_{opp} b \\ & =
(-1)^{\deg(a)\deg(b)
+ \deg(ab)\deg(m) + \deg(b)\deg(am) + \deg(a)\deg(m)} (m
\cdot_{opp} a) \cdot_{opp} b \\ & = (m
\cdot_{opp}
a) \cdot_{opp} b
\end{align*}
Bien sûr, nous aurions pu le montrer en utilisant la compatibilité entre les
contraintes d'associativité et de commutativité dans la catégorie monoïdale
symétrique des complexes de $R$-modules également.

\begin{lemma}
\label{lemma-left-right}
Soit $(A, \text{d})$ une $R$-algèbre différentielle graduée.
Le foncteur $M \mapsto M^{opp}$ de la catégorie des
$A$-modules différentiels gradués à gauche vers la catégorie des
$A^{opp}$-modules différentiels gradués à droite est une équivalence.
\end{lemma}

\begin{proof}
Démonstration omise.
\end{proof}

\noindent
Passons maintenant aux décalages. Soit $(A, \text{d})$ une algèbre différentielle
graduée. Soit $M$ un $A$-module différentiel gradué à gauche dont le complexe
sous-jacent de $R$-modules est noté $M^\bullet$.
Pour tout $k \in \mathbf{Z}$ nous définissons le {\it module décalé de $k$}
$M[k]$ comme suit
\begin{enumerate}
\item le complexe de $R$-modules sous-jacent à $M[k]$ est $M^\bullet[k]$
\item la structure de $A$-module est donnée par la multiplication
$$
A^n \times (M[k])^m \longrightarrow (M[k])^{n + m}
$$
est égale à $(-1)^{nk}$ multiplié par la multiplication
donnée $A^n \times M^{m + k} \to M^{n + m + k}$.
\end{enumerate}
Vérifions qu'avec ce choix la règle de Leibniz est satisfaite.
Soient $a \in A^n$ et $x \in M[k]^m = M^{m + k}$. En notant
$\cdot_{M[k]}$ le produit dans $M[k]$, on obtient
\begin{align*}
\text{d}_{M[k]}(a \cdot_{M[k]} x)
& =
(-1)^{k + nk} \text{d}_M(ax) \\
& =
(-1)^{k + nk} \text{d}(a) x + (-1)^{k + nk + n} a \text{d}_M(x) \\ &
=
\text{d}(a) \cdot_{M[k]} x + (-1)^{nk + n} a \text{d}_{M[k]}(x) \\ & =
\text{d}(a)
\cdot_{M[k]} x + (-1)^n a \cdot_{M[k]} \text{d}_{M[k]}(x)
\end{align*}
C'est ce que nous voulons puisque $a$ a un degré $n$ en tant qu'élément homogène de $A$. Nous
observons également qu'avec ces choix, nous pouvons considérer
l'application de multiplication comme un morphisme de complexes
$$
\text{Tot}(A^\bullet \otimes_R M^\bullet[k]) \to
\text{Tot}(A^\bullet \otimes_R M^\bullet)[k] \to
M^\bullet[k]
$$
où la première flèche est
celle de Compléments d'algèbre, section \ref{more-algebra-section-sign-rules}
(\ref{more-algebra-item-shift-tensor}), ce qui dans ce cas implique
exactement le signe que nous avons choisi ci-dessus. (En fait, nous aurions pu déduire que la
règle de Leibniz est respectée à partir de cette observation.)

\medskip\noindent
Avec la règle ci-dessus, nous avons des identifications canoniques
$$
(M[k])^{opp} = M^{opp}[k]
$$
de modules différentiels gradués à droite sur
$A^{opp}$ définis sans l'intervention de signes, en d'autres termes, l'équivalence du
lemme \ref{lemma-left-right} est compatible avec les foncteurs de décalage.

\medskip\noindent
Notre choix ci-dessus nécessite la définition suivante.

\begin{definition}
\label{definition-shift-graded-module}
Soit $R$ un anneau. Soit $A$ une algèbre graduée par $\mathbf{Z}$ sur $R$.
\begin{enumerate}
\item Étant donné un $A$-module gradué à droite $M$, nous définissons le
{\it décalé d'ordre $k$, en tant que $A$-module,} $M[k]$ comme étant
identique en tant que $A$-module à droite mais avec la graduation $(M[k])^n = M^{n + k}$.
\item Étant donné un $A$-module gradué à gauche $M$, nous
définissons le {\it module décalé d'ordre $k$ sur $A$} $M[k]$
comme le module de graduation $(M[k])^n = M^{n + k}$, dont la
multiplication $A^n \times (M[k])^m \to
(M[k])^{n + m}$ est égale à $(-1)^{nk}$ fois la multiplication
donnée $A^n \times M^{m + k} \to M^{n + m + k}$.
\end{enumerate}
\end{definition}

\noindent
Soit $(A, \text{d})$ une algèbre différentielle graduée. Soient
$f, g : M \to N$ des homomorphismes de $A$-modules différentiels gradués à gauche. Une
{\it homotopie entre $f$ et $g$} est une application de $A$-modules
gradués $h : M \to N[-1]$ (observez le décalage !) telle que
$$
f(x) - g(x) = \text{d}_N(h(x)) + h(\text{d}_M(x))
$$
pour tout $x \in M$. S'il existe une homotopie, nous disons alors que $f$ et
$g$ sont {\it homotopes}. Ainsi, $h$ est compatible avec la structure de
$A$-module (avec celle décalée sur $N$) et la graduation (avec la graduation décalée sur
$N$) mais pas avec la différentielle. Si $f = g$ et $h$ est une homotopie, alors
$h$ définit un morphisme $h : M \to N[-1]$ de $A$-modules différentiels
gradués à gauche.

\medskip\noindent
Avec la règle ci-dessus, nous constatons que $f, g : M \to N$ sont homotopes
si et seulement si les morphismes induits
$f^{opp}, g^{opp} : M^{opp} \to N^{opp}$ sont
homotopes en tant qu'homomorphismes de modules différentiels gradués à droite sur
$A^{opp}$ (avec la même homotopie).

\medskip\noindent
La catégorie homotopique, les cônes, les suites exactes courtes
admissibles, les triangles distingués sont tous définis exactement de la même manière
que pour les modules différentiels gradués à droite (et tout concorde sur
les complexes sous-jacents de $R$-modules avec les constructions pour les
complexes de $R$-modules). De cette manière, nous obtenons l'analogue de la
proposition \ref{proposition-homotopy-category-triangulated} pour les
modules à gauche également, ou nous pouvons la déduire en travaillant
avec des modules à droite sur l'algèbre opposée.



\section{Produit tensoriel}
\label{section-tensor-product}

\noindent
Soit $R$ un anneau. Soit $A$ une $R$-algèbre (voir la
section \ref{section-conventions}). Étant donné un $A$-module à droite $M$ et
un $A$-module à gauche $N$, il existe un {\it produit tensoriel}
$$
M \otimes_A N
$$
Ce produit tensoriel est un module sur $R$. Comme $R$-module, $M \otimes_A
N$ est engendré par des symboles $x \otimes y$ avec $x \in M$ et $y \in N$
soumis aux relations
$$
\begin{matrix}
(x_1 + x_2) \otimes y - x_1 \otimes y - x_2 \otimes y, \\
x \otimes (y_1 + y_2) - x \otimes y_1 - x \otimes y_2, \\
xa \otimes y - x \otimes ay
\end{matrix}
$$
pour $a \in A$, $x, x_1, x_2 \in M$ et $y, y_1, y_2 \in N$. Nous
énumérons certaines propriétés du produit tensoriel

\medskip\noindent
Dans chaque variable, le produit tensoriel est exact à droite, en fait il commute
avec les sommes directes et les colimites arbitraires.

\medskip\noindent
Le produit tensoriel $M \otimes_A N$ est le réceptacle de l'application
$A$-bilinéaire universelle $M \times N \to M \otimes_A N$, $(x, y) \mapsto x \otimes y$.
Autrement dit,
$$
\text{Bilinear}_A(M \times N, Q) = \Hom_R(M \otimes_A N, Q)
$$
pour tout $R$-module $Q$.

\medskip\noindent
Si $A$ est une algèbre graduée par $\mathbf{Z}$ et si $M$, $N$ sont des
$A$-modules gradués, alors $M \otimes_A N$ est un $R$-module gradué.
La composante de degré $n$, notée $(M \otimes_A N)^n$, de $M \otimes_A N$
est égale à
$$
\Coker\left(
\bigoplus\nolimits_{r + t + s = n}
M^r \otimes_R A^t \otimes_R N^s \to
\bigoplus\nolimits_{p + q = n} M^p \otimes_R N^q
\right)
$$
où l'application envoie $x \otimes a \otimes y$
vers $x \otimes ay - xa \otimes y$ pour
$x \in M^r$, $y \in N^s$ et $a \in A^t$ avec $r + s + t = n$. Dans
ce cas, l'application $M \times N \to M \otimes_A N$ est $A$-bilinéaire et
compatible avec les graduations et universelle dans le sens que
$$
\text{GradedBilinear}_A(M \times N, Q) =
\Hom_{\text{}R\text{-modules gradués}}(M \otimes_A N, Q)
$$
pour tout $R$-module gradué $Q$ avec une notion évidente d'application
bilinéaire graduée.

\medskip\noindent
Si $(A, \text{d})$ est une algèbre différentielle graduée et
$M$ et $N$ sont des $A$-modules différentiels gradués à gauche et à droite,
alors $M \otimes_A N$ est un $R$-module différentiel gradué avec une différentielle
$$
\text{d}(x \otimes y) =
\text{d}(x) \otimes y + (-1)^{\deg(x)}x \otimes \text{d}(y)
$$
pour $x \in M$ et $y \in N$ homogènes. Dans ce cas, l'application
$M \times N \to M \otimes_A N$ est $A$-bilinéaire, compatible avec les graduations,
et compatible avec les différentielles et universelle dans le sens que
$$
\text{DifferentialGradedBilinear}_A(M \times N, Q) =
\Hom_{\text{Comp}(R)}(M \otimes_A N, Q)
$$
pour tout $R$-module différentiel gradué $Q$ avec une notion évidente
d'application bilinéaire différentielle graduée.






\section{Complexes de Hom et modules différentiels gradués}
\label{section-hom-complexes}

\noindent
Nous encourageons le lecteur à sauter cette section.

\medskip\noindent
Soit $R$ un anneau et soit $M^\bullet$ un complexe de $R$-modules.
Considérons le complexe de $R$-modules
$$
E^\bullet = \Hom^\bullet(M^\bullet, M^\bullet)
$$
introduit dans
Compléments d'algèbre, section \ref{more-algebra-section-hom-complexes}.
D'après Compléments d'algèbre, lemme \ref{more-algebra-lemma-composition} il
existe une loi de composition canonique
$$
\text{Tot}(E^\bullet \otimes_R E^\bullet) \to E^\bullet
$$
qui est un morphisme de complexes. Ainsi, nous voyons que $E^\bullet$ avec
cette multiplication est une $R$-algèbre différentielle graduée que nous
noterons $(E, \text{d})$. De plus, en considérant $M^\bullet$ comme
$\Hom^\bullet(R, M^\bullet)$, nous voyons que la composition définit une multiplication
$$
\text{Tot}(E^\bullet \otimes_R M^\bullet) \to M^\bullet
$$
qui munit $M^\bullet$ d'une structure de $E$-module différentiel gradué à {\bf
gauche}, que nous noterons $M$.

\begin{lemma}
\label{lemma-left-module-structure}
Dans la situation ci-dessus, soit $A$ une $R$-algèbre différentielle graduée.
Donner une structure de $A$-module à gauche sur $M$ revient au même que
donner un homomorphisme $A \to E$ d'algèbres différentielles graduées sur $R$.
\end{lemma}

\begin{proof}
Démonstration omise. Observez qu'aucun signe n'intervient dans cette correspondance.
\end{proof}

\noindent
Nous poursuivons avec la discussion ci-dessus et nous supposons donné un
autre complexe $N^\bullet$ de $R$-modules. Considérons le complexe
de $R$-modules $\Hom^\bullet(M^\bullet, N^\bullet)$ introduit dans
Compléments d'algèbre, section \ref{more-algebra-section-hom-complexes}. Comme
ci-dessus, nous voyons que la composition
$$
\text{Tot}(\Hom^\bullet(M^\bullet, N^\bullet) \otimes_R E^\bullet)
\to \Hom^\bullet(M^\bullet, N^\bullet)
$$
définit une multiplication qui transforme $\Hom^\bullet(M^\bullet, N^\bullet)$
en un $E$-module différentiel gradué à {\bf droite}. En utilisant
le lemme \ref{lemma-left-module-structure},
nous concluons que, pour tout $A$-module différentiel gradué à gauche $M$
et tout complexe de $R$-modules $N^\bullet$, il existe un
$A$-module différentiel gradué canonique à droite dont le complexe
sous-jacent de $R$-modules est $\Hom^\bullet(M^\bullet, N^\bullet)$ et où
la multiplication
$$
\Hom^n(M^\bullet, N^\bullet) \times A^m \longrightarrow
\Hom^{n + m}(M^\bullet, N^\bullet)
$$
envoie $f = (f_{p, q})_{p + q = n}$ avec $f_{p, q} \in \Hom(M^{-q}, N^p)$
et $a \in A^m$ vers l'élément $f \cdot a = (f_{p, q} \circ a)$ où $f_{p,
q} \circ a$ est l'application
$$
M^{-q - m} \xrightarrow{a} M^{-q} \xrightarrow{f_{p, q}} N^p, \quad
x \longmapsto f_{p, q}(ax)
$$
sans l'intervention des signes. Utilisons la notation
$\Hom(M, N^\bullet)$ pour désigner ce $A$-module différentiel gradué à droite.

\begin{lemma}
\label{lemma-characterize-hom}
Soit $R$ un anneau. Soit $(A, \text{d})$ une $R$-algèbre différentielle graduée.
Soit $M'$ un $A$-module différentiel gradué à droite et soit
$M$ un $A$-module différentiel gradué à gauche.
Soit $N^\bullet$ un complexe de $R$-modules. Alors nous avons
$$
\Hom_{\text{Mod}_{(A, d)}}(M', \Hom(M, N^\bullet)) =
\Hom_{\text{Comp}(R)}(M' \otimes_A M, N^\bullet)
$$
où $M \otimes_A M$ est considéré comme un complexe de $R$-modules
comme dans la section \ref{section-tensor-product}.
\end{lemma}

\begin{proof}
Montrons que les deux côtés correspondent à des applications $A$-bilinéaires graduées
$$
M' \times M \longrightarrow N^\bullet
$$
compatibles avec les différentielles. Nous l'avons établi pour le membre de
droite dans la section \ref{section-tensor-product}. Étant donné un élément
$g$ du membre de gauche, l'égalité de
Compléments d'algèbre, lemme \ref{more-algebra-lemma-compose}, détermine
un morphisme de complexes de $R$-modules $g' :
\text{Tot}(M' \otimes_R M) \to N^\bullet$. En d'autres
termes, nous obtenons un morphisme $R$-bilinéaire
gradué $g'' : M' \times M \to N^\bullet$ compatible avec les différentielles. La
$A$-linéarité de $g$ se traduit immédiatement par la
$A$-bilinéarité de $g''$.
\end{proof}

\noindent
Soient $R$, $M^\bullet$, $E^\bullet$, $E$ et $M$ comme ci-dessus.
Cependant, supposons maintenant données une $R$-algèbre différentielle graduée $A$ et
une structure à {\bf droite} de $A$-module différentiel gradué sur $M$. Alors
nous pouvons considérer le morphisme
$$
\text{Tot}(A^\bullet \otimes_R M^\bullet)
\xrightarrow{\psi}
\text{Tot}(A^\bullet \otimes_R M^\bullet)
\to
M^\bullet
$$
où la première flèche est la contrainte de commutativité sur la
catégorie différentielle graduée des complexes de $R$-modules. Cela
correspond à une application
$$
\tau : A^\bullet \longrightarrow E^\bullet
$$
de complexes de $R$-modules. Rappelons que
$E^n = \prod_{p + q = n} \Hom_R(M^{-q}, M^p)$ et
écrivons $\tau(a) = (\tau_{p, q}(a))_{p + q = n}$ pour $a \in A^n$. Alors
nous voyons
$$
\tau_{p, q}(a) : M^{-q} \longrightarrow M^p,\quad
x \longmapsto (-1)^{\deg(a)\deg(x)}x a = (-1)^{-nq}xa
$$
Cela n'est pas compatible avec le produit sur $A$ comme le lecteur devrait s'y
attendre d'après la discussion de la section \ref{section-left-modules}. À
savoir, nous avons
$$
\tau(a a') = (-1)^{\deg(a)\deg(a')}\tau(a') \tau(a)
$$
Nous concluons que le lemme suivant est vrai

\begin{lemma}
\label{lemma-right-module-structure}
Dans la situation ci-dessus, soit $A$ une $R$-algèbre différentielle graduée.
Donner une structure de $A$-module à droite sur $M$ revient au même que
donner un homomorphisme $\tau : A \to E^{opp}$
d'algèbres différentielles graduées sur $R$.
\end{lemma}

\begin{proof}
Voir la discussion ci-dessus et noter que la construction de $\tau$ à partir
de l'application de multiplication $M^n \times A^m \to M^{n + m}$
utilise des signes.
\end{proof}

\noindent
Soient $R$, $M^\bullet$, $E^\bullet$, $E$, $A$ et $M$ comme ci-dessus et
soit une structure de $A$-module différentiel gradué à droite sur $M$
donnée comme dans le lemme. Dans ce cas, il existe un $A$-module différentiel
gradué à gauche canonique dont le complexe sous-jacent de $R$-modules est
$\Hom^\bullet(M^\bullet, N^\bullet)$. À savoir, pour la multiplication, nous pouvons
utiliser
\begin{align*}
\text{Tot}(A^\bullet \otimes_R \Hom^\bullet(M^\bullet, N^\bullet))
& \xrightarrow{\psi}
\text{Tot}(\Hom^\bullet(M^\bullet, N^\bullet) \otimes_R A^\bullet) \\ &
\xrightarrow{\tau}
\text{Tot}(\Hom^\bullet(M^\bullet, N^\bullet) \otimes_R
\Hom^\bullet(M^\bullet, M^\bullet)) \\ &
\to
\text{Tot}(\Hom^\bullet(M^\bullet, N^\bullet)
\end{align*}
La première flèche utilise la contrainte de commutativité sur la catégorie
des complexes de $R$-modules, la deuxième flèche est décrite ci-dessus, et la
troisième flèche est la loi de composition pour le complexe Hom. Chacune de ces
applications est un morphisme de complexes ; leur composé est donc un morphisme
de complexes. En fait, cette construction transforme $\Hom^\bullet(M^\bullet,
N^\bullet)$ en un $A$-module différentiel gradué à gauche (l'associativité de la multiplication
peut être démontrée en utilisant la structure monoïdale symétrique ou par un calcul direct en
utilisant les formules ci-dessous). Explicitons la multiplication
$$
A^n \times \Hom^m(M^\bullet, N^\bullet) \longrightarrow
\Hom^{n + m}(M^\bullet, N^\bullet)
$$
Elle envoie $a \in A^n$
et $f = (f_{p, q})_{p + q = m}$ avec $f_{p, q} \in \Hom(M^{-q}, N^p)$ sur
l'élément $a \cdot f$ dont les composantes sont
$$
(-1)^{nm}f_{p, q} \circ \tau_{-q, q + n}(a) =
(-1)^{nm - n(q + n)}f_{p, q} \circ a =
(-1)^{np + n} f_{p, q} \circ a
$$
dans $\Hom_R(M^{-q - n}, N^p)$ où $f_{p, q} \circ a$ est l'application
$$
M^{-q - n} \xrightarrow{a} M^{-q} \xrightarrow{f_{p, q}} N^p,\quad
x \longmapsto f_{p, q}(xa)
$$
Ici intervient un signe $(-1)^{np + n}$. Utilisons la notation
$\Hom(M, N^\bullet)$ pour désigner ce $A$-module différentiel gradué à gauche.

\begin{lemma}
\label{lemma-characterize-hom-other-side}
Soit $R$ un anneau. Soit $(A, \text{d})$ une $R$-algèbre différentielle graduée.
Soit $M$ un $A$-module différentiel gradué à droite et soit
$M'$ un $A$-module différentiel gradué à gauche.
Soit $N^\bullet$ un complexe de $R$-modules. Alors nous avons
$$
\Hom_{\text{}A\text{-modules différentiels gradués à gauche}}(M', \Hom(M,
N^\bullet)) = \Hom_{\text{Comp}(R)}(M \otimes_A M', N^\bullet)
$$
où $M \otimes_A M'$ est vu comme un complexe de $R$-modules
comme dans la section \ref{section-tensor-product}.
\end{lemma}

\begin{proof}
Montrons que les deux membres correspondent à des applications bilinéaires graduées sur $A$
$$
M \times M' \longrightarrow N^\bullet
$$
compatibles avec les différentielles. Nous avons vu que cela est vrai pour le
membre de droite dans la section \ref{section-tensor-product}. Étant donné un
élément $g$ du membre de gauche, l'égalité de
Compléments d'algèbre, lemme \ref{more-algebra-lemma-compose},
détermine un morphisme de complexes
$g' : \text{Tot}(M' \otimes_R M) \to N^\bullet$. Nous
précomposons avec la contrainte de commutativité pour obtenir
$$
\text{Tot}(M \otimes_R M') \xrightarrow{\psi}
\text{Tot}(M' \otimes_R M) \xrightarrow{g'}
N^\bullet
$$
Ce composé correspond à une application
$R$-bilinéaire graduée $g'' : M \times M' \to N^\bullet$ compatible avec les
différentielles. La $A$-linéarité de $g$ se traduit immédiatement par la $A$-bilinéarité de
$g''$. En effet, soient $x \in M^e$ et $x' \in (M')^{e'}$ et $a \in A^n$.
Alors d'une part nous avons
\begin{align*}
g''(x, ax')
& =
(-1)^{e(n + e')} g'(ax' \otimes x) \\ &
=
(-1)^{e(n + e')} g(ax')(x) \\ &
=
(-1)^{e(n + e')} (a \cdot g(x'))(x) \\ & =
(-1)^{e(n
+ e') + n(n + e + e') + n} g(x')(xa)
\end{align*}
et d'autre part nous avons
$$
g''(xa, x') = (-1)^{(e + n)e'} g'(x' \otimes xa) =
(-1)^{(e + n)e'} g(x')(xa)
$$
Ces deux expressions sont égales, comme le montre un calcul élémentaire modulo $2$ sur les exposants.
\end{proof}

\begin{remark}
\label{remark-evaluation-map-left}
Soit $R$ un anneau. Soit $A$ une $R$-algèbre différentielle graduée.
Soit $M$ un $A$-module différentiel gradué à gauche. Soit
$N^\bullet$ un complexe de $R$-modules. Les constructions ci-dessus produisent
un $A$-module différentiel gradué à droite $\Hom(M, N^\bullet)$, puis un
$A$-module différentiel gradué à gauche $\Hom(\Hom(M,
N^\bullet), N^\bullet)$. Nous affirmons qu'il existe une application
d'évaluation
$$
ev : M \longrightarrow \Hom(\Hom(M, N^\bullet), N^\bullet)
$$
dans la catégorie des $A$-modules différentiels gradués à gauche. Pour la définir,
d'après le lemme \ref{lemma-characterize-hom}, il suffit de construire un
couplage $A$-bilinéaire
$$
\Hom(M, N^\bullet) \times M \longrightarrow N^\bullet
$$
compatible avec la graduation et les différentielles. Pour cela nous prenons
$$
(f, x) \longmapsto f(x)
$$
Nous laissons au lecteur le soin de vérifier que ceci est compatible avec la graduation,
les différentielles, et bilinéaire sur $A$. L'application $ev$ sur les complexes de
$R$-modules sous-jacents est celle de Compléments d'algèbre, point (\ref{more-algebra-item-evaluation}).
\end{remark}

\begin{remark}
\label{remark-evaluation-map-right}
Soit $R$ un anneau. Soit $A$ une $R$-algèbre différentielle graduée.
Soit $M$ un $A$-module différentiel gradué à droite. Soit
$N^\bullet$ un complexe de $R$-modules. Les constructions ci-dessus produisent
un $A$-module différentiel gradué à gauche $\Hom(M, N^\bullet)$, puis un
$A$-module différentiel gradué à droite $\Hom(\Hom(M,
N^\bullet), N^\bullet)$. Nous affirmons qu'il existe une application d'évaluation
$$
ev : M \longrightarrow \Hom(\Hom(M, N^\bullet), N^\bullet)
$$
dans la catégorie des $A$-modules différentiels gradués à droite. Pour la définir,
d'après le lemme \ref{lemma-characterize-hom}, il suffit de construire un
couplage $A$-bilinéaire
$$
M \times \Hom(M, N^\bullet) \longrightarrow N^\bullet
$$
compatible avec la graduation et les différentielles. Pour cela, nous prenons
$$
(x, f) \longmapsto (-1)^{\deg(x)\deg(f)}f(x)
$$
Nous laissons au lecteur le soin de vérifier que cela est compatible avec la
graduation, les différentielles et bilinéaire sur $A$. L'application $ev$ sur les complexes de
$R$-modules sous-jacents est celle de Compléments d'algèbre, point (\ref{more-algebra-item-evaluation}).
\end{remark}

\begin{remark}
\label{remark-shift-dual}
Soit $R$ un anneau. Soit $A$ une algèbre différentielle graduée sur
$R$. Soient $M^\bullet$ et $N^\bullet$ des complexes de
$R$-modules. Soit $k \in \mathbf{Z}$ et considérons l'isomorphisme
$$
\Hom^\bullet(M^\bullet, N^\bullet)[-k]
\longrightarrow
\Hom^\bullet(M^\bullet[k], N^\bullet)
$$
de complexes de $R$-modules défini dans
Compléments d'algèbre, point (\ref{more-algebra-item-shift-hom}). Si
$M^\bullet$ a la structure d'un module différentiel gradué à gauche, resp.\ à droite,
sur $A$, alors il s'agit d'un homomorphisme de modules
différentiels gradués à droite, resp.\ à gauche, sur $A$ (avec les
structures de module telles que définies dans cette section).
Nous omettons la vérification ; nous avertissons le lecteur que la
structure de $A$-module sur le décalage d'un $A$-module gradué à gauche est
définie en utilisant un signe,
voir définition \ref{definition-shift-graded-module}.
\end{remark}






\section{Modules projectifs sur les algèbres}
\label{section-projectives-over-algebras}

\noindent
Dans cette section, nous étudions les modules projectifs sur des algèbres, par
analogie avec Algèbre, section
\ref{algebra-section-projective}. Cette section devrait probablement être déplacée ailleurs.

\medskip\noindent
Soit $R$ un anneau et soit $A$ une $R$-algèbre, voir la
section \ref{section-conventions} pour nos conventions. Il
est clair que $A$ est un $A$-module à droite projectif puisque
$\Hom_A(A, M) = M$ pour tout $A$-module à droite $M$ (et donc $\Hom_A(A, -)$ est
exact). Inversement, soit $P$ un $A$-module à droite projectif. Alors nous
pouvons choisir une surjection
$\bigoplus_{i \in I} A \to P$ en choisissant un ensemble $\{p_i\}_{i \in I}$
de générateurs de $P$ sur $A$. Puisque $P$ est projectif, il existe un
inverse à gauche de la surjection, et nous trouvons que $P$ est isomorphe à
un facteur direct d'un module libre, exactement comme dans le cas commutatif
(Algèbre, lemme \ref{algebra-lemma-characterize-projective}).

\medskip\noindent
Nous concluons
\begin{enumerate}
\item la catégorie des $A$-modules a suffisamment de projectifs,
\item $A$ est un $A$-module projectif.
\item chaque $A$-module est un quotient d'une somme directe de copies de $A$,
\item chaque $A$-module projectif est un facteur direct d'une somme
directe de copies de $A$.
\end{enumerate}



\section{Modules projectifs sur les algèbres graduées}
\label{section-projectives-over-graded-algebras}

\noindent
Dans cette section, nous étudions les modules gradués projectifs sur des algèbres
graduées, par analogie avec Algèbre, section \ref{algebra-section-projective}.

\medskip\noindent
Soit $R$ un anneau. Soit $A$ une algèbre $\mathbf{Z}$-graduée sur $R$.
Voir la section \ref{section-conventions} pour nos
conventions. Soit $\text{Mod}_A$ la catégorie des $A$-modules gradués à droite.
Pour un entier $k$, soit $A[k]$ le décalage de $A$. Pour un
$A$-module gradué à droite, nous avons
$$
\Hom_{\text{Mod}_A}(A[k], M) = M^{-k}
$$
Comme le foncteur $M \mapsto M^{-k}$ est exact sur $\text{Mod}_A$, nous
en concluons que $A[k]$ est un objet projectif de $\text{Mod}_A$.
Inversement, supposons que $P$ soit un objet projectif de $\text{Mod}_A$. En
choisissant un ensemble de générateurs homogènes de $P$ comme $A$-module, nous
pouvons trouver une surjection
$$
\bigoplus\nolimits_{i \in I} A[k_i] \longrightarrow P
$$
Ainsi, nous concluons qu'un objet projectif de $\text{Mod}_A$ est
un facteur direct d'une somme directe des décalages $A[k]$.

\medskip\noindent
Nous concluons
\begin{enumerate}
\item la catégorie des $A$-modules gradués a suffisamment d'objets projectifs,
\item $A[k]$ est un $A$-module projectif pour chaque $k \in \mathbf{Z}$,
\item chaque $A$-module gradué est un quotient d'une somme directe
de copies des modules $A[k]$ pour différents $k$.
\item tout $A$-module projectif est un facteur direct d'une somme
directe de copies des modules $A[k]$ pour des $k$ variables.
\end{enumerate}




\section{Modules projectifs et algèbres différentielles graduées}
\label{section-projective-over-differential-graded}

\noindent
Si $(A, \text{d})$ est une algèbre différentielle graduée et $P$ est
un objet de $\text{Mod}_{(A, \text{d})}$ alors nous disons
{\it $P$ est projectif en tant que $A$-module gradué} ou parfois
{\it $P$ est projectif gradué} pour signifier que $P$
est un objet projectif de la catégorie abélienne $\text{Mod}_A$ des
$A$-modules gradués comme dans
la section \ref{section-projectives-over-graded-algebras}.

\begin{lemma}
\label{lemma-target-graded-projective}
Soit $(A, \text{d})$ une algèbre différentielle graduée.
Soit $M \to P$ un homomorphisme surjectif de $A$-modules différentiels
gradués. Si $P$ est projectif en tant que $A$-module gradué, alors
$M \to P$ est un épimorphisme admissible.
\end{lemma}

\begin{proof}
Cela résulte immédiatement des définitions.
\end{proof}

\begin{lemma}
\label{lemma-hom-from-shift-free}
Soit $(A, d)$ une algèbre différentielle graduée. Alors nous avons
$$
\Hom_{\text{Mod}_{(A, \text{d})}}(A[k], M) =
\Ker(\text{d} : M^{-k} \to M^{-k + 1})
$$
et
$$
\Hom_{K(\text{Mod}_{(A, \text{d})})}(A[k], M) = H^{-k}(M)
$$
pour tout $A$-module différentiel gradué $M$.
\end{lemma}

\begin{proof}
Immédiat d'après les définitions.
\end{proof}







\section{Modules injectifs sur les algèbres}
\label{section-modules-noncommutative}

\noindent
Dans cette section, nous étudions les modules injectifs sur des
algèbres, par
analogie avec Compléments d'algèbre, section
\ref{more-algebra-section-injectives-modules}. Cette section devrait probablement être déplacée ailleurs.

\medskip\noindent
Soit $R$ un anneau et soit $A$ une
$R$-algèbre, voir la section \ref{section-conventions} pour nos conventions. Pour un
$A$-module à droite $M$, nous posons
$$
M^\vee = \Hom_\mathbf{Z}(M, \mathbf{Q}/\mathbf{Z})
$$
que nous considérons comme un $A$-module à gauche par la
multiplication $(a f)(x) = f(xa)$. À savoir, $((ab)f)(x) = f(xab) = (bf)(xa) = (a(bf))(x)$.
Inversement, si $M$ est un $A$-module à gauche, alors $M^\vee$ est un
$A$-module à droite. Puisque $\mathbf{Q}/\mathbf{Z}$ est un groupe abélien
injectif (Compléments d'algèbre, lemme \ref{more-algebra-lemma-injective-abelian}), le
foncteur $M \mapsto M^\vee$ est exact
(Compléments d'algèbre, lemme \ref{more-algebra-lemma-vee-exact}). De plus,
l'application d'évaluation $M \to (M^\vee)^\vee$ est
injective pour tous les modules $M$
(Compléments d'algèbre, lemme \ref{more-algebra-lemma-ev-injective}).

\medskip\noindent
Nous affirmons que $A^\vee$ est un $A$-module à droite injectif. À savoir, étant donné
un $A$-module à droite $N$ nous avons
$$
\Hom_A(N, A^\vee) =
\Hom_A(N, \Hom_\mathbf{Z}(A, \mathbf{Q}/\mathbf{Z})) = N^\vee
$$
et nous concluons parce que le foncteur $N \mapsto N^\vee$ est
exact. La seconde égalité tient parce que
$$
\Hom_\mathbf{Z}(N, \Hom_\mathbf{Z}(A, \mathbf{Q}/\mathbf{Z})) =
\Hom_\mathbf{Z}(N \otimes_\mathbf{Z} A, \mathbf{Q}/\mathbf{Z})
$$
d'après Algèbre, lemme \ref{algebra-lemma-hom-from-tensor-product}. À
l'intérieur de ce module, la $A$-linéarité sélectionne exactement les applications
bilinéaires $\varphi : N \times A \to \mathbf{Q}/\mathbf{Z}$ qui
ont la même valeur sur $x \otimes a$ et $xa \otimes 1$, c'est-à-dire qui
proviennent d'éléments de $N^\vee$.

\medskip\noindent
Enfin, pour tout $A$-module à droite $M$, nous pouvons choisir une
surjection $\bigoplus_{i \in I} A \to M^\vee$ pour obtenir une
injection $M \to (M^\vee)^\vee \to \prod_{i \in I} A^\vee$.

\medskip\noindent
Nous concluons
\begin{enumerate}
\item que la catégorie des $A$-modules possède suffisamment d'injectifs,
\item $A^\vee$ est un $A$-module injectif, et
\item tout $A$-module s'injecte dans un produit de copies de $A^\vee$.
\end{enumerate}





\section{Modules injectifs sur les algèbres graduées}
\label{section-modules-noncommutative-graded}

\noindent
Dans cette section, nous discutons des modules gradués injectifs sur les algèbres
graduées de manière
analogue à Compléments d'algèbre, section \ref{more-algebra-section-injectives-modules}.

\medskip\noindent
Soit $R$ un anneau. Soit $A$ une algèbre graduée par $\mathbf{Z}$ sur $R$. Voir
la section \ref{section-conventions} pour nos conventions. Si
$M$ est un $R$-module gradué, nous posons
$$
M^\vee =
\bigoplus\nolimits_{n \in \mathbf{Z}}
\Hom_\mathbf{Z}(M^{-n}, \mathbf{Q}/\mathbf{Z}) =
\bigoplus\nolimits_{n \in \mathbf{Z}} (M^{-n})^\vee
$$
comme un $R$-module gradué (aucun signe dans les actions de $R$ sur
les parties homogènes). Si $M$ a la structure d'un $A$-module à gauche
gradué, alors nous définissons une structure de $A$-module gradué à droite
sur $M^\vee$ en laissant $a \in A^m$ agir par
$$
(M^{-n})^\vee \to (M^{-n - m})^\vee, \quad
f \mapsto f \circ a
$$
comme dans la section
\ref{section-hom-complexes}. Si $M$ a la structure d'un $A$-module à droite
gradué, alors nous définissons une structure de $A$-module gradué à
gauche sur $M^\vee$ en laissant $a \in A^n$ agir par
$$
(M^{-m})^\vee \to (M^{-m - n})^\vee, \quad
f \mapsto (-1)^{nm}f \circ a
$$
comme dans la section \ref{section-hom-complexes} (le signe nous est imposé
parce que nous voulons utiliser la même formule pour le cas où
l'on travaille avec des modules différentiels gradués --- si l'on ne
s'intéresse qu'aux modules gradués, on peut omettre le signe ici). Sur la
catégorie des $A$-modules gradués (à gauche ou à droite), le foncteur $M
\mapsto M^\vee$ est exact (vérifier sur les composantes graduées). De
plus, il existe une application d'évaluation injective
$$
ev : M \longrightarrow (M^\vee)^\vee, \quad
ev^n = (-1)^n \text{ l'application d'évaluation }M^n \to ((M^n)^\vee)^\vee
$$
de modules gradués sur $R$,
voir Compléments d'algèbre, point (\ref{more-algebra-item-evaluation}).
Cette application d'évaluation est un homomorphisme de $A$-modules à gauche, resp.\ à droite,
si $M$ est un $A$-module à gauche, resp.\ à
droite ; voir les remarques \ref{remark-evaluation-map-left} et
\ref{remark-evaluation-map-right}. Enfin, étant donné $k \in \mathbf{Z}$, il existe un isomorphisme canonique
$$
M^\vee[-k] \longrightarrow (M[k])^\vee
$$
de $R$-modules gradués qui utilise un signe et qui, si $M$
est un $A$-module à gauche, resp. \ à droite, est un isomorphisme de
$A$-modules à droite, resp. \ à gauche. Voir la remarque \ref{remark-shift-dual}.

\medskip\noindent
Nous affirmons que $A^\vee$ est un objet injectif de la catégorie
$\text{Mod}_A$ des $A$-modules gradués à droite. En effet, étant donné un
$A$-module gradué à droite $N$ nous avons
$$
\Hom_{\text{Mod}_A}(N, A^\vee) =
\Hom_{\text{Comp}(\mathbf{Z})}(N \otimes_A A, \mathbf{Q}/\mathbf{Z})) =
(N^0)^\vee
$$
selon le lemme
\ref{lemma-characterize-hom} (appliqué au cas où toutes les différentielles sont nulles).
Nous en concluons parce que le foncteur $N \mapsto (N^0)^\vee = (N^\vee)^0$
est exact.

\medskip\noindent
Enfin, pour chaque $A$-module gradué à droite $M$ nous pouvons choisir une surjection
de $A$-modules gradués à gauche
$$
\bigoplus\nolimits_{i \in I} A[k_i] \to M^\vee
$$
où $A[k_i]$ désigne le décalage de $A$ par $k_i \in \mathbf{Z}$. Nous
faisons cela en choisissant des générateurs homogènes pour $M^\vee$. De cette
manière, nous obtenons une injection
$$
M \to (M^\vee)^\vee \to \prod A[k_i]^\vee = \prod A^\vee[-k_i]
$$
Remarquez que les produits dans la formule ci-dessus sont des produits
dans la catégorie des modules gradués (autrement dit, prendre les produits degré
par degré, puis la somme directe des composantes).

\medskip\noindent
Nous concluons que
\begin{enumerate}
\item la catégorie des $A$-modules gradués a assez d'injectifs,
\item pour chaque $k \in \mathbf{Z}$ le module $A^\vee[k]$ est injectif, et
\item chaque $A$-module s'injecte dans un produit dans la catégorie des
modules gradués de copies de décalages $A^\vee[k]$.
\end{enumerate}





\section{Modules injectifs et algèbres différentielles graduées}
\label{section-modules-noncommutative-differential-graded}

\noindent
Si $(A, \text{d})$ est une algèbre différentielle graduée et $I$ est
un objet de $\text{Mod}_{(A, \text{d})}$ alors nous disons
que {\it $I$ est injectif en tant que $A$-module gradué} ou
parfois {\it $I$ est injectif gradué} pour
signifier que $I$ est un objet injectif de la catégorie abélienne $\text{Mod}_A$ des
$A$-modules gradués.

\begin{lemma}
\label{lemma-source-graded-injective}
Soit $(A, \text{d})$ une algèbre différentielle
graduée. Soit $I \to M$ un homomorphisme injectif de $A$-modules
différentiels gradués. Si $I$ est injectif gradué, alors $I
\to M$ est un monomorphisme admissible.
\end{lemma}

\begin{proof}
Cela résulte immédiatement des définitions.
\end{proof}

\noindent
Soit $(A, \text{d})$ une algèbre différentielle graduée. Si $M$ est un
$A$-module différentiel gradué à gauche, resp.\ à droite, alors
$$
M^\vee = \Hom^\bullet(M^\bullet, \mathbf{Q}/\mathbf{Z})
$$
avec la structure de $A$-module construite
dans la section \ref{section-modules-noncommutative-graded} est un
$A$-module différentiel gradué à droite, resp.\ à gauche,
selon la discussion dans la section
\ref{section-hom-complexes}. Selon les
remarques \ref{remark-evaluation-map-left} et \ref{remark-evaluation-map-right},
l'application d'évaluation de la section \ref{section-modules-noncommutative-graded}
$$
M \longrightarrow (M^\vee)^\vee
$$
est un homomorphisme de $A$-modules différentiels gradués à gauche, resp.\ à droite.

\begin{lemma}
\label{lemma-map-into-dual}
Soit $(A, \text{d})$ une algèbre différentielle graduée. Si
$M$ est un $A$-module différentiel gradué à gauche et $N$ est un
$A$-module différentiel gradué à droite, alors
\begin{align*}
\Hom_{\text{Mod}_{(A, \text{d})}}(N, M^\vee)
& =
\Hom_{\text{Comp}(\mathbf{Z})}(N \otimes_A M, \mathbf{Q}/\mathbf{Z}) \\ &
=
\text{DifferentialGradedBilinear}_A(N \times M, \mathbf{Q}/\mathbf{Z})
\end{align*}
\end{lemma}

\begin{proof}
La première égalité découle du lemme
\ref{lemma-characterize-hom} et la deuxième égalité a été démontrée dans la section \ref{section-tensor-product}.
\end{proof}

\begin{lemma}
\label{lemma-hom-into-shift-dual-free}
Soit $(A, \text{d})$ une algèbre différentielle graduée. Alors nous avons
$$
\Hom_{\text{Mod}_{(A, \text{d})}}(M, A^\vee[k]) =
\Ker(\text{d} : (M^\vee)^k \to (M^\vee)^{k + 1})
$$
et
$$
\Hom_{K(\text{Mod}_{(A, \text{d})})}(M, A^\vee[k]) = H^k(M^\vee)
$$
comme foncteurs du $A$-module différentiel gradué $M$.
\end{lemma}

\begin{proof}
Cela est clair à partir de la discussion ci-dessus.
\end{proof}















\section{P-résolutions}
\label{section-P-resolutions}

\noindent
Cette section est l'analogue de
Catégories dérivées, section \ref{derived-section-unbounded}.

\medskip\noindent
Soit $(A, \text{d})$ une algèbre différentielle graduée.
Soit $P$ un $A$-module différentiel gradué. Nous disons
que $P$ {\it a la propriété (P)} s'il existe une filtration
$$
0 = F_{-1}P \subset F_0P \subset F_1P \subset \ldots \subset P
$$
par des sous-modules différentiels gradués, satisfaisant aux conditions suivantes :
\begin{enumerate}
\item $P = \bigcup F_pP$,
\item les inclusions $F_iP \to F_{i + 1}P$ sont des
monomorphismes admissibles,
\item les quotients $F_{i + 1}P/F_iP$ sont isomorphes en tant que modules
différentiels gradués sur $A$ à une somme directe de $A[k]$.
\end{enumerate}
En fait, la condition (2) est une conséquence de la condition (3),
voir le lemme \ref{lemma-target-graded-projective}. De plus, le lecteur
peut vérifier que, en tant que $A$-module gradué, $P$ sera isomorphe à une
somme directe de décalages de $A$.

\begin{lemma}
\label{lemma-property-P-sequence}
Soit $(A, \text{d})$ une algèbre différentielle graduée.
Soit $P$ un module différentiel gradué sur $A$. Si $F_\bullet$ est une
filtration comme dans la propriété (P), alors nous obtenons
une suite exacte courte admissible
$$
0 \to
\bigoplus\nolimits F_iP \to
\bigoplus\nolimits F_iP \to P \to 0
$$
de modules différentiels gradués sur $A$.
\end{lemma}

\begin{proof}
La deuxième application est la somme directe des applications
d'inclusion. La première application sur le facteur direct $F_iP$ de la source
est la somme de l'identité $F_iP \to F_iP$ et de l'opposée de l'application
d'inclusion $F_iP \to F_{i + 1}P$. Choisissons des homomorphismes $s_i : F_{i + 1}P \to F_iP$ de
modules gradués sur $A$ qui sont des inverses à gauche des applications
d'inclusion. La composition donne des applications $s_{j, i} : F_jP \to F_iP$ pour tout
$j > i$. Ensuite, un inverse à gauche de la première flèche envoie $x \in
F_jP$ sur $(s_{j, 0}(x), s_{j, 1}(x), \ldots, s_{j, j - 1}(x), 0, \ldots)$ dans
$\bigoplus F_iP$.
\end{proof}

\noindent
Le lemme suivant montre que les modules différentiels gradués ayant la
propriété (P) sont la notion duale des modules K-injectifs
(c'est-à-dire qu'ils sont K-projectifs dans un certain
sens). Voir Catégories dérivées, définition \ref{derived-definition-K-injective}.

\begin{lemma}
\label{lemma-property-P-K-projective}
Soit $(A, \text{d})$ une algèbre différentielle graduée.
Soit $P$ un $A$-module différentiel gradué ayant la propriété (P).
Alors
$$
\Hom_{K(\text{Mod}_{(A, \text{d})})}(P, N) = 0
$$
pour tous les $A$-modules différentiels gradués acycliques $N$.
\end{lemma}

\begin{proof}
Nous allons utiliser que $K(\text{Mod}_{(A, \text{d})})$ est une
catégorie triangulée (proposition
\ref{proposition-homotopy-category-triangulated}). Soit $F_\bullet$ une filtration sur $P$ comme dans la
propriété (P). La suite exacte courte du lemme
\ref{lemma-property-P-sequence} produit un triangle distingué. Donc, d'après
Catégories dérivées, lemme \ref{derived-lemma-representable-homological}, il
suffit de montrer que
$$
\Hom_{K(\text{Mod}_{(A, \text{d})})}(F_iP, N) = 0
$$
pour tous les $A$-modules différentiels gradués acycliques $N$ et tous les
$i$. Chacun des modules différentiels gradués $F_iP$ possède une filtration finie par des
monomorphismes admissibles, dont les quotients successifs sont des sommes
directes de décalages $A[k]$. Ainsi, il suffit de prouver que
$$
\Hom_{K(\text{Mod}_{(A, \text{d})})}(A[k], N) = 0
$$
pour tous les $A$-modules différentiels gradués acycliques $N$ et tous
les $k$. Cela découle du lemme \ref{lemma-hom-from-shift-free}.
\end{proof}

\begin{lemma}
\label{lemma-good-quotient}
Soit $(A, \text{d})$ une algèbre différentielle
graduée. Soit $M$ un $A$-module différentiel gradué. Il existe un homomorphisme
$P \to M$ de $A$-modules différentiels gradués avec les propriétés
suivantes
\begin{enumerate}
\item $P \to M$ est surjectif,
\item $\Ker(\text{d}_P) \to \Ker(\text{d}_M)$ est surjectif, et
\item $P$ s'insère dans une suite exacte courte admissible
$0 \to P' \to P \to P'' \to 0$ où $P'$, $P''$ sont des sommes directes
de décalages de $A$.
\end{enumerate}
\end{lemma}

\begin{proof}
Soit $P_k$ le $A$-module libre avec des générateurs $x, y$ en degrés $k$ et
$k + 1$. Définissons la structure de $A$-module différentiel gradué sur
$P_k$ en posant $\text{d}(x) = y$ et $\text{d}(y) = 0$. Pour chaque élément $m
\in M^k$, il existe un homomorphisme $P_k \to M$ envoyant $x$
vers $m$ et $y$ vers $\text{d}(m)$. Ainsi, nous voyons qu'il existe
une surjection d'une somme directe de copies de $P_k$ vers $M$.
Cela produit clairement $P \to M$ ayant les propriétés (1) et (3).
Pour obtenir la propriété (2), remarquons que si $m \in
\Ker(\text{d}_M)$ a le degré $k$, alors il existe une application $A[k] \to M$ envoyant
$1$ vers $m$. Par conséquent, nous pouvons atteindre (2) en ajoutant une somme
directe de copies de décalages de $A$.
\end{proof}

\begin{lemma}
\label{lemma-resolve}
Soit $(A, \text{d})$ une algèbre différentielle
graduée. Soit $M$ un $A$-module différentiel gradué. Il existe un homomorphisme
$P \to M$ de $A$-modules différentiels gradués tel que
\begin{enumerate}
\item $P \to M$ soit un quasi-isomorphisme, et
\item $P$ ait la propriété (P).
\end{enumerate}
\end{lemma}

\begin{proof}
Posons $M = M_0$. Nous choisissons par récurrence des suites exactes courtes
$$
0 \to M_{i + 1} \to P_i \to M_i \to 0
$$
où les applications $P_i \to M_i$ sont choisies comme dans le lemme
\ref{lemma-good-quotient}. Cela donne une « résolution »
$$
\ldots \to P_2 \xrightarrow{f_2} P_1 \xrightarrow{f_1} P_0 \to M \to 0
$$
Ensuite, nous posons
$$
P = \bigoplus\nolimits_{i \geq 0} P_i
$$
comme $A$-module, avec une graduation
donnée par $P^n = \bigoplus_{a + b = n} P_{-a}^b$
et une différentielle (comme dans la construction du complexe total associé à
un double complexe) par
$$
\text{d}_P(x) = f_{-a}(x) + (-1)^a \text{d}_{P_{-a}}(x)
$$
pour $x \in P_{-a}^b$. Avec ces conventions, $P$ est en effet un $A$-module
différentiel gradué. En rappelant que chaque $P_i$ a une filtration à deux crans $0
\to P_i' \to P_i \to P_i'' \to 0$, nous posons
$$
F_{2i}P = \bigoplus\nolimits_{i \geq j \geq 0} P_j
\subset
\bigoplus\nolimits_{i \geq 0} P_i = P
$$
et nous ajoutons $P'_{i + 1}$ à $F_{2i}P$ pour obtenir $F_{2i +
1}$. Ce sont des sous-modules différentiels gradués et les quotients successifs
sont des sommes directes de décalages de
$A$. D'après le lemme \ref{lemma-target-graded-projective},
nous voyons que les inclusions $F_iP \to F_{i + 1}P$ sont des monomorphismes
admissibles. Enfin, nous devons montrer que l'application $P \to M$
(donnée par l'augmentation $P_0 \to M$) est un quasi-isomorphisme. Cela découle
d'Homologie, lemme \ref{homology-lemma-good-resolution-gives-qis}.
\end{proof}





\section{I-résolutions}
\label{section-I-resolutions}

\noindent
Cette section est le dual de la section sur les P-résolutions.

\medskip\noindent
Soit $(A, \text{d})$ une algèbre différentielle graduée.
Soit $I$ un $A$-module différentiel gradué. Nous disons
que $I$ {\it a la propriété (I)} s'il existe une filtration
$$
I = F_0I \supset F_1I \supset F_2I \supset \ldots \supset 0
$$
par sous-modules différentiels gradués telle que
\begin{enumerate}
\item $I = \lim I/F_pI$,
\item les applications $I/F_{i + 1}I \to I/F_iI$ soient des épimorphismes admissibles,
\item les quotients $F_iI/F_{i + 1}I$ soient isomorphes comme
$A$-modules différentiels gradués à des produits des modules $A^\vee[k]$ construits
à la section \ref{section-modules-noncommutative-differential-graded}.
\end{enumerate}
En fait, la condition (2) est une conséquence de la condition (3),
voir le lemme \ref{lemma-source-graded-injective}. Le lecteur peut vérifier
qu'en tant que module gradué, $I$ sera isomorphe à un produit de $A^\vee[k]$.

\begin{lemma}
\label{lemma-property-I-sequence}
Soit $(A, \text{d})$ une algèbre différentielle graduée.
Soit $I$ un $A$-module différentiel gradué. Si $F_\bullet$ est une
filtration comme dans la propriété (I), alors on obtient
une suite exacte courte admissible
$$
0 \to I \to
\prod\nolimits I/F_iI \to
\prod\nolimits I/F_iI \to 0
$$
de $A$-modules différentiels gradués.
\end{lemma}

\begin{proof}
Démonstration omise. Indication : c'est le dual du lemme \ref{lemma-property-P-sequence}.
\end{proof}

\noindent
Le lemme suivant montre que les modules différentiels gradués avec
la propriété (I) sont l'analogue des modules K-injectifs.
Voir Catégories dérivées, définition \ref{derived-definition-K-injective}.

\begin{lemma}
\label{lemma-property-I-K-injective}
Soit $(A, \text{d})$ une algèbre différentielle graduée.
Soit $I$ un $A$-module différentiel gradué avec la propriété (I).
Alors
$$
\Hom_{K(\text{Mod}_{(A, \text{d})})}(N, I) = 0
$$
pour tous les $A$-modules différentiels gradués acycliques $N$.
\end{lemma}

\begin{proof}
Nous utiliserons que $K(\text{Mod}_{(A, \text{d})})$ est une
catégorie triangulée (proposition \ref{proposition-homotopy-category-triangulated}).
Soit $F_\bullet$ une filtration sur $I$ comme dans la
propriété (I). La suite exacte courte du lemme
\ref{lemma-property-I-sequence} produit un triangle distingué. Ainsi, par
Catégories dérivées, lemme \ref{derived-lemma-representable-homological}, il
suffit de montrer que
$$
\Hom_{K(\text{Mod}_{(A, \text{d})})}(N, I/F_iI) = 0
$$
pour tous les $A$-modules différentiels gradués acycliques $N$ et tous les
$i$. Chacun des modules différentiels gradués $I/F_iI$ possède une filtration finie
par monomorphismes admissibles, dont les quotients successifs
sont des produits de $A^\vee[k]$. Ainsi il suffit de prouver que
$$
\Hom_{K(\text{Mod}_{(A, \text{d})})}(N, A^\vee[k]) = 0
$$
pour tous les $A$-modules différentiels gradués acycliques $N$ et
tous les $k$. Cela découle du lemme
\ref{lemma-hom-into-shift-dual-free} et du fait que $(-)^\vee$ est un foncteur exact.
\end{proof}

\begin{lemma}
\label{lemma-good-sub}
Soit $(A, \text{d})$ une algèbre différentielle
graduée. Soit $M$ un $A$-module différentiel gradué. Il existe un homomorphisme
$M \to I$ de $A$-modules différentiels gradués avec les propriétés
suivantes
\begin{enumerate}
\item $M \to I$ est injectif,
\item $\Coker(\text{d}_M) \to \Coker(\text{d}_I)$ est injectif,
et
\item $I$ s'insère dans une suite exacte courte
admissible $0 \to I' \to I \to I'' \to 0$ où $I'$, $I''$ sont des produits
de décalages de $A^\vee$.
\end{enumerate}
\end{lemma}

\begin{proof}
Nous utiliserons les foncteurs $N \mapsto N^\vee$ (des modules
différentiels gradués à gauche vers ceux à droite, et des modules différentiels
gradués à droite vers ceux à gauche)
construits dans la section
\ref{section-modules-noncommutative-differential-graded} et toutes leurs propriétés.
Pour chaque $k \in \mathbf{Z}$, soit $Q_k$ le $A$-module libre à gauche avec des
générateurs $x, y$ dans les degrés $k$ et $k + 1$. Définissons la structure de $A$-module
différentiel gradué à gauche sur $Q_k$ en posant $\text{d}(x) = y$ et $\text{d}(y) =
0$. En raisonnant exactement comme dans la démonstration du lemme
\ref{lemma-good-quotient}, nous trouvons une surjection
$$
\bigoplus\nolimits_{i \in I} Q_{k_i} \longrightarrow M^\vee
$$
de modules différentiels gradués à gauche sur $A$. Ensuite, nous pouvons considérer l'injection
$$
M \to (M^\vee)^\vee \to (\bigoplus\nolimits_{i \in I} Q_{k_i})^\vee =
\prod\nolimits_{i \in I} I_{k_i}
$$
où $I_k = Q_{-k}^\vee$ est le $A$-module différentiel gradué à droite « dual ».
De plus, la suite exacte courte $0 \to A[-k - 1] \to Q_k \to A[-k] \to 0$
produit une suite exacte courte
$0 \to A^\vee[k] \to I_k \to A^\vee[k + 1] \to 0$.

\medskip\noindent
Le résultat du paragraphe précédent produit $M \to I$ ayant
les propriétés (1) et (3). Pour obtenir la propriété (2), supposons que
$\overline{m} \in \Coker(\text{d}_M)$ soit un élément non nul de degré
$k$. Choisissons une application $\lambda : M^k \to \mathbf{Q}/\mathbf{Z}$
qui s'annule sur $\Im(M^{k - 1} \to M^k)$ mais pas sur $m$. Selon le
lemme \ref{lemma-hom-into-shift-dual-free}, cela
correspond à un homomorphisme $M \to A^\vee[k]$ de $A$-modules
différentiels gradués qui ne s'annule pas sur $m$. Par conséquent,
nous pouvons obtenir (2) en ajoutant un
produit de copies de décalages de $A^\vee$.
\end{proof}

\begin{lemma}
\label{lemma-right-resolution}
Soit $(A, \text{d})$ une algèbre différentielle
graduée. Soit $M$ un $A$-module différentiel gradué. Il existe un homomorphisme
$M \to I$ de $A$-modules différentiels gradués tel que
\begin{enumerate}
\item $M \to I$ soit un quasi-isomorphisme, et
\item $I$ possède la propriété (I).
\end{enumerate}
\end{lemma}

\begin{proof}
Posons $M = M_0$. Nous choisissons par récurrence des suites exactes courtes
$$
0 \to M_i \to I_i \to M_{i + 1} \to 0
$$
où les applications $M_i \to I_i$ sont choisies comme dans le lemme
\ref{lemma-good-sub}. Cela donne une « résolution »
$$
0 \to M \to I_0 \xrightarrow{f_0} I_1 \xrightarrow{f_1} I_1 \to \ldots
$$
Notons $I$ le $A$-module différentiel gradué dont les composantes graduées sont
$$
I^n = \prod\nolimits_{i \geq 0} I^{n - i}_i
$$
et dont la différentielle est définie par
$$
\text{d}_I(x) = f_i(x) + (-1)^i \text{d}_{I_i}(x)
$$
pour $x \in I_i^{n - i}$. Avec ces conventions, $I$ est en effet un $A$-module
différentiel gradué. En rappelant que chaque $I_i$ possède une filtration à deux crans
$0 \to I_i' \to I_i \to I_i'' \to 0$, nous posons
$$
F_{2i}I^n = \prod\nolimits_{j \geq i} I^{n - j}_j
\subset
\prod\nolimits_{i \geq 0} I^{n - i}_i = I^n
$$
et nous ajoutons un facteur $I'_{i + 1}$ à $F_{2i}I$ pour obtenir $F_{2i +
1}I$. Ce sont des sous-modules différentiels gradués et les quotients
successifs sont des produits de décalages de
$A^\vee$. D'après le lemme \ref{lemma-source-graded-injective}, nous
voyons que les inclusions $F_{i + 1}I \to F_iI$ sont des monomorphismes
admissibles. Enfin, nous devons montrer que l'application $M \to I$
(donnée par l'augmentation $M \to I_0$) est un quasi-isomorphisme. Cela découle
d'Homologie, lemme \ref{homology-lemma-good-right-resolution-gives-qis}.
\end{proof}






\section{La catégorie dérivée}
\label{section-derived}

\noindent
Rappelons que les notions de modules différentiels gradués
acycliques et de quasi-isomorphisme de modules différentiels gradués ont un
sens (voir section \ref{section-modules}).

\begin{lemma}
\label{lemma-acyclic}
Soit $(A, \text{d})$ une algèbre différentielle
graduée. La sous-catégorie pleine $\text{Ac}$ de $K(\text{Mod}_{(A,
\text{d})})$ composée de modules acycliques est une sous-catégorie triangulée
strictement pleine et saturée de $K(\text{Mod}_{(A,
\text{d})})$. Le système multiplicatif saturé correspondant
(voir Catégories dérivées, lemme \ref{derived-lemma-operations})
de $K(\text{Mod}_{(A, \text{d})})$ est la classe $\text{Qis}$ des
quasi-isomorphismes. En particulier, le noyau du foncteur de
localisation
$$
Q : K(\text{Mod}_{(A, \text{d})}) \to
\text{Qis}^{-1}K(\text{Mod}_{(A, \text{d})})
$$
est $\text{Ac}$. De plus, le foncteur $H^0$ se factorise par $Q$.
\end{lemma}

\begin{proof}
Nous savons que $H^0$ est un foncteur homologique par la suite
exacte longue d'homologie (\ref{equation-les}). Le
noyau de $H^0$ est la sous-catégorie des objets acycliques et les
flèches qui induisent des isomorphismes sur tous les $H^i$ sont
les quasi-isomorphismes. Ainsi, ce lemme est un cas particulier
de Catégories dérivées, lemme \ref{derived-lemma-acyclic-general}.

\medskip\noindent
Remarque de théorie des ensembles. La construction de la localisation
donnée dans Catégories dérivées,
proposition
\ref{derived-proposition-construct-localization}, exige que la catégorie triangulée donnée soit « petite », c'est-à-dire
que la collection de ses objets forme un ensemble. Soit $V_\alpha$ un
univers partiel (comme dans Ensembles, section \ref{sets-section-sets-hierarchy})
contenant $(A, \text{d})$ et où la cofinalité de $\alpha$ est supérieure à
$\aleph_0$ (voir Ensembles,
proposition \ref{sets-proposition-exist-ordinals-large-cofinality}). Alors nous pouvons
considérer la catégorie $\text{Mod}_{(A, \text{d}), \alpha}$ des $A$-modules
différentiels gradués contenus dans $V_\alpha$. Une vérification simple
montre que toutes les constructions utilisées dans la démonstration de la
proposition \ref{proposition-homotopy-category-triangulated} fonctionnent à
l'intérieur de $\text{Mod}_{(A, \text{d}), \alpha}$ (car au pire
nous prenons des sommes directes finies de modules différentiels gradués). Ainsi, nous
obtenons une catégorie triangulée
$\text{Qis}_\alpha^{-1}K(\text{Mod}_{(A, \text{d}), \alpha})$. Nous verrons
ci-dessous que si $\beta > \alpha$, alors les foncteurs de transition
$$
\text{Qis}_\alpha^{-1}K(\text{Mod}_{(A, \text{d}), \alpha})
\longrightarrow
\text{Qis}_\beta^{-1}K(\text{Mod}_{(A, \text{d}), \beta})
$$
sont pleinement fidèles, puisque les ensembles de morphismes dans les
catégories quotients se calculent à l'aide de morphismes entre P-résolutions dans les
catégories homotopiques (la construction d'une P-résolution dans la démonstration du lemme
\ref{lemma-resolve} prend des sommes directes dénombrables ainsi que des sommes directes
indexées sur des sous-ensembles du module donné). Le lecteur doit donc considérer la
catégorie du lemme comme l'union de ces sous-catégories.
\end{proof}

\noindent
En tenant compte de la remarque de théorie des ensembles à la fin de la
démonstration du lemme précédent, nous définissons la catégorie dérivée comme suit.

\begin{definition}
\label{definition-unbounded-derived-category}
Soit $(A, \text{d})$ une algèbre différentielle graduée.
Soient $\text{Ac}$ et $\text{Qis}$ comme dans le lemme
\ref{lemma-acyclic}. La {\it catégorie dérivée de $(A, \text{d})$} est la catégorie
triangulée
$$
D(A, \text{d}) =
K(\text{Mod}_{(A, \text{d})})/\text{Ac} =
\text{Qis}^{-1}K(\text{Mod}_{(A, \text{d})}).
$$
Nous notons $H^0 : D(A, \text{d}) \to \text{Mod}_R$ le foncteur unique
dont la composition avec le foncteur quotient redonne le foncteur
$H^0$ défini ci-dessus.
\end{definition}

\noindent
Voici le lemme promis calculant les ensembles de morphismes dans
la catégorie dérivée.

\begin{lemma}
\label{lemma-hom-derived}
Soit $(A, \text{d})$ une algèbre différentielle graduée.
Soient $M$ et $N$ des $A$-modules différentiels gradués.
\begin{enumerate}
\item Soit $P \to M$ une P-résolution comme dans
le lemme \ref{lemma-resolve}. Alors
$$
\Hom_{D(A, \text{d})}(M, N) =
\Hom_{K(\text{Mod}_{(A, \text{d})})}(P, N)
$$
\item Soit $N \to I$ une I-résolution comme dans
le lemme \ref{lemma-right-resolution}. Alors
$$
\Hom_{D(A, \text{d})}(M, N) =
\Hom_{K(\text{Mod}_{(A, \text{d})})}(M, I)
$$
\end{enumerate}
\end{lemma}

\begin{proof}
Soit $P \to M$ comme en (1). Puisque $P \to M$ est un quasi-isomorphisme, nous voyons que
$$
\Hom_{D(A, \text{d})}(P, N) = \Hom_{D(A, \text{d})}(M, N)
$$
par définition de la catégorie dérivée. Un
morphisme $f : P \to N$ dans $D(A, \text{d})$ est égal à
$s^{-1}f'$ où $f' : P \to N'$ est un morphisme et où
$s : N \to N'$ est un quasi-isomorphisme. Choisissons un triangle distingué
$$
N \to N' \to Q \to N[1]
$$
Comme $s$ est un quasi-isomorphisme, nous voyons que $Q$ est acyclique.
Ainsi $\Hom_{K(\text{Mod}_{(A, \text{d})})}(P, Q[k]) = 0$ pour tout $k$ par le
lemme \ref{lemma-property-P-K-projective}. Puisque
$\Hom_{K(\text{Mod}_{(A, \text{d})})}(P, -)$ est
cohomologique, nous concluons que nous pouvons relever $f' : P \to N'$ de manière
unique en un morphisme $f : P \to N$. Cela termine la démonstration.

\medskip\noindent
La démonstration de (2) est duale à celle de (1) en
utilisant le lemme \ref{lemma-property-I-K-injective} au
lieu du lemme \ref{lemma-property-P-K-projective}.
\end{proof}

\begin{lemma}
\label{lemma-derived-products}
Soit $(A, \text{d})$ une algèbre différentielle graduée. Alors
\begin{enumerate}
\item $D(A, \text{d})$ a à la fois des sommes directes et des produits,
\item Les sommes directes sont obtenues en prenant des sommes directes de modules
différentiels gradués,
\item Les produits sont obtenus en prenant les produits de modules
différentiels gradués.
\end{enumerate}
\end{lemma}

\begin{proof}
Nous utiliserons que $\text{Mod}_{(A, \text{d})}$ est une catégorie abélienne
avec des sommes directes et des produits arbitraires, et qu'ils donnent
lieu à des sommes directes et des produits dans $K(\text{Mod}_{(A,
\text{d})})$. Voir les lemmes \ref{lemma-dgm-abelian} et \ref{lemma-homotopy-direct-sums}.

\medskip\noindent
Soit $M_j$ une famille de $A$-modules différentiels gradués.
Considérons la somme directe graduée $M = \bigoplus M_j$ qui est un
$A$-module différentiel gradué de manière évidente. Pour un
$A$-module différentiel gradué $N$, choisissons un quasi-isomorphisme $N \to
I$ où $I$ est un $A$-module différentiel gradué avec la propriété (I). Voir le
lemme \ref{lemma-right-resolution}. En utilisant
le lemme \ref{lemma-hom-derived}, nous avons
\begin{align*}
\Hom_{D(A, \text{d})}(M, N)
& =
\Hom_{K(A, \text{d})}(M, I) \\ &
=
\prod \Hom_{K(A, \text{d})}(M_j, I) \\ & =
\prod
\Hom_{D(A, \text{d})}(M_j, N)
\end{align*}
d'où l'existence de sommes directes dans $D(A, \text{d})$ comme indiqué
dans la partie (2) du lemme.

\medskip\noindent
Soit $M_j$ une famille de $A$-modules différentiels gradués.
Considérons le produit $M = \prod M_j$ de $A$-modules différentiels gradués. Pour un
$A$-module différentiel gradué $N$, choisissons un quasi-isomorphisme $P \to
N$ où $P$ est un $A$-module différentiel gradué ayant la propriété (P). Voir le
lemme \ref{lemma-resolve}. En utilisant
le lemme \ref{lemma-hom-derived}, nous avons
\begin{align*}
\Hom_{D(A, \text{d})}(N, M)
& =
\Hom_{K(A, \text{d})}(P, M) \\ &
=
\prod \Hom_{K(A, \text{d})}(P, M_j) \\ & =
\prod
\Hom_{D(A, \text{d})}(N, M_j)
\end{align*}
d'où l'existence de sommes directes dans $D(A, \text{d})$ comme indiqué
dans la partie (3) du lemme.
\end{proof}

\begin{remark}
\label{remark-P-resolution}
Soit $R$ un anneau. Soit $(A, \text{d})$ une $R$-algèbre différentielle graduée. En
utilisant des P-résolutions, nous pouvons parfois réduire les énoncés concernant
des objets généraux de $D(A, \text{d})$ à des énoncés concernant $A[k]$. À
savoir, soit $T$ une propriété d'objets de $D(A, \text{d})$ et supposons que
\begin{enumerate}
\item si $K_i$, $i \in I$ est une famille d'objets de $D(A, \text{d})$
et $T(K_i)$ vaut pour tous les $i \in I$, alors $T(\bigoplus K_i)$,
\item si $K \to L \to M \to K[1]$ est un triangle distingué de
$D(A, \text{d})$ et $T$ vaut pour deux, alors $T$
vaut pour le troisième objet, et
\item $T(A[k])$ vaut pour tous les $k \in \mathbf{Z}$.
\end{enumerate}
Alors $T$ vaut pour tous les objets de $D(A, \text{d})$. Cela découle
clairement des lemmes \ref{lemma-property-P-sequence} et \ref{lemma-resolve}.
\end{remark}







\section{Le foncteur delta canonique}
\label{section-canonical-delta-functor}

\noindent
Soit $(A, \text{d})$ une algèbre différentielle
graduée. Considérons le
foncteur $\text{Mod}_{(A, \text{d})} \to K(\text{Mod}_{(A,
\text{d})})$. Ce foncteur n'est {\bf pas} un $\delta$-foncteur en
général. Cependant, il s'avère que le
foncteur $\text{Mod}_{(A, \text{d})} \to D(A, \text{d})$ est
un $\delta$-foncteur. Pour voir cela, nous devons définir
les morphismes $\delta$ associés à une suite exacte courte
$$
0 \to K \xrightarrow{a} L \xrightarrow{b} M \to 0
$$
dans la catégorie abélienne $\text{Mod}_{(A, \text{d})}$.
Considérons le cône $C(a)$ du morphisme $a$. Nous avons $C(a) = L \oplus K$ et
nous définissons $q : C(a) \to M$ via la projection sur $L$ suivie de $b$.
On obtient donc un homomorphisme de $A$-modules différentiels gradués
$$
q : C(a) \longrightarrow M.
$$
Il est clair que $q \circ i = b$ où $i$ est comme
dans la définition
\ref{definition-cone}. Notons que, comme $a$ est injective, le noyau de $q$ est identifié avec le
cône de $\text{id}_K$ qui est acyclique. Ainsi, nous voyons
que $q$ est un quasi-isomorphisme. Selon le
lemme \ref{lemma-the-same-up-to-isomorphisms}, le
triangle
$$
(K, L, C(a), a, i, -p)
$$
est un triangle distingué dans $K(\text{Mod}_{(A, \text{d})})$. Comme
le foncteur de localisation
$K(\text{Mod}_{(A, \text{d})}) \to D(A, \text{d})$ est exact,
nous voyons que $(K, L, C(a), a, i, -p)$ est un triangle distingué
dans $D(A, \text{d})$. Puisque $q$ est un quasi-isomorphisme, nous
voyons que $q$ est un isomorphisme dans $D(A, \text{d})$. Ainsi,
nous déduisons que
$$
(K, L, M, a, b, -p \circ q^{-1})
$$
est un triangle distingué de $D(A,
\text{d})$. Cela suggère le lemme suivant.

\begin{lemma}
\label{lemma-derived-canonical-delta-functor}
Soit $(A, \text{d})$ une algèbre différentielle graduée. Le foncteur
$\text{Mod}_{(A, \text{d})} \to D(A, \text{d})$
défini possède la structure naturelle d'un $\delta$-foncteur, avec
$$
\delta_{K \to L \to M} = - p \circ q^{-1}
$$
avec $p$ et $q$ comme expliqué ci-dessus.
\end{lemma}

\begin{proof}
Nous avons déjà vu que ce choix conduit à un triangle distingué
chaque fois qu'une suite exacte courte de complexes est donnée.
Nous devons montrer la fonctorialité de cette construction,
voir Catégories dérivées, définition \ref{derived-definition-delta-functor}.
Cela découle du lemme \ref{lemma-functorial-cone} avec un peu de
travail. Comparer avec
Catégories dérivées, lemme \ref{derived-lemma-derived-canonical-delta-functor}.
\end{proof}

\begin{lemma}
\label{lemma-homotopy-colimit}
Soit $(A, \text{d})$ une algèbre différentielle graduée.
Soit $M_n$ un système de modules différentiels gradués. Alors la colimite
dérivée $\text{hocolim} M_n$ dans $D(A, \text{d})$ est représentée
par le module différentiel gradué $\colim M_n$.
\end{lemma}

\begin{proof}
Posons $M = \colim M_n$. Nous avons une suite exacte de modules différentiels gradués
$$
0 \to \bigoplus M_n \to \bigoplus M_n \to M \to 0
$$
par Catégories dérivées, lemme \ref{derived-lemma-compute-colimit}
(appliqué aux complexes sous-jacents de groupes abéliens).
Les sommes directes sont des sommes directes dans
$D(\mathcal{A})$ par le lemme
\ref{lemma-derived-products}. Ainsi, le résultat découle de la
définition des colimites dérivées dans
Catégories dérivées, définition
\ref{derived-definition-derived-colimit} et du fait qu'une suite exacte courte de
complexes donne un triangle distingué
(lemme \ref{lemma-derived-canonical-delta-functor}).
\end{proof}








\section{Catégories linéaires}
\label{section-linear}

\noindent
Juste les définitions.

\begin{definition}
\label{definition-linear-category}
Soit $R$ un anneau. Une {\it catégorie $R$-linéaire $\mathcal{A}$} est une
catégorie où chaque ensemble de morphismes est muni d'une structure de $R$-module
et où, pour $x, y, z \in \Ob(\mathcal{A})$, la loi de composition
$$
\Hom_\mathcal{A}(y, z) \times \Hom_\mathcal{A}(x, y)
\longrightarrow
\Hom_\mathcal{A}(x, z)
$$
est $R$-bilinéaire.
\end{definition}

\noindent
Ainsi la composition détermine une application $R$-linéaire
$$
\Hom_\mathcal{A}(y, z) \otimes_R \Hom_\mathcal{A}(x, y)
\longrightarrow
\Hom_\mathcal{A}(x, z)
$$
de $R$-modules. Notons que nous ne supposons pas que les catégories $R$-linéaires
soient additives.

\begin{definition}
\label{definition-functor-linear-categories}
Soit $R$ un anneau. Un {\it foncteur de catégories $R$-linéaires}, ou
un {\it foncteur $R$-linéaire} est un foncteur $F : \mathcal{A} \to
\mathcal{B}$ où pour tous objets $x, y$ de $\mathcal{A}$
l'application $F : \Hom_\mathcal{A}(x, y) \to \Hom_\mathcal{B}(F(x), F(y))$ est
un homomorphisme de $R$-modules.
\end{definition}







\section{Catégories graduées}
\label{section-graded}

\noindent
Juste quelques définitions.

\begin{definition}
\label{definition-graded-category}
Soit $R$ un anneau. Une {\it catégorie graduée $\mathcal{A}$
sur $R$} est une catégorie où chaque ensemble de morphismes est doté de la
structure d'un $R$-module gradué et où pour
$x, y, z \in \Ob(\mathcal{A})$ la composition est $R$-bilinéaire et induit un
homomorphisme
$$
\Hom_\mathcal{A}(y, z) \otimes_R \Hom_\mathcal{A}(x, y)
\longrightarrow
\Hom_\mathcal{A}(x, z)
$$
de modules gradués sur $R$ (c'est-à-dire préservant les degrés).
\end{definition}

\noindent
Dans cette situation, nous notons $\Hom_\mathcal{A}^i(x, y)$ la composante de
degré $i$ de l'objet gradué $\Hom_\mathcal{A}(x, y)$, de sorte que
$$
\Hom_\mathcal{A}(x, y) =
\bigoplus\nolimits_{i \in \mathbf{Z}} \Hom_\mathcal{A}^i(x, y)
$$
est la décomposition en somme directe en parties graduées.

\begin{definition}
\label{definition-functor-graded-categories}
Soit $R$ un anneau. Un {\it foncteur de catégories graduées sur $R$}, ou un
{\it foncteur gradué}
est un foncteur $F : \mathcal{A} \to \mathcal{B}$ où, pour tous les objets
$x, y$ de $\mathcal{A}$,
l'application $F : \Hom_\mathcal{A}(x, y) \to \Hom_\mathcal{A}(F(x), F(y))$ est
un homomorphisme de $R$-modules gradués.
\end{definition}

\noindent
Étant donnée une catégorie graduée, nous nous intéressons
souvent à la catégorie « usuelle » des applications de degré $0$.
Voici une définition formelle.

\begin{definition}
\label{definition-H0-of-graded-category}
Soit $R$ un anneau. Soit $\mathcal{A}$ une catégorie graduée
sur $R$. On considère {\it $\mathcal{A}^0$} comme la catégorie ayant
les mêmes objets que $\mathcal{A}$ et avec
$$
\Hom_{\mathcal{A}^0}(x, y) = \Hom^0_\mathcal{A}(x, y)
$$
la composante de degré $0$ du module gradué des morphismes de
$\mathcal{A}$.
\end{definition}

\begin{definition}
\label{definition-graded-direct-sum}
Soit $R$ un anneau. Soit $\mathcal{A}$ une catégorie graduée sur $R$. Une
somme directe $(x, y, z, i, j, p, q)$ dans $\mathcal{A}$ (notation comme
dans Homologie, remarque \ref{homology-remark-direct-sum})
est une {\it somme directe graduée} si $i, j, p, q$ sont homogènes
de degré $0$.
\end{definition}

\begin{example}[Catégorie graduée des objets gradués]
\label{example-graded-category-graded-objects}
Soit $\mathcal{B}$ une catégorie additive. Rappelons que nous avons défini
la catégorie $\text{Gr}(\mathcal{B})$ des objets gradués de $\mathcal{B}$ en
Homologie, définition \ref{homology-definition-graded}. Dans cet
exemple, nous allons construire une catégorie graduée
$\text{Gr}^{gr}(\mathcal{B})$ sur $R = \mathbf{Z}$ dont la
catégorie associée $\text{Gr}^{gr}(\mathcal{B})^0$ redonne
$\text{Gr}(\mathcal{B})$. Comme objets de
$\text{Gr}^{gr}(\mathcal{B})$, nous prenons des
objets gradués de $\mathcal{B}$. Alors, étant donnés les objets gradués $A =
(A^i)$ et $B = (B^i)$ de $\mathcal{B}$ nous posons
$$
\Hom_{\text{Gr}^{gr}(\mathcal{B})}(A, B) =
\bigoplus\nolimits_{n \in \mathbf{Z}} \Hom^n(A, B)
$$
où la composante graduée de degré $n$ est le groupe abélien des applications
homogènes de degré $n$ de $A$ vers $B$. Explicitement, nous avons
$$
\Hom^n(A, B) = \prod\nolimits_{p + q = n} \Hom_\mathcal{B}(A^{-q}, B^p)
$$
(observez l'inversion des indices et observez que nous avons un produit ici et non
une somme directe). En d'autres termes, un morphisme de degré $n$
$f$ de $A$ vers $B$ peut être vu comme un système $f = (f_{p, q})$ où
$p, q \in \mathbf{Z}$, $p + q = n$, et où $f_{p, q} :
A^{-q} \to B^p$ est un morphisme de $\mathcal{B}$. Étant donnés
des objets gradués $A$, $B$, $C$ de $\mathcal{B}$, la
composition des morphismes dans $\text{Gr}^{gr}(\mathcal{B})$ est définie via les
applications
$$
\Hom^m(B, C) \times \Hom^n(A, B) \longrightarrow \Hom^{n + m}(A, C)
$$
par simple composition $(g, f) \mapsto g \circ f$ de morphismes
homogènes d'objets gradués. En termes de composantes, nous avons
$$
(g \circ f)_{p, r} = g_{p, q} \circ f_{-q, r}
$$
où $q$ est tel que $p + q = m$ et $-q + r = n$.
\end{example}

\begin{example}[Catégorie graduée des modules gradués]
\label{example-gm-gr-cat}
Soit $A$ une algèbre graduée par $\mathbf{Z}$ sur un anneau $R$. Nous allons
construire une catégorie graduée $\text{Mod}^{gr}_A$ sur $R$ dont la catégorie associée
$(\text{Mod}^{gr}_A)^0$ est la catégorie des $A$-modules gradués. Comme objets de
$\text{Mod}^{gr}_A$, nous prenons des $A$-modules droits gradués (voir section
\ref{section-projectives-over-algebras}). Étant donnés des $A$-modules
gradués $L$ et $M$, nous posons
$$
\Hom_{\text{Mod}^{gr}_A}(L, M) =
\bigoplus\nolimits_{n \in \mathbf{Z}} \Hom^n(L, M)
$$
où $\Hom^n(L, M)$ est l'ensemble des applications de $A$-modules à droite $L \to
M$ qui sont homogènes de degré $n$, c'est-à-dire $f(L^i) \subset M^{i + n}$
pour tout $i \in \mathbf{Z}$. En termes de composantes, nous avons que
$$
\Hom^n(L, M) \subset \prod\nolimits_{p + q = n} \Hom_R(L^{-q}, M^p)
$$
(observer l'inversion des indices) est le sous-ensemble constitué de
ceux $f = (f_{p, q})$ tels que
$$
f_{p, q}(m a) = f_{p - i, q + i}(m)a
$$
pour $a \in A^i$ et $m \in L^{-q - i}$. Pour les $A$-modules gradués $K$,
$L$, $M$, nous définissons la composition dans $\text{Mod}^{gr}_A$ via les
applications
$$
\Hom^m(L, M) \times \Hom^n(K, L) \longrightarrow \Hom^{n + m}(K, M)
$$
par simple composition d'applications de modules à droite sur $A$ : $(g, f) \mapsto g \circ f$.
\end{example}

\begin{remark}
\label{remark-graded-shift-functors}
Soit $R$ un anneau. Soit $\mathcal{D}$ une catégorie $R$-linéaire dotée d'une
collection de foncteurs $R$-linéaires $[n] : \mathcal{D} \to \mathcal{D}$, $x
\mapsto x[n]$ indexés par $n \in \mathbf{Z}$ tels que $[n]
\circ [m] = [n + m]$ et $[0] = \text{id}_\mathcal{D}$ (égalité comme
foncteurs). Cela nous permet de construire une catégorie graduée $\mathcal{D}^{gr}$ sur
$R$ avec les mêmes objets que $\mathcal{D}$ en posant
$$
\Hom_{\mathcal{D}^{gr}}(x, y) =
\bigoplus\nolimits_{n \in \mathbf{Z}} \Hom_\mathcal{D}(x, y[n])
$$
pour $x, y$ dans $\mathcal{D}$. Observez que $(\mathcal{D}^{gr})^0 = \mathcal{D}$
(voir la définition \ref{definition-H0-of-graded-category}). De plus, la catégorie
graduée $\mathcal{D}^{gr}$ hérite de foncteurs $R$-linéaires gradués
$[n]$ satisfaisant $[n] \circ [m] = [n + m]$ et $[0] = \text{id}_{\mathcal{D}^{gr}}$
avec la propriété que
$$
\Hom_{\mathcal{D}^{gr}}(x, y[n]) = \Hom_{\mathcal{D}^{gr}}(x, y)[n]
$$
comme $R$-modules gradués, de manière compatible avec la composition des morphismes.

\medskip\noindent
Inversement, supposons donnée une catégorie graduée $\mathcal{A}$ sur $R$
dotée d'une collection de foncteurs gradués $R$-linéaires
$[n]$ satisfaisant $[n] \circ [m] = [n + m]$ et $[0] =
\text{id}_\mathcal{A}$, qui sont de plus équipés d'isomorphismes
$$
\Hom_\mathcal{A}(x, y[n]) = \Hom_\mathcal{A}(x, y)[n]
$$
comme $R$-modules gradués, de manière compatible avec la composition des
morphismes. Ensuite, le lecteur montre facilement que $\mathcal{A} = (\mathcal{A}^0)^{gr}$.

\medskip\noindent
Voici deux exemples de la relation
$\mathcal{D} \leftrightarrow \mathcal{A}$ que nous avons établie ci-dessus :
\begin{enumerate}
\item Soit $\mathcal{B}$ une catégorie additive. Si
$\mathcal{D} = \text{Gr}(\mathcal{B})$, alors
$\mathcal{A} = \text{Gr}^{gr}(\mathcal{B})$ comme dans
l'exemple \ref{example-graded-category-graded-objects}.
\item Si $A$ est un anneau gradué et $\mathcal{D} = \text{Mod}_A$
est la catégorie des modules gradués à droite sur $A$,
alors $\mathcal{A} = \text{Mod}^{gr}_A$, voir exemple \ref{example-gm-gr-cat}.
\end{enumerate}
\end{remark}






\section{Catégories différentielles graduées}
\label{section-dga-categories}

\noindent
Notons que si $R$ est un anneau, alors $R$ est une algèbre
différentielle graduée sur lui-même (avec $R = R^0$ bien sûr). Dans ce cas, un
$R$-module différentiel gradué est la même chose qu'un complexe de
$R$-modules. En particulier, étant donnés deux $R$-modules différentiels
gradués $M$ et $N$, nous notons $M \otimes_R N$ le $R$-module
différentiel gradué correspondant au complexe total associé au
double complexe obtenu par le produit tensoriel des complexes de $R$-modules
associés à $M$ et $N$.

\begin{definition}
\label{definition-dga-category}
Soit $R$ un anneau. Une {\it catégorie différentielle graduée
$\mathcal{A}$ sur $R$} est une catégorie où chaque ensemble de morphismes est doté de la
structure d'un $R$-module différentiel gradué et où pour
$x, y, z \in \Ob(\mathcal{A})$ la composition est $R$-bilinéaire et induit un
homomorphisme
$$
\Hom_\mathcal{A}(y, z) \otimes_R \Hom_\mathcal{A}(x, y)
\longrightarrow
\Hom_\mathcal{A}(x, z)
$$
de $R$-modules différentiels gradués.
\end{definition}

\noindent
La condition finale de la définition signifie ce qui suit :
si $f \in \Hom_\mathcal{A}^n(x, y)$ et
$g \in \Hom_\mathcal{A}^m(y, z)$ sont homogènes
de degrés $n$ et $m$, alors
$$
\text{d}(g \circ f) = \text{d}(g) \circ f + (-1)^mg \circ \text{d}(f)
$$
dans $\Hom_\mathcal{A}^{n + m + 1}(x, z)$. Cela découle de la règle du
signe pour la différentielle sur le complexe total d'un double complexe,
voir Homologie, définition \ref{homology-definition-associated-simple-complex}.

\begin{definition}
\label{definition-functor-dga-categories}
Soit $R$ un anneau. Un {\it foncteur de catégories différentielles graduées sur $R$}
est un foncteur $F : \mathcal{A} \to \mathcal{B}$ où pour tous les objets
$x, y$ de $\mathcal{A}$,
l'application $F : \Hom_\mathcal{A}(x, y) \to \Hom_\mathcal{A}(F(x), F(y))$
est un homomorphisme de $R$-modules différentiels gradués.
\end{definition}

\noindent
Étant donnée une catégorie différentielle graduée, nous nous intéressons souvent
aux catégories correspondantes de complexes et à la catégorie
homotopique. Voici une définition formelle.

\begin{definition}
\label{definition-homotopy-category-of-dga-category}
Soit $R$ un anneau. Soit $\mathcal{A}$ une catégorie différentielle
graduée sur $R$. On définit
\begin{enumerate}
\item la {\it catégorie des complexes de $\mathcal{A}$}\footnote{Cette
terminologie n'est peut-être pas usuelle.} comme la
catégorie $\text{Comp}(\mathcal{A})$ dont les objets sont les mêmes que les
objets de $\mathcal{A}$ et avec
$$
\Hom_{\text{Comp}(\mathcal{A})}(x, y) =
\Ker(d : \Hom^0_\mathcal{A}(x, y) \to \Hom^1_\mathcal{A}(x, y))
$$
\item la {\it catégorie homotopique de $\mathcal{A}$} comme la
catégorie $K(\mathcal{A})$ dont les objets sont les mêmes que les
objets de $\mathcal{A}$ et avec
$$
\Hom_{K(\mathcal{A})}(x, y) = H^0(\Hom_\mathcal{A}(x, y))
$$
\end{enumerate}
\end{definition}

\noindent
Notre utilisation du symbole $K(\mathcal{A})$ est non standard, mais au moins
elle est compatible avec l'utilisation de $K(-)$ dans d'autres chapitres du projet Champs.

\begin{definition}
\label{definition-dg-direct-sum}
Soit $R$ un anneau. Soit $\mathcal{A}$ une catégorie différentielle graduée sur
$R$. Une somme directe $(x, y, z, i, j, p, q)$ dans $\mathcal{A}$ (notation comme dans
Homologie, remarque \ref{homology-remark-direct-sum})
est une {\it somme directe différentielle graduée} si $i, j, p, q$ sont homogènes
de degré $0$ et fermés, c'est-à-dire $\text{d}(i) = 0$, etc.
\end{definition}

\begin{lemma}
\label{lemma-functorial}
Soit $R$ un anneau. Un foncteur $F : \mathcal{A} \to \mathcal{B}$
de catégories différentielles graduées sur $R$ induit des
foncteurs $\text{Comp}(\mathcal{A}) \to
\text{Comp}(\mathcal{B})$ et $K(\mathcal{A}) \to K(\mathcal{B})$.
\end{lemma}

\begin{proof}
Démonstration omise.
\end{proof}

\begin{example}[Catégorie différentielle graduée des complexes]
\label{example-category-complexes}
Soit $\mathcal{B}$ une catégorie additive. Nous construirons une
catégorie différentielle graduée $\text{Comp}^{dg}(\mathcal{B})$ sur $R
= \mathbf{Z}$ dont la catégorie associée de complexes est
$\text{Comp}(\mathcal{B})$ et dont la catégorie homotopique associée est
$K(\mathcal{B})$. Comme objets de $\text{Comp}^{dg}(\mathcal{B})$, nous prenons des
complexes de $\mathcal{B}$. Étant donnés les complexes
$A^\bullet$ et $B^\bullet$ de $\mathcal{B}$, nous désignons parfois également par
$A^\bullet$ et $B^\bullet$ les objets gradués correspondants de
$\mathcal{B}$ (c'est-à-dire en oubliant la différentielle). En utilisant
cet abus de notation, nous posons
$$
\Hom_{\text{Comp}^{dg}(\mathcal{B})}(A^\bullet, B^\bullet) =
\Hom_{\text{Gr}^{gr}(\mathcal{B})}(A^\bullet, B^\bullet) =
\bigoplus\nolimits_{n \in \mathbf{Z}} \Hom^n(A, B)
$$
comme $\mathbf{Z}$-module gradué, avec les notations et définitions
de l'exemple \ref{example-graded-category-graded-objects}. En
d'autres termes, la composante de degré $n$
est le groupe abélien des morphismes homogènes de degré $n$ d'objets gradués
$$
\Hom^n(A^\bullet, B^\bullet) =
\prod\nolimits_{p + q = n} \Hom_\mathcal{B}(A^{-q}, B^p)
$$
Notons l'inversion des indices et remarquons que nous avons un produit
direct et non une somme directe. Pour un
élément $f \in \Hom^n(A^\bullet, B^\bullet)$ de degré $n$, nous posons
$$
\text{d}(f) = \text{d}_B \circ f - (-1)^n f \circ \text{d}_A
$$
Le signe est exactement
comme dans Plus sur l'algèbre, section
\ref{more-algebra-section-sign-rules}. Pour comprendre cela, nous pensons à $\text{d}_B$ et $\text{d}_A$
comme des morphismes d'objets gradués de $\mathcal{B}$ homogènes de degré $1$ et
nous utilisons la composition dans la catégorie $\text{Gr}^{gr}(\mathcal{B})$
dans le membre de droite. En termes de composantes, si $f = (f_{p, q})$ avec
$f_{p, q} : A^{-q} \to B^p$, nous avons
\begin{equation}
\label{equation-differential-hom-complex}
\text{d}(f_{p, q}) =
\text{d}_B \circ f_{p, q} - (-1)^{p + q} f_{p, q} \circ \text{d}_A
\end{equation}
Notons que le premier terme de cette expression se trouve
dans $\Hom_\mathcal{B}(A^{-q}, B^{p + 1})$ et que le deuxième terme se
trouve dans $\Hom_\mathcal{B}(A^{-q - 1}, B^p)$. Le lecteur vérifie que
\begin{enumerate}
\item $\text{d}$ est de carré nul,
\item un élément $f$ dans $\Hom^n(A^\bullet, B^\bullet)$
a $\text{d}(f) = 0$ si et seulement si le morphisme
$f : A^\bullet \to B^\bullet[n]$ d'objets gradués de $\mathcal{B}$ est en
fait un morphisme de complexes,
\item en particulier, la catégorie des complexes
de $\text{Comp}^{dg}(\mathcal{B})$ est égale à $\text{Comp}(\mathcal{B})$,
\item le morphisme de complexes défini par $f$ comme dans
(2) est homotope à zéro si et seulement si $f = \text{d}(g)$ pour
un certain $g \in \Hom^{n - 1}(A^\bullet, B^\bullet)$.
\item en particulier, nous obtenons un isomorphisme canonique
$$
\Hom_{K(\mathcal{B})}(A^\bullet, B^\bullet)
\longrightarrow
H^0(\Hom_{\text{Comp}^{dg}(\mathcal{B})}(A^\bullet, B^\bullet))
$$
et la catégorie homotopique de $\text{Comp}^{dg}(\mathcal{B})$ est égale à
$K(\mathcal{B})$.
\end{enumerate}
Étant donnés les complexes $A^\bullet$, $B^\bullet$, $C^\bullet$, nous définissons
la composition
$$
\Hom^m(B^\bullet, C^\bullet) \times \Hom^n(A^\bullet, B^\bullet)
\longrightarrow
\Hom^{n + m}(A^\bullet, C^\bullet)
$$
par composition $(g, f) \mapsto g \circ f$ dans la catégorie graduée
$\text{Gr}^{gr}(\mathcal{B})$, voir
exemple \ref{example-graded-category-graded-objects}. Cela
définit un homomorphisme de modules différentiels gradués
$$
\Hom_{\text{Comp}^{dg}(\mathcal{B})}(B^\bullet, C^\bullet)
\otimes_R
\Hom_{\text{Comp}^{dg}(\mathcal{B})}(A^\bullet, B^\bullet)
\longrightarrow
\Hom_{\text{Comp}^{dg}(\mathcal{B})}(A^\bullet, C^\bullet)
$$
comme requis dans la définition \ref{definition-dga-category}
parce que
\begin{align*}
\text{d}(g \circ f) & =
\text{d}_C \circ g \circ f - (-1)^{n + m} g \circ f \circ \text{d}_A \\
& =
\left(\text{d}_C \circ g - (-1)^m g \circ \text{d}_B\right) \circ f +
(-1)^m g \circ \left(\text{d}_B \circ f - (-1)^n f \circ \text{d}_A\right) \\ & =
\text{d}(g)
\circ f + (-1)^m g \circ \text{d}(f)
\end{align*}
comme voulu.
\end{example}

\begin{lemma}
\label{lemma-additive-functor-induces-dga-functor}
Soit $F : \mathcal{B} \to \mathcal{B}'$ un foncteur additif entre des
catégories additives. Alors $F$ induit un foncteur de catégories
différentielles graduées
$$
F : \text{Comp}^{dg}(\mathcal{B}) \to \text{Comp}^{dg}(\mathcal{B}')
$$
de l'exemple \ref{example-category-complexes}
induisant les foncteurs usuels sur la catégorie des complexes et les
catégories d'homotopie.
\end{lemma}

\begin{proof}
Démonstration omise.
\end{proof}

\begin{example}[Catégorie différentielle graduée des modules différentiels gradués]
\label{example-dgm-dg-cat}
Soit $(A, \text{d})$ une algèbre différentielle graduée sur un anneau $R$. Nous allons
construire une catégorie différentielle graduée $\text{Mod}^{dg}_{(A, \text{d})}$ sur
$R$ dont la catégorie des complexes est $\text{Mod}_{(A, \text{d})}$ et dont la
catégorie homotopique est $K(\text{Mod}_{(A, \text{d})})$. Comme objets de
$\text{Mod}^{dg}_{(A, \text{d})}$, nous prenons les
$A$-modules différentiels gradués. Étant donnés des $A$-modules
différentiels gradués $L$ et $M$, nous définissons
$$
\Hom_{\text{Mod}^{dg}_{(A, \text{d})}}(L, M) =
\Hom_{\text{Mod}^{gr}_A}(L, M) = \bigoplus \Hom^n(L, M)
$$
en tant que $R$-module gradué où le membre de droite est défini comme dans
l'exemple \ref{example-gm-gr-cat}. En d'autres termes, la composante de degré $n$
$\Hom^n(L, M)$ est le $R$-module des morphismes de $A$-modules à droite homogènes
de degré $n$. Pour un élément $f \in \Hom^n(L, M)$, nous posons
$$
\text{d}(f) = \text{d}_M \circ f - (-1)^n f \circ \text{d}_L
$$
Pour comprendre cela, nous considérons $\text{d}_M$ et $\text{d}_L$ comme des
morphismes de $R$-modules gradués et nous utilisons la composition de
morphismes de $R$-modules gradués. Il est clair que $\text{d}(f)$ est homogène de degré
$n + 1$ en tant que morphisme de $R$-modules gradués, et il est $A$-linéaire
parce que
\begin{align*}
\text{d}(f)(xa)
& =
\text{d}_M(f(x) a) - (-1)^n f (\text{d}_L(xa)) \\
& =
\text{d}_M(f(x)) a + (-1)^{\deg(x) + n} f(x) \text{d}(a) -
(-1)^n f(\text{d}_L(x)) a - (-1)^{n + \deg(x)} f(x) \text{d}(a) \\ & =
\text{d}(f)(x) a
\end{align*}
comme souhaité (observez que ce calcul ne fonctionnerait pas sans le signe
dans la définition de notre différentielle sur $\Hom$). Des formules similaires
à celles de l'exemple \ref{example-category-complexes} valent pour la
différentielle de $f$ en termes de composantes. Le
lecteur vérifie (de la même manière que dans
l'exemple \ref{example-category-complexes}) que
\begin{enumerate}
\item $\text{d}$ est de carré nul,
\item un élément $f$ dans $\Hom^n(L, M)$ a $\text{d}(f) = 0$ si et seulement si
$f : L \to M[n]$ est un homomorphisme de $A$-modules différentiels gradués,
\item en particulier, la catégorie des complexes
de $\text{Mod}^{dg}_{(A, \text{d})}$ est $\text{Mod}_{(A, \text{d})}$,
\item l'homomorphisme défini par $f$ comme dans (2) est homotope à zéro
si et seulement si $f = \text{d}(g)$ pour un certain
$g \in \Hom^{n - 1}(L, M)$.
\item en particulier, nous obtenons un isomorphisme canonique
$$
\Hom_{K(\text{Mod}_{(A, \text{d})})}(L, M)
\longrightarrow
H^0(\Hom_{\text{Mod}^{dg}_{(A, \text{d})}}(L, M))
$$
et la catégorie homotopique de $\text{Mod}^{dg}_{(A, \text{d})}$
est $K(\text{Mod}_{(A, \text{d})})$.
\end{enumerate}
Étant donnés des $A$-modules différentiels gradués $K$, $L$, $M$, nous définissons
la composition
$$
\Hom^m(L, M) \times \Hom^n(K, L) \longrightarrow \Hom^{n + m}(K, M)
$$
par composition de morphismes homogènes de $A$-modules à droite $(g, f) \mapsto g \circ f$.
Cela définit un homomorphisme de modules différentiels gradués
$$
\Hom_{\text{Mod}^{dg}_{(A, \text{d})}}(L, M) \otimes_R
\Hom_{\text{Mod}^{dg}_{(A, \text{d})}}(K, L) \longrightarrow
\Hom_{\text{Mod}^{dg}_{(A, \text{d})}}(K, M)
$$
comme requis
dans la définition \ref{definition-dga-category}
parce que
\begin{align*}
\text{d}(g \circ f) & =
\text{d}_M \circ g \circ f - (-1)^{n + m} g \circ f \circ \text{d}_K \\
& =
\left(\text{d}_M \circ g - (-1)^m g \circ \text{d}_L\right) \circ f +
(-1)^m g \circ \left(\text{d}_L \circ f - (-1)^n f \circ \text{d}_K\right) \\ & =
\text{d}(g)
\circ f + (-1)^m g \circ \text{d}(f)
\end{align*}
comme voulu.
\end{example}

\begin{lemma}
\label{lemma-homomorphism-induces-dga-functor}
Soit $\varphi : (A, \text{d}) \to (E, \text{d})$ un homomorphisme
d'algèbres différentielles graduées. Alors $\varphi$ induit un foncteur de catégories
différentielles graduées
$$
F :
\text{Mod}^{dg}_{(E, \text{d})}
\longrightarrow
\text{Mod}^{dg}_{(A, \text{d})}
$$
de l'exemple \ref{example-dgm-dg-cat} induisant des foncteurs de restriction
évidents sur les catégories de modules différentiels gradués et les catégories d'homotopie.
\end{lemma}

\begin{proof}
Démonstration omise.
\end{proof}

\begin{lemma}
\label{lemma-construction}
Soit $R$ un anneau. Soit $\mathcal{A}$ une catégorie différentielle
graduée sur $R$. Soit $x$ un objet de $\mathcal{A}$. Soit
$$
(E, \text{d}) = \Hom_\mathcal{A}(x, x)
$$
l'algèbre différentielle graduée sur $R$ des endomorphismes de $x$.
Nous obtenons un foncteur
$$
\mathcal{A} \longrightarrow \text{Mod}^{dg}_{(E, \text{d})},\quad
y \longmapsto \Hom_\mathcal{A}(x, y)
$$
de catégories différentielles graduées en laissant $E$ agir
sur $\Hom_\mathcal{A}(x, y)$ via la composition dans $\mathcal{A}$.
Ce foncteur induit des foncteurs
$$
\text{Comp}(\mathcal{A}) \to \text{Mod}_{(A, \text{d})}
\quad\text{et}\quad
K(\mathcal{A}) \to K(\text{Mod}_{(A, \text{d})})
$$
par une application du lemme \ref{lemma-functorial}.
\end{lemma}

\begin{proof}
La démonstration est immédiate.
\end{proof}









\section{Construction de catégories triangulées}
\label{section-review}

\noindent
Dans cette section, nous examinons le cadre le plus général auquel s'appliquent
les démonstrations de Catégories dérivées,
proposition
\ref{derived-proposition-homotopy-category-triangulated} et de la proposition \ref{proposition-homotopy-category-triangulated}.

\medskip\noindent
Soit $R$ un anneau. Soit $\mathcal{A}$ une catégorie différentielle graduée sur
$R$. Pour que notre argument fonctionne, nous imposons quelques axiomes sur $\mathcal{A}$ :
\begin{enumerate}
\item[(A)] $\mathcal{A}$ possède un objet nul et des sommes
directes différentielles graduées de deux
objets (comme dans la définition
\ref{definition-dg-direct-sum}). \item[(B)] il existe des foncteurs $[n] : \mathcal{A} \longrightarrow
\mathcal{A}$ de catégories différentielles graduées tels que
$[0] = \text{id}_\mathcal{A}$ et $[n + m] = [n] \circ [m]$ et des
isomorphismes donnés
$$
\Hom_\mathcal{A}(x, y[n]) = \Hom_\mathcal{A}(x, y)[n]
$$
de $R$-modules différentiels gradués compatibles avec la composition.
\end{enumerate}

\noindent
Étant donnée notre catégorie différentielle graduée $\mathcal{A}$, nous disons que
\begin{enumerate}
\item une suite $x \to y \to z$ de morphismes de $\text{Comp}(\mathcal{A})$
est une {\it suite exacte courte admissible} s'il existe un
isomorphisme $y \cong x \oplus z$ dans la catégorie graduée sous-jacente tel
que $x \to z$ et $y \to z$ soient des (co)projections.
\item un morphisme $x\to y$ de $\text{Comp}(\mathcal{A})$ est un
{\it monomorphisme admissible} s'il se prolonge en
une suite exacte courte admissible $x\to y\to z$.
\item un morphisme $y\to z$ de $\text{Comp}(\mathcal{A})$ est un
{\it épimorphisme admissible} s'il se prolonge en
une suite exacte courte admissible $x\to y\to z$.
\end{enumerate}
Le lemme suivant nous indique qu'une suite exacte courte admissible donne un
triangle, à condition que nous ayons les axiomes (A) et (B).

\begin{lemma}
\label{lemma-get-triangle}
Soit $\mathcal{A}$ une catégorie différentielle graduée satisfaisant aux
axiomes (A) et (B). Étant donnée une suite exacte courte admissible $x \to
y \to z$, on obtient (voir la démonstration) un triangle
$$
x \to y \to z \to x[1]
$$
dans $\text{Comp}(\mathcal{A})$ avec la propriété que toute composée de deux flèches
consécutives de $z[-1] \to x \to y \to z \to x[1]$ est nulle dans $K(\mathcal{A})$.
\end{lemma}

\begin{proof}
Choisissons un diagramme
$$
\xymatrix{
x \ar[rr]_1 \ar[rd]_a & & x \\
& y \ar[ru]_\pi \ar[rd]^b & \\ z
\ar[rr]^1 \ar[ru]^s & & z
}
$$
donnant l'isomorphisme des objets gradués $y \cong x \oplus z$ comme dans la
définition d'une suite exacte courte admissible. Voici quelques équations qui sont
valables dans cette situation
\begin{enumerate}
\item $1 = \pi a$ et donc $\text{d}(\pi) a = 0$,
\item $1 = b s$ et donc $b \text{d}(s) = 0$,
\item $1 = a \pi + s b$ et donc $a \text{d}(\pi) + \text{d}(s) b = 0$,
\item $\pi s = 0$ et donc $\text{d}(\pi)s + \pi \text{d}(s) = 0$,
\item $\text{d}(s) = a \pi \text{d}(s)$ parce
que $\text{d}(s) = (a \pi + s b)\text{d}(s)$ et $b\text{d}(s) = 0$,
\item $\text{d}(\pi) = \text{d}(\pi) s b$ parce
que $\text{d}(\pi) = \text{d}(\pi)(a \pi + s b)$ et $\text{d}(\pi)a = 0$,
\item $\text{d}(\pi \text{d}(s)) = 0$ parce que si nous le
postcomposons avec le monomorphisme $a$ nous
obtenons $\text{d}(a\pi \text{d}(s)) = \text{d}(\text{d}(s)) = 0$, et
\item $\text{d}(\text{d}(\pi)s) = 0$ car, d'après (4), c'est
l'opposé de $\text{d}(\pi\text{d}(s))$ qui est $0$ par (7).
\end{enumerate}
Nous avons utilisé à plusieurs reprises que $\text{d}(a) = 0$, $\text{d}(b)
= 0$, et que $\text{d}(1) = 0$. Par (7) nous voyons que
$$
\delta = \pi \text{d}(s) = - \text{d}(\pi) s : z \to x[1]
$$
est un morphisme dans $\text{Comp}(\mathcal{A})$. D'après (5), nous
voyons que la composition $a \delta = a \pi \text{d}(s) =
\text{d}(s)$ est homotope à zéro. D'après (6), nous voyons que la
composition $\delta b = - \text{d}(\pi)sb = \text{d}(-\pi)$ est homotope à zéro.
\end{proof}

\noindent
Outre les axiomes (A) et (B), nous avons besoin d'un axiome concernant
l'existence des cônes. Nous formalisons tout comme suit.

\begin{situation}
\label{situation-ABC}
Ici $R$ est un anneau et $\mathcal{A}$ est une catégorie différentielle graduée
sur $R$ satisfaisant aux axiomes (A), (B), et
\begin{enumerate}
\item[(C)] étant donnée une flèche $f : x \to y$ de degré $0$
avec $\text{d}(f) = 0$, il existe une suite exacte courte admissible
$y \to c(f) \to x[1]$ dans $\text{Comp}(\mathcal{A})$ telle que
l'application $x[1] \to y[1]$ du lemme \ref{lemma-get-triangle} soit égale à $f[1]$.
\end{enumerate}
\end{situation}

\noindent
Nous appellerons $c(f)$ un {\it cône} du morphisme $f$. Si (A),
(B) et (C) sont satisfaits, alors les
cônes sont fonctoriels dans un sens faible.

\begin{lemma}
\label{lemma-cone}
\begin{slogan}
La catégorie homotopique est une catégorie
triangulée. Ce lemme prouve une partie des axiomes d'une catégorie triangulée.
\end{slogan}
Dans la situation \ref{situation-ABC}, supposons que
$$
\xymatrix{
x_1 \ar[r]_{f_1} \ar[d]_a & y_1 \ar[d]^b \\
x_2 \ar[r]^{f_2} & y_2
}
$$
soit un diagramme de $\text{Comp}(\mathcal{A})$ commutatif à homotopie
près. Alors il existe un morphisme $c : c(f_1) \to c(f_2)$ qui donne lieu à un
morphisme de triangles
$$
(a, b, c) : (x_1, y_1, c(f_1)) \to (x_1, y_1, c(f_1))
$$
dans $K(\mathcal{A})$.
\end{lemma}

\begin{proof}
L'hypothèse signifie qu'il existe un morphisme $h : x_1 \to y_2$ de degré
$-1$ tel que $\text{d}(h) = b f_1 - f_2 a$. Choisissons des
isomorphismes $c(f_i) = y_i \oplus x_i[1]$ d'objets gradués compatibles avec les
morphismes $y_i \to c(f_i) \to x_i[1]$. Désignons par
$a_i : y_i \to c(f_i)$, $b_i : c(f_i) \to x_i[1]$, $s_i : x_i[1] \to c(f_i)$ et
$\pi_i : c(f_i) \to y_i$ les morphismes donnés. Rappelons que $x_i[1] \to
y_i[1]$ est donné par $\pi_i \text{d}(s_i)$. Selon l'axiome (C), cela
signifie que
$$
f_i = \pi_i \text{d}(s_i) = - \text{d}(\pi_i) s_i
$$
(nous identifions $\Hom(x_i, y_i)$ avec $\Hom(x_i[1], y_i[1])$ en
utilisant le foncteur de décalage $[1]$).
Posons $c = a_2 b \pi_1 + s_2 a b_1 + a_2hb$. Puis, en utilisant les
égalités trouvées dans la démonstration du lemme \ref{lemma-get-triangle}, nous
obtenons
\begin{align*}
\text{d}(c)
& =
a_2 b \text{d}(\pi_1) + \text{d}(s_2) a b_1 + a_2 \text{d}(h) b_1 \\ &
= -
a_2 b f_1 b_1 + a_2 f_2 a b_1 + a_2 (b f_1 - f_2 a) b_1 \\ & =
0
\end{align*}
(où nous avons utilisé en particulier
que $\text{d}(\pi_1) = \text{d}(\pi_1) s_1 b_1 = f_1 b_1$ et
$\text{d}(s_2) = a_2 \pi_2 \text{d}(s_2) = a_2 f_2$).
Ainsi $c$ est un morphisme de degré $0$ $c : c(f_1) \to c(f_2)$ de $\mathcal{A}$
compatible avec les morphismes donnés $y_i \to c(f_i) \to x_i[1]$.
\end{proof}

\noindent
Dans la situation \ref{situation-ABC}, nous disons qu'un
triangle $(x, y, z, f, g, h)$ dans
$K(\mathcal{A})$ est un {\it triangle distingué} s'il existe une suite
exacte courte admissible $x' \to y' \to z'$ telle que $(x,
y, z, f, g, h)$ soit isomorphe comme triangle dans $K(\mathcal{A})$ au
triangle $(x', y', z', x' \to y', y' \to z', \delta)$ construit
dans le lemme \ref{lemma-get-triangle}. Nous montrerons ci-dessous que
$$
\boxed{
K(\mathcal{A})\text{ est une catégorie triangulée}
}
$$
Ce résultat, bien qu'il ne soit pas aussi général qu'on pourrait le penser,
s'applique à un certain nombre de généralisations naturelles des cas couverts jusqu'à présent
dans le projet Champs. Voici quelques exemples :
\begin{enumerate}
\item Soit $(X, \mathcal{O}_X)$ un espace annelé. Soit $(A, d)$ un
faisceau d'algèbres différentielles graduées sur $\mathcal{O}_X$. Soit
$\mathcal{A}$ la catégorie différentielle graduée des modules
différentiels gradués sur $A$. Alors $K(\mathcal{A})$ est une catégorie triangulée.
\item Soit $(\mathcal{C}, \mathcal{O})$ un site annelé. Soit $(A, d)$ un
faisceau d'algèbres différentielles graduées sur $\mathcal{O}$.
Soit $\mathcal{A}$ la catégorie différentielle graduée des $A$-modules
différentiels gradués. Alors $K(\mathcal{A})$ est une catégorie triangulée.
Voir Faisceaux Différentiels Gradués, proposition
\ref{sdga-proposition-homotopy-category-triangulated}.
\item On trouvera deux exemples d'une autre nature dans Exemples, section
\ref{examples-section-nongraded-differential-graded}.
\end{enumerate}

\noindent
Le lemme élémentaire suivant joue un rôle essentiel dans la construction.

\begin{lemma}
\label{lemma-id-cone-null}
Dans la situation
\ref{situation-ABC}, étant donné tout objet $x$ de $\mathcal{A}$, et le cône $C(1_x)$
du morphisme identité $1_x : x \to x$, le morphisme identité sur
$C(1_x)$ est homotope à zéro.
\end{lemma}

\begin{proof}
Considérons la suite exacte courte admissible donnée par l'axiome (C).
$$
\xymatrix{
x \ar@<0.5ex>[r]^a &
C(1_x) \ar@<0.5ex>[l]^{\pi} \ar@<0.5ex>[r]^b &
x[1]\ar@<0.5ex>[l]^s
}
$$
Alors, d'après le lemme \ref{lemma-get-triangle}, en identifiant les ensembles
de morphismes par décalage, nous avons $1_x=\pi d(s)=-d(\pi)s$ où $s$ est
considéré comme un morphisme dans $\Hom_{\mathcal{A}}^{-1}(x,C(1_x))$. Par
conséquent, $a=a\pi d(s)=d(s)$ en utilisant la formule (5) du lemme
\ref{lemma-get-triangle}, et $b=-d(\pi)sb=-d(\pi)$ par la formule (6) du lemme \ref{lemma-get-triangle}.
Ainsi
$$
1_{C(1_x)} = a\pi + sb = d(s)\pi - sd(\pi) = d(s\pi)
$$
puisque $s$ est de degré $-1$.
\end{proof}

\noindent
Une version plus générale du lemme ci-dessus apparaîtra
dans le lemme \ref{lemma-cone-homotopy}. Le lemme suivant est
l'analogue du lemme \ref{lemma-make-commute-map}.

\begin{lemma}
\label{lemma-homo-change}
Dans la situation \ref{situation-ABC}, soit un diagramme
$$
\xymatrix{x\ar[r]^f\ar[d]_a & y\ar[d]^b\\
z\ar[r]^g & w}
$$
dans $\text{Comp}(\mathcal{A})$ commutatif à homotopie près. Alors
\begin{enumerate}
\item Si $f$ est un monomorphisme admissible, alors $b$ est
homotope à un morphisme $b'$ qui fait commuter le diagramme.
\item Si $g$ est un épimorphisme admissible, alors $a$ est
homotope à un morphisme $a'$ qui fait commuter le diagramme.
\end{enumerate}
\end{lemma}

\begin{proof}
Pour prouver (1), on observe que l'hypothèse implique qu'il existe un
$h\in\Hom_{\mathcal{A}}(x,w)$ de degré $-1$ tel que $bf-ga=d(h)$.
Puisque $f$ est un monomorphisme admissible, il existe un morphisme
$\pi : y \to x$ dans la catégorie $\mathcal{A}$ de degré $0$.
Soit $b' = b - d(h\pi)$. Alors
\begin{align*}
b'f = bf - d(h\pi)f
= &
bf - d(h\pi f) \quad (\text{puisque }d(f) = 0) \\ = &
bf-d(h)
\\ =
&
ga
\end{align*}
comme voulu. La démonstration pour (2) est omise.
\end{proof}

\noindent
Le lemme suivant est l'analogue du lemme \ref{lemma-make-injective}.

\begin{lemma}
\label{lemma-factor}
Dans la situation \ref{situation-ABC}, soit $\alpha : x \to
y$ un morphisme dans $\text{Comp}(\mathcal{A})$. Alors il existe
une factorisation dans $\text{Comp}(\mathcal{A})$ :
$$
\xymatrix{
x \ar[r]^{\tilde{\alpha}} &
\tilde{y} \ar@<0.5ex>[r]^{\pi} &
y\ar@<0.5ex>[l]^s
}
$$
avec les propriétés suivantes :
\begin{enumerate}
\item $\tilde{\alpha}$ est un monomorphisme admissible, et
$\pi\tilde{\alpha}=\alpha$.
\item Il existe un morphisme
$s:y\to\tilde{y}$ dans $\text{Comp}(\mathcal{A})$
tel que $\pi s=1_y$ et que $s\pi$ soit homotope à $1_{\tilde{y}}$.
\end{enumerate}
\end{lemma}

\begin{proof}
D'après l'axiome (A), nous pouvons prendre pour $\tilde{y}$ la somme directe
différentielle graduée de $y$ et $C(1_x)$ ; il existe donc un diagramme
$$
\xymatrix@C=3pc{
y \ar@<0.5ex>[r]^s &
y\oplus C(1_x) \ar@<0.5ex>[l]^{\pi} \ar@<0.5ex>[r]^{p} &
C(1_x)\ar@<0.5ex>[l]^t
}
$$
où tous les morphismes sont de degré zéro, et dans
$\text{Comp}(\mathcal{A})$. Soit $\tilde{y} = y \oplus C(1_x)$. Alors
$1_{\tilde{y}} = s\pi + tp$. Considérons maintenant le diagramme
$$
\xymatrix{
x \ar[r]^{\tilde{\alpha}} &
\tilde{y} \ar@<0.5ex>[r]^{\pi} &
y\ar@<0.5ex>[l]^s
}
$$
où $\tilde{\alpha}$ est induit par le morphisme $x\xrightarrow{\alpha}y$
et le morphisme naturel $x\to C(1_x)$ s'insérant dans la suite
exacte courte admissible
$$
\xymatrix{
x \ar@<0.5ex>[r] &
C(1_x) \ar@<0.5ex>[l] \ar@<0.5ex>[r] &
x[1]\ar@<0.5ex>[l]
}
$$
Ainsi, le morphisme $C(1_x)\to x$ de degré 0 dans ce
diagramme, ainsi que le morphisme nul $y\to x$, induisent un morphisme de
degré 0 $\beta : \tilde{y} \to x$. Ensuite, $\tilde{\alpha}$ est un
monomorphisme admissible puisqu'il s'intègre dans la suite exacte
courte admissible
$$
\xymatrix{
x\ar[r]^{\tilde{\alpha}} &
\tilde{y} \ar[r] &
x[1]
}
$$
De plus, $\pi\tilde{\alpha} = \alpha$ par la construction de
$\tilde{\alpha}$, et $\pi s = 1_y$ par le premier diagramme. Il reste à montrer
que $s\pi$ est homotope à $1_{\tilde{y}}$. Écrivons $1_x$ sous la forme
$d(h)$ pour un certain morphisme de degré $-1$. Ensuite, notre
dernière affirmation découle de
\begin{align*}
1_{\tilde{y}} - s\pi =
& tp
\\ = &
t(dh)p\quad\text{(d'après
le lemme \ref{lemma-id-cone-null})} \\ =
&
d(thp)
\end{align*}
puisque $dt = dp = 0$, et $t$ est de degré zéro.
\end{proof}

\noindent
Le lemme suivant est l'analogue du lemme \ref{lemma-sequence-maps-split}.

\begin{lemma}
\label{lemma-analogue-sequence-maps-split}
Dans la situation
\ref{situation-ABC}, soit $x_1 \to x_2 \to \ldots \to
x_n$ une suite de morphismes composables dans $\text{Comp}(\mathcal{A})$.
Alors il existe un diagramme commutatif dans $\text{Comp}(\mathcal{A})$ :
$$
\xymatrix{x_1\ar[r] & x_2\ar[r] & \ldots\ar[r] & x_n\\
y_1\ar[r]\ar[u] & y_2\ar[r]\ar[u] & \ldots\ar[r] & y_n\ar[u]}
$$
de telle sorte que chaque $y_i\to y_{i+1}$ soit un monomorphisme
admissible et chaque $y_i\to x_i$ soit une équivalence d'homotopie.
\end{lemma}

\begin{proof}
Le cas de $n=1$ est trivial : on prend simplement $y_1 = x_1$ et le
morphisme identité sur $x_1$ est en particulier une équivalence d'homotopie. Le
cas $n = 2$ est donné par le lemme \ref{lemma-factor}. Supposons que nous
ayons construit le diagramme jusqu'à $x_{n - 1}$. Nous
appliquons le lemme \ref{lemma-factor} à la composition
$y_{n - 1} \to x_{n-1} \to x_n$ pour obtenir $y_n$.
Ensuite, $y_{n - 1} \to y_n$ sera un monomorphisme admissible, et $y_n
\to x_n$ une équivalence d'homotopie.
\end{proof}

\noindent
Le lemme suivant est l'analogue du lemme \ref{lemma-nilpotent}.

\begin{lemma}
\label{lemma-triseq}
Dans la situation \ref{situation-ABC}, soient $x_i \to y_i \to
z_i$ des morphismes dans $\mathcal{A}$ ($i=1,2,3$) tels
que $x_2 \to y_2\to z_2$ soit une suite exacte courte admissible.
Soient $b : y_1 \to y_2$ et $b' : y_2\to y_3$ des morphismes
dans $\text{Comp}(\mathcal{A})$ tels que
$$
\vcenter{
\xymatrix{
x_1 \ar[d]_0 \ar[r] &
y_1 \ar[r] \ar[d]_b &
z_1 \ar[d]_0 \\ x_2
\ar[r] & y_2 \ar[r] & z_2 } }
\quad\text{et}\quad
\vcenter{
\xymatrix{
x_2
\ar[d]^0
\ar[r] & y_2 \ar[r]
\ar[d]^{b'} & z_2 \ar[d]^0
\\ x_3 \ar[r] &
y_3 \ar[r] & z_3
}
}
$$
commutent à homotopie près. Alors $b'\circ b$ est homotope à $0$.
\end{lemma}

\begin{proof}
Par le lemme \ref{lemma-homo-change}, nous pouvons remplacer $b$ et
$b'$ par des morphismes homotopes $\tilde{b}$ et $\tilde{b}'$, de telle sorte que
le carré de droite du diagramme de gauche commute et que le carré de gauche du
diagramme de droite commute. Notons $b = \tilde{b} + d(h)$ et $b'=\tilde{b}'+d(h')$
pour des morphismes de degré $-1$ $h$ et $h'$ dans $\mathcal{A}$. D'où
$$
b'b = \tilde{b}'\tilde{b} + d(\tilde{b}'h + h'\tilde{b} + h'd(h))
$$
puisque $d(\tilde{b})=d(\tilde{b}')=0$, c'est-à-dire que $b'b$ est
homotope à $\tilde{b}'\tilde{b}$. Nous voulons maintenant montrer que
$\tilde{b}'\tilde{b}=0$. Comme $x_2\xrightarrow{f} y_2\xrightarrow{g} z_2$ est une suite
exacte courte admissible, il existe des morphismes de degré
$0$ $\pi : y_2 \to x_2$ et $s : z_2 \to y_2$ tels que
$\text{id}_{y_2} = f\pi + sg$. Par conséquent
$$
\tilde{b}'\tilde{b} = \tilde{b}'(f\pi + sg)\tilde{b} = 0
$$
puisque $g\tilde{b} = 0$ et $\tilde{b}'f = 0$ d'après
les deux carrés commutatifs.
\end{proof}

\noindent
Le lemme suivant est l'analogue du
lemme \ref{lemma-triangle-independent-splittings}.

\begin{lemma}
\label{lemma-analogue-triangle-independent-splittings}
Dans la situation
\ref{situation-ABC}, soit $0 \to x \to y \to z \to 0$ une suite exacte
courte admissible dans $\text{Comp}(\mathcal{A})$. Le triangle
$$
\xymatrix{x\ar[r] & y\ar[r] & z\ar[r]^{\delta} & x[1]}
$$
avec $\delta : z \to x[1]$ tel que défini dans le lemme
\ref{lemma-get-triangle} est, à isomorphisme canonique près dans $K(\mathcal{A})$, indépendant des
choix faits dans le lemme \ref{lemma-get-triangle}.
\end{lemma}

\begin{proof}
Supposons que $\delta$ soit défini par le scindage
$$
\xymatrix{
x \ar@<0.5ex>[r]^{a} &
y \ar@<0.5ex>[r]^b\ar@<0.5ex>[l]^{\pi} & z
\ar@<0.5ex>[l]^s
}
$$
et que $\delta'$ soit défini par le scindage avec $\pi',s'$ à
la place de $\pi,s$. Ensuite
$$
s'-s = (a\pi + sb)(s'-s) = a\pi s'
$$
puisque $bs' = bs = 1_z$ et $\pi s = 0$. De même,
$$
\pi' - \pi = (\pi' - \pi)(a\pi + sb) = \pi'sb
$$
Puisque $\delta = \pi d(s)$ et $\delta' = \pi'd(s')$
d'après la construction du lemme \ref{lemma-get-triangle}, nous pouvons calculer
$$
\delta' = \pi'd(s') = (\pi + \pi'sb)d(s + a\pi s') = \delta + d(\pi s')
$$
en utilisant $\pi a = 1_x$, $ba = 0$ et $\pi'sbd(s') = \pi'sba\pi d(s') = 0$
selon la formule (5) dans le lemme \ref{lemma-get-triangle}.
\end{proof}

\noindent
Le lemme suivant est l'analogue du lemme \ref{lemma-rotate-cone}.

\begin{lemma}
\label{lemma-restate-axiom-c}
Dans la situation
\ref{situation-ABC}, soit $f: x \to y$ un morphisme dans
$\text{Comp}(\mathcal{A})$. Le triangle $(y, c(f), x[1], i, p, f[1])$ est le triangle
associé à la suite exacte courte admissible
$$
\xymatrix{y\ar[r] & c(f) \ar[r] & x[1]}
$$
où le cône $c(f)$ est défini comme dans le lemme \ref{lemma-get-triangle}.
\end{lemma}

\begin{proof}
Cela découle de l'axiome (C).
\end{proof}

\noindent
Le lemme suivant est l'analogue du lemme \ref{lemma-rotate-triangle}.

\begin{lemma}
\label{lemma-cone-rotate-isom}
Dans la situation \ref{situation-ABC}, supposons que $\alpha : x \to y$ et $\beta : y \to
z$ définissent une suite exacte courte admissible
$$
\xymatrix{
x \ar[r] &
y\ar[r] &
z
}
$$
dans $\text{Comp}(\mathcal{A})$. Soit $(x, y, z, \alpha, \beta,
\delta)$ le triangle associé dans $K(\mathcal{A})$. Alors les triangles
$$
(z[-1], x, y, \delta[-1], \alpha, \beta)
\quad\text{et}\quad
(z[-1], x, c(\delta[-1]), \delta[-1], i, p)
$$
sont isomorphes.
\end{lemma}

\begin{proof}
Nous avons un diagramme de la forme
$$
\xymatrix{
z[-1]\ar[r]^{\delta[-1]}\ar[d]^1 &
x\ar@<0.5ex>[r]^{\alpha}\ar[d]^1 &
y\ar@<0.5ex>[r]^{\beta}\ar@{.>}[d]\ar@<0.5ex>[l]^{\tilde{\alpha}} &
z\ar[d]^1\ar@<0.5ex>[l]^{\tilde\beta} \\
z[-1] \ar[r]^{\delta[-1]} &
x\ar@<0.5ex>[r]^i &
c(\delta[-1]) \ar@<0.5ex>[r]^p\ar@<0.5ex>[l]^{\tilde i} &
z\ar@<0.5ex>[l]^{\tilde p}
}
$$
avec les scindages de $\alpha, \beta, i$, et $p$ données
par $\tilde{\alpha}, \tilde{\beta}, \tilde{i},$ et $\tilde{p}$ respectivement.
Définissons un morphisme $y \to c(\delta[-1])$ par
$i\tilde{\alpha} + \tilde{p}\beta$ et un morphisme
$c(\delta[-1]) \to y$ par $\alpha \tilde{i} + \tilde{\beta} p$.
Vérifions d'abord que ceux-ci définissent des morphismes dans $\text{Comp}(\mathcal{A})$.
Nous remarquons que, d'après les identités du lemme
\ref{lemma-get-triangle}, nous avons la
relation $\delta[-1] = \tilde{\alpha}d(\tilde{\beta}) = -d(\tilde{\alpha})\tilde{\beta}$ et
la relation $\delta[-1] = \tilde{i}d(\tilde{p})$. Ensuite
\begin{align*}
d(\tilde{\alpha})
& =
d(\tilde{\alpha})\tilde{\beta}\beta \\ &
=
-\delta[-1]\beta
\end{align*}
où nous avons utilisé l'équation (6)
du lemme \ref{lemma-get-triangle} pour la première égalité et
la remarque précédente pour la seconde. De même, nous obtenons
$d(\tilde{p}) = i\delta[-1]$. Ainsi
\begin{align*}
d(i\tilde{\alpha} + \tilde{p}\beta)
& =
d(i)\tilde{\alpha} + id(\tilde{\alpha}) +
d(\tilde{p})\beta + \tilde{p}d(\beta) \\ & =
id(\tilde{\alpha})
+ d(\tilde{p})\beta \\ & =
-i\delta[-1]\beta
+ i\delta[-1]\beta \\ &
=
0
\end{align*}
donc $i\tilde{\alpha} + \tilde{p}\beta$ est en effet un morphisme
de $\text{Comp}(\mathcal{A})$. Par un calcul similaire,
$\alpha \tilde{i} + \tilde{\beta} p$ est également un morphisme de
$\text{Comp}(\mathcal{A})$. Il est immédiat que ces morphismes s'insèrent
dans le diagramme commutatif. Nous calculons :
\begin{align*}
(i\tilde{\alpha} + \tilde{p}\beta)(\alpha \tilde{i} + \tilde{\beta} p)
& =
i\tilde{\alpha}\alpha\tilde{i} + i\tilde{\alpha}\tilde{\beta}p +
\tilde{p}\beta\alpha\tilde{i} + \tilde{p}\beta\tilde{\beta}p \\ & =
i\tilde{i}
+ \tilde{p}p \\ &
=
1_{c(\delta[-1])}
\end{align*}
où nous avons librement utilisé les identités
du lemme \ref{lemma-get-triangle}. De même, nous calculons
$(\alpha \tilde{i} + \tilde{\beta} p)(i\tilde{\alpha} + \tilde{p}\beta) = 1_y$, donc nous
concluons $y \cong c(\delta[-1])$. Par conséquent, les deux triangles en question
sont isomorphes.
\end{proof}

\noindent
Le lemme suivant est l'analogue du
lemme \ref{lemma-third-isomorphism}.

\begin{lemma}
\label{lemma-analogue-third-isomorphism}
Dans la situation \ref{situation-ABC}, soient $f_1 : x_1 \to y_1$
et $f_2 : x_2 \to y_2$ des morphismes dans $\text{Comp}(\mathcal{A})$. Soit
$$
(a,b,c): (x_1,y_1,c(f_1), f_1, i_1, p_1) \to (x_2,y_2, c(f_2), f_2, i_1, p_1)
$$
un morphisme quelconque de triangles dans $K(\mathcal{A})$. Si
$a$ et $b$ sont des équivalences d'homotopie, alors il en est de même pour $c$.
\end{lemma}

\begin{proof}
Puisque $a$ et $b$ sont des équivalences d'homotopie, elles sont
inversibles dans $K(\mathcal{A})$, notons donc $a^{-1}$ et $b^{-1}$ leurs inverses dans
$K(\mathcal{A})$, ce qui nous donne un diagramme commutatif
$$
\xymatrix{
x_2\ar[d]^{a^{-1}}\ar[r]^{f_2} &
y_2\ar[d]^{b^{-1}}\ar[r]^{i_2} &
c(f_2)\ar[d]^{c'} \\
x_1\ar[r]^{f_1} & y_1
\ar[r]^{i_1} &
c(f_1)
}
$$
où l'application $c'$ est définie via le lemme \ref{lemma-cone} appliqué au carré
commutatif de gauche du diagramme ci-dessus. Puisque le diagramme commute
dans $K(\mathcal{A})$, il suffit, d'après le lemme \ref{lemma-triseq}, de
prouver ce qui suit : étant donné un morphisme de triangles
$(1,1,c): (x,y,c(f),f,i,p)\to (x,y,c(f),f,i,p)$ dans
$K(\mathcal{A})$, l'application $c$ est un isomorphisme dans
$K(\mathcal{A})$. Nous avons les diagrammes commutatifs dans $K(\mathcal{A})$:
$$
\vcenter{
\xymatrix{
y\ar[d]^{1}\ar[r] &
c(f)\ar[d]^{c}\ar[r] &
x[1]\ar[d]^{1} \\
y\ar[r] & c(f)
\ar[r] & x[1] }
}
\quad\Rightarrow\quad
\vcenter{
\xymatrix{
y\ar[d]^{0}\ar[r]
&
c(f)\ar[d]^{c-1}\ar[r]
& x[1]\ar[d]^{0}
\\ y\ar[r] & c(f)
\ar[r]
&
x[1]
}
}
$$
Comme les lignes sont des suites exactes courtes admissibles, nous
obtenons l'identité $(c-1)^2 = 0$ par le lemme \ref{lemma-triseq},
dont nous concluons que $2-c$ est l'inverse de $c$ dans $K(\mathcal{A})$ de
sorte que $c$ est un isomorphisme.
\end{proof}

\noindent
Le lemme suivant est l'analogue du
lemme \ref{lemma-the-same-up-to-isomorphisms}.

\begin{lemma}
\label{lemma-cone-homotopy}
Dans la situation \ref{situation-ABC}.
\begin{enumerate}
\item Soit une suite exacte courte admissible
$x\xrightarrow{\alpha} y\xrightarrow{\beta} z$. Il
existe alors une équivalence d'homotopie
$e:C(\alpha)\to z$ telle que le diagramme
\begin{equation}
\label{equation-cone-isom-triangle}
\vcenter{
\xymatrix{
x\ar[r]^{\alpha}\ar[d] &
y\ar[r]^{b}\ar[d] &
C(\alpha)\ar[r]^{-c}\ar@{.>}[d]^{e} &
x[1]\ar[d] \\
x\ar[r]^{\alpha} &
y\ar[r]^{\beta} &
z\ar[r]^{\delta} & x[1]
}
}
\end{equation}
définit un isomorphisme de triangles dans $K(\mathcal{A})$. Ici,
$y\xrightarrow{b}C(\alpha)\xrightarrow{c}x[1]$ est
la suite exacte courte admissible donnée comme dans l'axiome (C).
\item Étant donné un
morphisme $\alpha : x \to y$ dans
$\text{Comp}(\mathcal{A})$, soit $x \xrightarrow{\tilde{\alpha}} \tilde{y} \to y$
la factorisation donnée comme dans le lemme \ref{lemma-factor}, où le
monomorphisme admissible $x \xrightarrow{\tilde{\alpha}} y$ s'étend à la
suite exacte courte admissible
$$
\xymatrix{
x \ar[r]^{\tilde{\alpha}} &
\tilde{y} \ar[r] & z
}
$$
Alors il existe un isomorphisme de triangles
$$
\xymatrix{
x \ar[r]^{\tilde{\alpha}} \ar[d] &
\tilde{y} \ar[r] \ar[d] & z
\ar[r]^{\delta} \ar@{.>}[d]^{e} & x[1]
\ar[d] \\ x
\ar[r]^{\alpha} & y
\ar[r] &
C(\alpha) \ar[r]^{-c} &
x[1]
}
$$
où le triangle supérieur est le
triangle associé à la suite $x
\xrightarrow{\tilde{\alpha}} \tilde{y} \to z$.
\end{enumerate}
\end{lemma}

\begin{proof}
Pour (1), nous considérons le diagramme plus complet, \emph{sans} le
changement de signe sur $c$ :
$$
\xymatrix{
x\ar@<0.5ex>[r]^{\alpha} \ar[d] &
y\ar@<0.5ex>[l]^{\pi} \ar@<0.5ex>[r]^{b}\ar[d] &
C(\alpha)\ar@<0.5ex>[l]^{p} \ar@<0.5ex>[r]^{c}\ar@{.>}@<0.5ex>[d]^{e} &
x[1]\ar@<0.5ex>[l]^{\sigma} \ar[d]\ar@<0.5ex>[r]^{\alpha} &
y[1]\ar@<0.5ex>[l]^{\pi} \\
x\ar@<0.5ex>[r]^{\alpha} &
y\ar@<0.5ex>[r]^{\beta} \ar@<0.5ex>[l]^{\pi} &
z\ar[r]^{\delta}\ar@<0.5ex>[l]^{s} \ar@{.>}@<0.5ex>[u]^{f} &
x[1]
}
$$
où la suite exacte courte admissible
$x\xrightarrow{\alpha} y\xrightarrow{\beta} z$ est
munie du scindage $\pi$, $s$, et la suite exacte courte admissible
$y\xrightarrow{b}C(\alpha)\xrightarrow{c}x[1]$ est munie du scindage $p$, $\sigma$.
Notons que (en identifiant les ensembles de morphismes par décalage)
$$
\alpha = pd(\sigma) = -d(p)\sigma,\quad
\delta = \pi d(s) = -d(\pi)s
$$
par la construction dans le lemme \ref{lemma-get-triangle}.

\medskip\noindent
Nous définissons $e=\beta p$ et $f=bs-\sigma\delta$. Nous vérifions d'abord qu'ils
sont des morphismes dans $\text{Comp}(\mathcal{A})$. Pour montrer que $d(e)=\beta
d(p)$ s'annule, il suffit de montrer que $\beta d(p)b$ et $\beta d(p)\sigma$
s'annulent tous les deux ; or
$$
\beta d(p)b = \beta d(pb) = \beta d(1_y) = 0,\quad
\beta d(p)\sigma = -\beta\alpha = 0
$$
De même, pour vérifier que $d(f)=bd(s)-d(\sigma)\delta$ s'annule, il
suffit de vérifier que les post-compositions par $p$ et $c$ s'annulent toutes
deux ; or
\begin{align*}
pbd(s) - pd(\sigma)\delta
= &
d(s)-\alpha\delta = d(s)-\alpha\pi d(s) = 0 \\
cbd(s)-cd(\sigma)\delta =
&
-cd(\sigma)\delta = -d(c\sigma)\delta = 0
\end{align*}
La commutativité des deux carrés de gauche du
diagramme \ref{equation-cone-isom-triangle} découle directement de la définition.
Avant de prouver la commutativité du carré de droite (à homotopie près), nous
vérifions d'abord que $e$ est une équivalence d'homotopie. Clairement,
$$
ef=\beta p (bs-\sigma\delta)=\beta s=1_z
$$
Pour vérifier que $fe$ est homotope à $1_{C(\alpha)}$, nous observons d'abord
$$
b\alpha = bpd(\alpha) = d(\sigma),\quad
\alpha c = -d(p)\sigma c = -d(p),\quad
d(\pi)p = d(\pi)s\beta p = -\delta\beta p
$$
En utilisant ces identités, nous calculons
\begin{align*}
1_{C(\alpha)} = &
bp + \sigma c
\quad (\text{d'après }y \xrightarrow{b} C(\alpha) \xrightarrow{c} x[1]) \\ = &
b(\alpha\pi
+ s\beta)p + \sigma(\pi\alpha)c \quad
(\text{d'après }x \xrightarrow{\alpha} y \xrightarrow{\beta} z) \\ = &
d(\sigma)\pi
p + bs\beta p - \sigma\pi d(p) \quad
(\text{d'après les deux premières identités ci-dessus}) \\ = &
d(\sigma)\pi
p + bs\beta p - \sigma\delta\beta p +
\sigma\delta\beta p - \sigma\pi d(p) \\ = & (bs -
\sigma\delta)\beta
p + d(\sigma)\pi p - \sigma d(\pi)p -
\sigma\pi d(p)\quad (\text{d'après la
troisième identité ci-dessus}) \\ = & fe +
d(\sigma
\pi p)
\end{align*}
puisque $\sigma \in \Hom^{-1}(x, C(\alpha))$
(cf. démonstration du lemme \ref{lemma-id-cone-null}).
Par conséquent, $e$ et $f$ sont des inverses
homotopiques. Enfin, pour vérifier que le carré
droit du diagramme \ref{equation-cone-isom-triangle} commute à homotopie
près, il suffit de vérifier que $-cf=\delta$. Ceci découle de
$$
-cf = -c(bs-\sigma\delta) = c\sigma\delta = \delta
$$
puisque $cb=0$.

\medskip\noindent
Pour (2), considérons la factorisation
$x\xrightarrow{\tilde{\alpha}}\tilde{y}\to y$
donnée comme dans le lemme \ref{lemma-factor}, donc le deuxième
morphisme est une équivalence d'homotopie. D'après les lemmes
\ref{lemma-cone} et \ref{lemma-analogue-third-isomorphism},
il existe un isomorphisme de triangles entre
$$
x \xrightarrow{\alpha} y \to C(\alpha) \to x[1]
\quad\text{et}\quad
x \xrightarrow{\tilde{\alpha}} \tilde{y} \to C(\tilde{\alpha}) \to x[1]
$$
Puisque nous pouvons composer des isomorphismes de triangles, en
remplaçant $\alpha$ par $\tilde{\alpha}$, $y$ par $\tilde{y}$, et $C(\alpha)$
par $C(\tilde{\alpha})$, nous pouvons supposer que $\alpha$ est un monomorphisme
admissible. Dans ce cas, le résultat découle de (1).
\end{proof}

\noindent
Le lemme suivant est l'analogue du
lemme \ref{lemma-homotopy-category-pre-triangulated}.

\begin{lemma}
\label{lemma-analogue-homotopy-category-pre-triangulated}
Dans la situation \ref{situation-ABC}, la catégorie homotopique
$K(\mathcal{A})$ avec ses foncteurs de translation naturels et ses triangles distingués
est une catégorie pré-triangulée.
\end{lemma}

\begin{proof}
Nous allons vérifier les axiomes TR1, TR2 et TR3.

\medskip\noindent
Démonstration de TR1. Par définition, tout triangle isomorphe à un triangle
distingué est distingué. Puisque
$$
\xymatrix{x\ar[r]^{1_x} & x\ar[r] & 0}
$$
est une suite exacte courte admissible, $(x, x, 0, 1_x, 0, 0)$
est un triangle distingué. De plus, étant donné un morphisme
$\alpha : x \to y$ dans $\text{Comp}(\mathcal{A})$, le triangle
donné par $(x, y, c(\alpha), \alpha, i, -p)$ est distingué d'après le
lemme \ref{lemma-cone-homotopy}.

\medskip\noindent
Démonstration de TR2. Soit $(x,y,z,\alpha,\beta,\gamma)$ un triangle
et supposons que $(y,z,x[1],\beta,\gamma,-\alpha[1])$ soit
distingué. Alors il existe une suite exacte courte admissible
$0 \to x' \to y' \to z' \to 0$ telle que le triangle associé
$(x',y',z',\alpha',\beta',\gamma')$ soit isomorphe à
$(y,z,x[1],\beta,\gamma,-\alpha[1])$. Après rotation, nous concluons que
$(x,y,z,\alpha,\beta,\gamma)$ est isomorphe à $(z'[-1],x',y',
\gamma'[-1], \alpha',\beta')$. D'après le lemme
\ref{lemma-cone-rotate-isom}, nous déduisons
que $(z'[-1],x',y', \gamma'[-1], \alpha',\beta')$ est isomorphe à
$(z'[-1],x',c(\gamma'[-1]), \gamma'[-1], i, p)$. En composant les
deux isomorphismes avec les changements de signe indiqués dans le
diagramme suivant :
$$
\xymatrix@C=3pc{
x\ar[r]^{\alpha}\ar[d] &
y\ar[r]^{\beta}\ar[d] &
z\ar[r]^{\gamma}\ar[d] &
x[1]\ar[d] \\
z'[-1]\ar[r]^{-\gamma'[-1]}\ar[d]_{-1_{z'[-1]}} & x
\ar[r]^{\alpha'}\ar@{=}[d] & y'
\ar[r]^{\beta'} \ar[d] &
z'\ar[d]^{-1_{z'}} \\
z'[-1]\ar[r]^{\gamma'[-1]} & x
\ar[r]^{\alpha'} &
c(\gamma'[-1]) \ar[r]^{-p} &
z'
}
$$
Nous concluons que $(x,y,z,\alpha,\beta,\gamma)$ est distingué par
le lemme \ref{lemma-cone-homotopy} (2). Inversement, supposons que
$(x,y,z,\alpha,\beta,\gamma)$ soit distingué, de sorte que
selon le lemme \ref{lemma-cone-homotopy} (1), il soit isomorphe
à un triangle de la forme $(x',y', c(\alpha'), \alpha', i,
-p)$ pour un certain morphisme $\alpha': x' \to y'$ dans
$\text{Comp}(\mathcal{A})$. Le triangle obtenu par rotation
$(y,z,x[1],\beta,\gamma, -\alpha[1])$ est isomorphe au triangle $(y',c(\alpha'), x'[1], i, -p,
-\alpha[1])$ qui est isomorphe à $(y',c(\alpha'), x'[1], i, p,
\alpha[1])$. D'après le lemme \ref{lemma-restate-axiom-c}, ce triangle est
distingué, d'où il suit que $(y,z,x[1], \beta,\gamma, -\alpha[1])$ est
distingué.

\medskip\noindent
Démonstration de TR3 : Supposons que $(x,y,z,
\alpha,\beta,\gamma)$ et $(x',y',z',\alpha',\beta',\gamma')$ soient des triangles
distingués de $\text{Comp}(\mathcal{A})$ et que $f: x \to x'$ et $g: y
\to y'$ soient des morphismes tels que
$\alpha' \circ f = g \circ \alpha$. D'après le
lemme \ref{lemma-cone-homotopy}, nous pouvons supposer que
$(x,y,z,\alpha,\beta,\gamma)= (x,y,c(\alpha),\alpha, i, -p)$ et
$(x',y',z', \alpha',\beta',\gamma')= (x',y',c(\alpha'), \alpha',i',-p')$. Appliquons maintenant
le lemme \ref{lemma-cone} et nous
avons terminé.
\end{proof}

\noindent
Le lemme suivant est l'analogue du lemme \ref{lemma-two-split-injections}.

\begin{lemma}
\label{lemma-dgc-analogue-tr4}
Dans la situation \ref{situation-ABC}, étant donnés des monomorphismes
admissibles $x \xrightarrow{\alpha} y$, $y \xrightarrow{\beta} z$ dans
$\mathcal{A}$, il existe des triangles
distingués $(x,y,q_1,\alpha,p_1,\delta_1)$, $(x,z,q_2,\beta\alpha,p_2,\delta_2)$ et
$(y,z,q_3,\beta,p_3,\delta_3)$ pour lesquels TR4 est vérifié.
\end{lemma}

\begin{proof}
Étant donnés les monomorphismes admissibles $x\xrightarrow{\alpha}
y$ et $y\xrightarrow{\beta}z$, nous pouvons trouver des triangles
distingués, en les prolongeant en suites exactes courtes admissibles,
$$
\xymatrix{
x\ar@<0.5ex>[r]^{\alpha} &
y\ar@<0.5ex>[l]^{\pi_1} \ar@<0.5ex>[r]^{p_1} &
q_1 \ar[r]^{\delta_1} \ar@<0.5ex>[l]^{s_1} &
x[1]
}
$$
$$
\xymatrix{
x\ar@<0.5ex>[r]^{\beta\alpha} &
z\ar@<0.5ex>[l]^{\pi_1\pi_3} \ar@<0.5ex>[r]^{p_2} &
q_2 \ar[r]^{\delta_2} \ar@<0.5ex>[l]^{s_2} &
x[1]
}
$$
$$
\xymatrix{
y\ar@<0.5ex>[r]^{\beta} &
z\ar@<0.5ex>[l]^{\pi_3} \ar@<0.5ex>[r]^{p_3} &
q_3 \ar[r]^{\delta_3} \ar@<0.5ex>[l]^{s_3} &
x[1]
}
$$
Dans ces diagrammes, les applications $\delta_i$ sont définies comme
$\delta_i = \pi_i d(s_i)$ de manière analogue aux applications définies dans le lemme
\ref{lemma-get-triangle}. Elles s'insèrent dans
le diagramme commutatif suivant, dont les flèches sont en traits pleins
$$
\xymatrix@C=5pc@R=3pc{
x\ar@<0.5ex>[r]^{\alpha} \ar@<0.5ex>[dr]^{\beta\alpha} &
y\ar@<0.5ex>[d]^{\beta} \ar@<0.5ex>[l]^{\pi_1} \ar@<0.5ex>[r]^{p_1} &
q_1 \ar[r]^{\delta_1} \ar@<0.5ex>[l]^{s_1} \ar@{.>}[dd]^{p_2\beta s_1} &
x[1] \\ & z
\ar@<0.5ex>[u]^{\pi_3}\ar@<0.5ex>[d]^{p_3}
\ar@<0.5ex>[dr]^{p_2}
\ar@<0.5ex>[ul]^{\pi_1\pi_3} & & \\ &
q_3\ar@<0.5ex>[u]^{s_3}
\ar[d]^{\delta_3} & q_2
\ar@{.>}[l]^{p_3s_2} \ar@<0.5ex>[ul]^{s_2} \ar[dr]^{\delta_2} \\ & y[1] &
&
x[1]}
$$
où nous avons défini les flèches en pointillés comme indiqué.
Clairement, leur composition $p_3s_2p_2\beta s_1 = 0$ puisque
$s_2p_2 = 0$. Nous affirmons qu'elles sont toutes deux des morphismes de
$\text{Comp}(\mathcal{A})$. Nous pouvons vérifier cela en utilisant les équations
du lemme \ref{lemma-get-triangle} :
$$
d(p_2\beta s_1) = p_2\beta d(s_1) = p_2\beta\alpha\pi_1 d(s_1) = 0
$$
puisque $p_2\beta\alpha = 0$, et
$$
d(p_3s_2) = p_3d(s_2) = p_3\beta\alpha\pi_1\pi_3 d(s_2) = 0
$$
puisque $p_3\beta = 0$. Pour vérifier que $q_1\to q_2\to
q_3$ est une suite exacte courte admissible, il reste à montrer
que, dans la catégorie graduée sous-jacente, $q_2 = q_1\oplus q_3$
avec les deux morphismes ci-dessus comme coprojection et projection.
Pour ce faire, observons que dans la catégorie graduée
sous-jacente $\mathcal{C}$, il y a
$$
y = x\oplus q_1,\quad
z = y\oplus q_3 = x\oplus q_1\oplus q_3
$$
où $\pi_1\pi_3$ donne le morphisme de projection sur le premier
facteur : $x\oplus q_1\oplus q_3\to z$. Par l'axiome (A) sur
$\mathcal{A}$, $\mathcal{C}$ est une catégorie additive, par conséquent
nous pouvons
appliquer Homologie, lemme \ref{homology-lemma-additive-cat-biproduct-kernel}
et conclure que
$$
\Ker(\pi_1\pi_3) = q_1\oplus q_3
$$
dans $\mathcal{C}$. Une autre application
d'Homologie, lemme \ref{homology-lemma-additive-cat-biproduct-kernel} à $z =
x\oplus q_2$ donne $\Ker(\pi_1\pi_3) = q_2$. Par conséquent
$q_2\cong q_1\oplus q_3$ dans $\mathcal{C}$. Il est clair
que les morphismes en pointillés définis ci-dessus donnent la
coprojection et la projection.

\medskip\noindent
Enfin, nous devons vérifier que le morphisme
$\delta : q_3 \to q_1[1]$ induit par la suite exacte
courte admissible $q_1\to q_2\to q_3$ coïncide avec
$p_1\delta_3$. Selon la construction du
lemme \ref{lemma-get-triangle}, le morphisme $\delta$ est donné par
\begin{align*}
p_1\pi_3s_2d(p_2s_3)
= &
p_1\pi_3s_2p_2d(s_3) \\ =
&
p_1\pi_3(1-\beta\alpha\pi_1\pi_3)d(s_3) \\ = &
p_1\pi_3d(s_3)\quad
(\text{puisque }\pi_3\beta = 0) \\ =
&
p_1\delta_3
\end{align*}
comme souhaité. La démonstration est complète.
\end{proof}

\noindent
En rassemblant tout, nous obtenons finalement l'analogue
de la proposition \ref{proposition-homotopy-category-triangulated}.

\begin{proposition}
\label{proposition-ABC-homotopy-category-triangulated}
Dans la situation \ref{situation-ABC}, la catégorie homotopique
$K(\mathcal{A})$ avec ses foncteurs de translation naturels et ses triangles distingués est
une catégorie triangulée.
\end{proposition}

\begin{proof}
D'après le lemme \ref{lemma-analogue-homotopy-category-pre-triangulated}, nous
savons que $K(\mathcal{A})$ est pré-triangulée. En
combinant les lemmes \ref{lemma-analogue-sequence-maps-split}
et \ref{lemma-dgc-analogue-tr4}
avec Catégories dérivées, lemme \ref{derived-lemma-easier-axiom-four},
nous concluons que $K(\mathcal{A})$ est une catégorie triangulée.
\end{proof}

\begin{lemma}
\label{lemma-functor-between-ABC}
Soit $R$ un anneau. Soit $F : \mathcal{A} \to \mathcal{B}$ un foncteur
entre des catégories différentielles graduées sur $R$ satisfaisant aux
axiomes (A), (B) et (C), tel que $F(x[1]) =
F(x)[1]$. Alors $F$ induit un foncteur
exact $K(\mathcal{A}) \to K(\mathcal{B})$ de catégories triangulées.
\end{lemma}

\begin{proof}
En effet, si $x \to y \to z$ est une suite exacte courte
admissible dans $\text{Comp}(\mathcal{A})$, alors $F(x) \to F(y) \to
F(z)$ est une suite exacte courte admissible dans
$\text{Comp}(\mathcal{B})$. De plus, le morphisme de liaison $\delta = \pi\text{d}(s) : z \to x[1]$
construit dans le lemme \ref{lemma-get-triangle} produit le morphisme
$F(\delta) : F(z) \to F(x[1]) = F(x)[1]$ qui est égal au morphisme de
liaison $F(\pi) \text{d}(F(s))$ pour la suite exacte courte admissible $F(x)
\to F(y) \to F(z)$.
\end{proof}





\section{Bimodules}
\label{section-bimodules}

\noindent
Nous continuons la discussion commencée dans la section \ref{section-tensor-product}.

\begin{definition}
\label{definition-bimodule}
Bimodules. Soit $R$ un anneau.
\begin{enumerate}
\item Soient $A$ et $B$ des $R$-algèbres. Un {\it $(A, B)$-bimodule}
est un $R$-module $M$ muni d'applications $R$-bilinéaires
$$
A \times M \to M, (a, x) \mapsto ax
\quad\text{et}\quad
M \times B \to M, (x, b) \mapsto xb
$$
vérifiant les propriétés suivantes
\begin{enumerate}
\item $a'(ax) = (a'a)x$ et $(xb)b' = x(bb')$,
\item $a(xb) = (ax)b$, et
\item $1 x = x = x 1$.
\end{enumerate}
\item Soient $A$ et $B$ des algèbres graduées par $\mathbf{Z}$ sur
$R$. Un {\it $(A, B)$-bimodule gradué} est un $(A, B)$-bimodule $M$ qui
possède une graduation $M = \bigoplus M^n$
telle que $A^n M^m \subset M^{n + m}$ et $M^n B^m \subset M^{n + m}$.
\item Soient $A$ et $B$ des algèbres différentielles graduées
sur $R$. Un {\it $(A, B)$-bimodule différentiel gradué} est un $(A,
B)$-bimodule gradué qui est muni d'une
différentielle $\text{d} : M \to M$ homogène de degré $1$
telle que $\text{d}(ax) = \text{d}(a)x + (-1)^{\deg(a)}a\text{d}(x)$ et
$\text{d}(xb) = \text{d}(x)b + (-1)^{\deg(x)}x\text{d}(b)$ pour
des éléments homogènes $a \in A$, $x \in M$, $b \in B$.
\end{enumerate}
\end{definition}

\noindent
Remarquons qu'un $(A, B)$-bimodule différentiel gradué $M$
est la même chose qu'un module différentiel gradué à droite
sur $B$ qui est également un module différentiel gradué à
gauche sur $A$, avec des graduations et des différentielles qui
coïncident, et tel que la structure de module sur $A$ commute
avec la structure de module sur $B$. Voici un énoncé précis.

\begin{lemma}
\label{lemma-what-makes-a-bimodule-dg}
Soit $R$ un anneau. Soient $(A, \text{d})$ et $(B, \text{d})$ des
algèbres différentielles graduées sur $R$. Soit $M$ un $B$-module
différentiel gradué droit. Il existe une correspondance $1$ à $1$ entre
les structures de bimodule $(A, B)$ sur $M$ compatibles avec la structure
de $B$-module différentiel gradué donnée et les homomorphismes
$$
A
\longrightarrow
\Hom_{\text{Mod}^{dg}_{(B, \text{d})}}(M, M)
$$
d'algèbres différentielles graduées sur $R$.
\end{lemma}

\begin{proof}
Soit $\mu : A \times M \to M$ une application définissant une structure de
$A$-module différentiel gradué à gauche sur le complexe de $R$-modules $M^\bullet$
sous-jacent à $M$. D'après le lemme \ref{lemma-left-module-structure}, la structure $\mu$
correspond à une application $\gamma : A \to \Hom^\bullet(M^\bullet,
M^\bullet)$ d'algèbres différentielles graduées sur $R$. L'assertion du lemme est
simplement que $\mu$ commute avec l'action de $B$ si et seulement si $\gamma$ est à
valeurs dans
$$
\Hom_{\text{Mod}^{dg}_{(B, \text{d})}}(M, M) \subset
\Hom^\bullet(M^\bullet, M^\bullet)
$$
Nous omettons le calcul détaillé.
\end{proof}

\noindent
Soit $M$ un $(A, B)$-bimodule différentiel gradué. Rappelons d'après la
section \ref{section-left-modules} que la structure de module différentiel
gradué à gauche sur $A$ correspond à une structure de module différentiel
gradué à droite sur $A^{opp}$. Comme les structures de module sur $A$ et sur $B$
commutent, cela donne à $M$ la structure d'un module différentiel
gradué sur $A^{opp} \otimes_R B$ :
$$
x \cdot (a \otimes b) = (-1)^{\deg(a)\deg(x)} axb
$$
Inversement, si nous avons un $A^{opp} \otimes_R B$-module différentiel gradué $M$,
alors nous pouvons utiliser la formule ci-dessus pour obtenir un $(A, B)$-bimodule
différentiel gradué.

\begin{lemma}
\label{lemma-bimodule-over-tensor}
Soit $R$ un anneau. Soient $(A, \text{d})$ et $(B, \text{d})$
des algèbres différentielles graduées sur $R$. La construction
ci-dessus définit une équivalence de catégories
$$
\begin{matrix}
\text{bimodules différentiels
gradués sur}\\ (A, B)\text{}
\end{matrix}
\longleftrightarrow
\begin{matrix}
\text{modules différentiels gradués à
droite sur }\\ A^{opp} \otimes_R B\text{}
\end{matrix}
$$
\end{lemma}

\begin{proof}
Cela résulte immédiatement de la discussion ci-dessus.
\end{proof}

\noindent
Soit $R$ un anneau. Soient $(A, \text{d})$ et $(B, \text{d})$ des
$R$-algèbres différentielles
graduées. Soit $P$ un bimodule différentiel gradué sur $(A, B)$. On dit
que $P$ {\it a la propriété (P)} s'il existe une filtration
$$
0 = F_{-1}P \subset F_0P \subset F_1P \subset \ldots \subset P
$$
par des $(A, B)$-bimodules différentiels gradués, telle que
\begin{enumerate}
\item $P = \bigcup F_pP$,
\item les inclusions $F_iP \to F_{i + 1}P$ sont scindées comme morphismes
de $(A, B)$-bimodules gradués,
\item les quotients $F_{i + 1}P/F_iP$ sont isomorphes comme $(A,
B)$-bimodules différentiels gradués à une somme directe de $(A \otimes_R B)[k]$.
\end{enumerate}

\begin{lemma}
\label{lemma-bimodule-resolve}
Soit $R$ un anneau. Soient $(A, \text{d})$ et $(B, \text{d})$ des
$R$-algèbres différentielles graduées. Soit $M$ un bimodule différentiel
gradué sur $(A, B)$. Il existe un homomorphisme $P \to M$ de
bimodules différentiels gradués sur $(A, B)$ qui est un quasi-isomorphisme tel que
$P$ possède la propriété (P) comme définie ci-dessus.
\end{lemma}

\begin{proof}
Cela résulte immédiatement des lemmes
\ref{lemma-bimodule-over-tensor} et \ref{lemma-resolve}.
\end{proof}

\begin{lemma}
\label{lemma-bimodule-property-P-sequence}
Soit $R$ un anneau. Soient $(A, \text{d})$ et $(B, \text{d})$ des
$R$-algèbres différentielles graduées. Soit $P$ un
$(A, B)$-bimodule différentiel gradué possédant la propriété (P)
avec la filtration correspondante $F_\bullet$, alors nous obtenons une
suite exacte courte
$$
0 \to
\bigoplus\nolimits F_iP \to
\bigoplus\nolimits F_iP \to P \to 0
$$
de $(A, B)$-bimodules différentiels gradués, qui est scindée comme suite
de $(A, B)$-bimodules gradués.
\end{lemma}

\begin{proof}
Cela résulte immédiatement des lemmes
\ref{lemma-bimodule-over-tensor} et \ref{lemma-property-P-sequence}.
\end{proof}















\section{Bimodules et produit tensoriel}
\label{section-bimodules-tensor}

\noindent
Soit $R$ un anneau. Soient $A$ et $B$ des $R$-algèbres. Soit $M$ un module à
droite sur $A$. Soit $N$ un bimodule sur $(A, B)$. Alors $M
\otimes_A N$ est un module à droite sur $B$.

\medskip\noindent
Si, dans la situation du paragraphe précédent, $A$
et $B$ sont des algèbres graduées par $\mathbf{Z}$,
$M$ est un module gradué sur $A$, et $N$ est un $(A, B)$-bimodule gradué,
alors $M \otimes_A N$ est un module gradué à droite sur $B$. La construction
est fonctorielle en $M$ et définit un foncteur
$$
- \otimes_A N :
\text{Mod}^{gr}_A
\longrightarrow
\text{Mod}^{gr}_B
$$
de catégories graduées comme dans l'exemple \ref{example-gm-gr-cat}. En
effet, si $M$ et $M'$ sont des $A$-modules gradués et que $f : M \to M'$ est un
homomorphisme de $A$-modules homogène de degré $n$,
alors $f \otimes \text{id}_N : M \otimes_A N \to M' \otimes_A N$
est un homomorphisme de $B$-modules homogène de degré $n$.

\medskip\noindent
Si dans la situation du paragraphe précédent $(A,
\text{d})$ et $(B, \text{d})$ sont des algèbres différentielles graduées, $M$
est un module différentiel gradué sur $A$, et $N$ est un bimodule
différentiel gradué sur $(A, B)$, alors $M \otimes_A N$ est un module
différentiel gradué à droite sur $B$.

\begin{lemma}
\label{lemma-tensor}
Soit $R$ un anneau. Soient $(A, \text{d})$ et $(B, \text{d})$
des algèbres différentielles graduées sur $R$. Soit $N$ un
$(A, B)$-bimodule différentiel gradué. Alors
$M \mapsto M \otimes_A N$ définit un foncteur
$$
- \otimes_A N :
\text{Mod}^{dg}_{(A, \text{d})}
\longrightarrow
\text{Mod}^{dg}_{(B, \text{d})}
$$
de catégories différentielles graduées. Ce foncteur induit des foncteurs
$$
\text{Mod}_{(A, \text{d})} \to \text{Mod}_{(B, \text{d})}
\quad\text{et}\quad
K(\text{Mod}_{(A, \text{d})}) \to K(\text{Mod}_{(B, \text{d})})
$$
par une application du lemme \ref{lemma-functorial}.
\end{lemma}

\begin{proof}
Ci-dessus, nous avons vu comment la construction définit un foncteur des
catégories graduées sous-jacentes. Ainsi, il suffit de montrer que la
construction est compatible avec les différentielles. Soient $M$ et $M'$ des
$A$-modules différentiels gradués et soit $f : M \to M'$ un homomorphisme de
$A$-modules homogène de degré $n$. Alors nous avons
$$
\text{d}(f) = \text{d}_{M'} \circ f - (-1)^n f \circ \text{d}_M
$$
D'un autre côté, nous avons
$$
\text{d}(f \otimes \text{id}_N) =
\text{d}_{M' \otimes_A N} \circ
(f \otimes \text{id}_N) -
(-1)^n
(f \otimes \text{id}_N) \circ \text{d}_{M \otimes_A N}
$$
En appliquant ceci à un élément $x \otimes y$ avec $x \in M$ et $y
\in N$ homogènes, nous obtenons
\begin{align*}
\text{d}(f \otimes \text{id}_N)(x \otimes y)
= &
\text{d}_{M'}(f(x)) \otimes y + (-1)^{n + \deg(x)}f(x) \otimes \text{d}_N(y) \\ &
- (-1)^n f(\text{d}_M(x)) \otimes y -
(-1)^{n + \deg(x)}f(x) \otimes \text{d}_N(y) \\ = &
\text{d}(f)
(x \otimes y)
\end{align*}
Ainsi, nous voyons que $\text{d}(f) \otimes \text{id}_N =
\text{d}(f \otimes \text{id}_N)$ et la démonstration est complète.
\end{proof}

\begin{remark}
\label{remark-shift-tensor-no-sign}
Soit $R$ un anneau. Soient $(A, \text{d})$ et $(B, \text{d})$
des algèbres différentielles graduées sur $R$. Soit $N$ un
bimodule différentiel gradué sur $(A, B)$. Soit $M$ un module différentiel gradué
à droite sur $A$. Alors pour chaque $k \in \mathbf{Z}$ il existe
un isomorphisme
$$
(M \otimes_A N)[k] \longrightarrow
M[k] \otimes_A N
$$
de modules différentiels gradués à droite sur $B$, défini sans intervention
de signes, voir Compléments d'algèbre, section \ref{more-algebra-section-sign-rules}.
\end{remark}

\noindent
Si nous avons un anneau $R$ et des $R$-algèbres $A$, $B$, et
$C$, un $A$-module à droite $M$, un $(A, B)$-bimodule $N$, et un
$(B, C)$-bimodule $N'$, alors
$N \otimes_B N'$ est un $(A, C)$-bimodule et nous avons
$$
(M \otimes_A N) \otimes_B N' = M \otimes_A (N \otimes_B N')
$$
Cette égalité reste valable dans le cas gradué et dans le cas différentiel
gradué. Voir Compléments d'algèbre, section \ref{more-algebra-section-sign-rules} pour les
règles des signes.






\section{Bimodules et Hom interne}
\label{section-bimodules-hom}

\noindent
Soit $R$ un anneau. Si $A$ est une $R$-algèbre (voir nos conventions à la
section \ref{section-conventions}) et que $M$, $M'$ sont des $A$-modules à
droite, alors nous définissons
$$
\Hom_A(M, M') = \{f : M \to M' \mid f \text{ est }A\text{-linéaire}\}
$$
comme d'habitude.

\medskip\noindent
Soit $R$ un anneau. Soient $A$ et $B$ des $R$-algèbres. Soit
$N$ un $(A, B)$-bimodule. Soit $N'$ un module à droite sur $B$.
Dans cette situation, nous considérerons
$$
\Hom_B(N, N')
$$
comme $A$-module à droite, par précomposition.

\medskip\noindent
Soit $R$ un anneau. Soient $A$ et $B$ des algèbres graduées par $\mathbf{Z}$ sur
$R$. Soit $N$ un $(A, B)$-bimodule gradué. Soit $N'$ un module gradué à droite sur
$B$. Dans cette situation, nous considérerons le $R$-module gradué
$$
\Hom_{\text{Mod}^{gr}_B}(N, N')
$$
défini dans l'exemple \ref{example-gm-gr-cat} comme un $A$-module gradué à droite
en utilisant la précomposition. La construction est fonctorielle en $N'$ et définit
un foncteur
$$
\Hom_{\text{Mod}^{gr}_B}(N, -) :
\text{Mod}^{gr}_B
\longrightarrow
\text{Mod}^{gr}_A
$$
de catégories graduées comme dans l'exemple \ref{example-gm-gr-cat}. En
effet, si $N_1$ et $N_2$ sont des $B$-modules gradués et que $f : N_1 \to N_2$ est un
homomorphisme de $B$-modules homogène de degré $n$, alors l'application
induite $\Hom_{\text{Mod}^{gr}_B}(N, N_1) \to \Hom_{\text{Mod}^{gr}_B}(N, N_2)$
est un homomorphisme de $A$-modules homogène de degré $n$.

\medskip\noindent
Soit $R$ un anneau. Soient $A$ et $B$ des algèbres différentielles graduées
par $\mathbf{Z}$ sur $R$. Soit $N$ un bimodule différentiel gradué sur $(A,
B)$. Soit $N'$ un module différentiel gradué à droite sur $B$. Dans cette
situation, nous considérerons le $R$-module différentiel gradué
$$
\Hom_{\text{Mod}^{dg}_{(B, \text{d})}}(N, N')
$$
défini dans l'exemple \ref{example-dgm-dg-cat} comme un $A$-module
différentiel gradué à droite en utilisant la précomposition comme dans le cas gradué.
Ceci est compatible avec les différentielles car la multiplication est la
composition
$$
\Hom_{\text{Mod}^{dg}_B}(N, N') \otimes_R A \to
\Hom_{\text{Mod}^{dg}_B}(N, N') \otimes_R
\Hom_{\text{Mod}^{dg}_B}(N, N) \to
\Hom_{\text{Mod}^{dg}_B}(N, N')
$$
La première flèche utilise
l'application du lemme \ref{lemma-what-makes-a-bimodule-dg} et la
deuxième flèche est la composition dans la catégorie différentielle
graduée $\text{Mod}^{dg}_{(B, \text{d})}$.

\begin{lemma}
\label{lemma-hom}
Soit $R$ un anneau. Soient $(A, \text{d})$ et $(B, \text{d})$
des algèbres différentielles graduées sur $R$. Soit $N$ un
$(A, B)$-bimodule différentiel gradué. La construction ci-dessus
définit un foncteur
$$
\Hom_{\text{Mod}^{dg}_{(B, \text{d})}}(N, -) :
\text{Mod}^{dg}_{(B, \text{d})}
\longrightarrow
\text{Mod}^{dg}_{(A, \text{d})}
$$
de catégories différentielles graduées. Ce foncteur induit des foncteurs
$$
\text{Mod}_{(B, \text{d})} \to \text{Mod}_{(A, \text{d})}
\quad\text{et}\quad
K(\text{Mod}_{(B, \text{d})}) \to K(\text{Mod}_{(A, \text{d})})
$$
par une application du lemme \ref{lemma-functorial}.
\end{lemma}

\begin{proof}
Ci-dessus, nous avons vu comment la construction définit un foncteur des
catégories graduées sous-jacentes. Il suffit donc de montrer que la construction est
compatible avec les différentielles. Soient $N_1$ et $N_2$ des $B$-modules
différentiels gradués. Écrivons
$$
H_{12} = \Hom_{\text{Mod}^{dg}_{(B, \text{d})}}(N_1, N_2),\quad
H_1 = \Hom_{\text{Mod}^{dg}_{(B, \text{d})}}(N, N_1),\quad
H_2 = \Hom_{\text{Mod}^{dg}_{(B, \text{d})}}(N, N_2)
$$
Considérons la composition
$$
c : H_{12} \otimes_R H_1 \longrightarrow H_2
$$
dans la catégorie différentielle graduée $\text{Mod}^{dg}_{(B,
\text{d})}$. Soit $f : N_1 \to N_2$ un homomorphisme de $B$-modules qui est
homogène de degré $n$, en d'autres termes, $f \in
H_{12}^n$. Le foncteur dans le lemme envoie $f$ sur $c_f : H_1 \to H_2$, $g \mapsto c(f,
g)$. De même pour $\text{d}(f)$. D'autre part, la différentielle sur
$$
\Hom_{\text{Mod}^{dg}_{(A, \text{d})}}(H_1, H_2)
$$
envoie $c_f$ sur $\text{d}_{H_2} \circ c_f - (-1)^n c_f \circ \text{d}_{H_1}$. Comme
$c$ est un morphisme de complexes de $R$-modules, nous avons
$\text{d} c(f, g) = c(\text{d}f, g) + (-1)^n c(f, \text{d}g)$. Par
conséquent, nous voyons que
\begin{align*}
(\text{d}c_f)(g)
& =
\text{d}c(f,g) - (-1)^n c(f, \text{d}g) \\
& =
c(\text{d}f, g) + (-1)^n c(f, \text{d}g) - (-1)^n c(f, \text{d}g) \\ & =
c(\text{d}f,
g) = c_{\text{d}f}(g)
\end{align*}
et la démonstration est complète.
\end{proof}

\begin{remark}
\label{remark-shift-hom-no-sign}
Soit $R$ un anneau. Soient $(A, \text{d})$ et $(B, \text{d})$
des algèbres différentielles graduées sur $R$. Soit $N$ un
bimodule différentiel gradué sur $(A, B)$. Soit $N'$ un module différentiel gradué
à droite sur $B$. Alors pour chaque $k \in \mathbf{Z}$ il existe
un isomorphisme
$$
\Hom_{\text{Mod}^{gr}_B}(N, N')[k]
\longrightarrow
\Hom_{\text{Mod}^{gr}_B}(N, N'[k])
$$
de modules différentiels gradués à droite sur $A$, défini sans intervention
de signes, voir Compléments d'algèbre, section \ref{more-algebra-section-sign-rules}.
\end{remark}

\begin{lemma}
\label{lemma-tensor-hom-adjunction}
Soit $R$ un anneau. Soient $A$ et $B$ des $R$-algèbres.
Soit $M$ un module à droite sur $A$, $N$ un bimodule sur $(A, B)$, et $N'$
un module à droite sur $B$. Alors nous avons un isomorphisme canonique
$$
\Hom_B(M \otimes_A N, N') = \Hom_A(M, \Hom_B(N, N'))
$$
de $R$-modules.
Si $A$, $B$, $M$, $N$, $N'$ sont compatiblement gradués, alors nous avons un
isomorphisme canonique
$$
\Hom_{\text{Mod}_B^{gr}}(M \otimes_A N, N') =
\Hom_{\text{Mod}_A^{gr}}(M, \Hom_{\text{Mod}_B^{gr}}(N, N'))
$$
de $R$-modules gradués. Si
$A$, $B$, $M$, $N$, $N'$ sont munis de structures différentielles graduées compatibles,
alors nous avons un isomorphisme canonique
$$
\Hom_{\text{Mod}^{dg}_{(B, \text{d})}}(M \otimes_A N, N') =
\Hom_{\text{Mod}^{dg}_{(A, \text{d})}}(M,
\Hom_{\text{Mod}^{dg}_{(B, \text{d})}}(N, N'))
$$
de complexes de $R$-modules.
\end{lemma}

\begin{proof}
Démonstration omise. Indication : dans le cas non gradué, interpréter les deux membres
comme des applications $A$-bilinéaires $\psi : M \times N \to N'$ qui sont
$B$-linéaires à droite. Dans le cas (différentiel) gradué, utiliser
l'isomorphisme de Compléments d'algèbre, lemme
\ref{more-algebra-lemma-compose} et vérifier qu'il est compatible avec les structures
de modules. On peut aussi utiliser l'isomorphisme du lemme
\ref{lemma-characterize-hom} et montrer qu'il est compatible avec les structures de $B$-modules.
\end{proof}







\section{Hom dérivé}
\label{section-restriction}

\noindent
Cette section est analogue à
Compléments d'algèbre, section \ref{more-algebra-section-RHom}.

\medskip\noindent
Soit $R$ un anneau. Soient $(A, \text{d})$ et $(B, \text{d})$
des algèbres différentielles graduées sur $R$. Soit $N$ un
bimodule différentiel gradué $(A, B)$. Considérons le foncteur
\begin{equation}
\label{equation-restriction}
\Hom_{\text{Mod}^{dg}_{(B, \text{d})}}(N, -) :
\text{Mod}_{(B, \text{d})}
\longrightarrow
\text{Mod}_{(A, \text{d})}
\end{equation}
de la section \ref{section-bimodules-hom}.

\begin{lemma}
\label{lemma-restriction-homotopy}
Le foncteur (\ref{equation-restriction}) définit un foncteur
exact $K(\text{Mod}_{(B, \text{d})}) \to K(\text{Mod}_{(A, \text{d})})$
de catégories triangulées.
\end{lemma}

\begin{proof}
Via le lemme \ref{lemma-hom}
et la remarque
\ref{remark-shift-hom-no-sign}, cela découle du principe général du
lemme \ref{lemma-functor-between-ABC}.
\end{proof}

\noindent
Rappelons que nous avons un foncteur
exact de catégories triangulées
$$
\Hom_{\text{Mod}^{dg}_{(B, \text{d})}}(N, -) :
K(\text{Mod}_{(B, \text{d})}) \to K(\text{Mod}_{(A, \text{d})})
$$
voir le lemme \ref{lemma-restriction-homotopy}. Considérons le diagramme
$$
\xymatrix{
K(\text{Mod}_{(B, \text{d})}) \ar[d] \ar[rr]_{\text{voir ci-dessus}} \ar[rrd]_F & &
K(\text{Mod}_{(A, \text{d})}) \ar[d] \\
D(B, \text{d}) \ar@{..>}[rr] & &
D(A, \text{d})
}
$$
Nous aimerions construire une flèche pointillée en
tant que {\it foncteur dérivé à droite} de la composition
$F$. ({\it Avertissement} : dans la plupart des cas intéressants, le diagramme ne
commutera pas.) Plus précisément, dans le
cadre général de Catégories dérivées, section
\ref{derived-section-derived-functors}, nous voulons
calculer le foncteur dérivé à droite de $F$ par rapport au système multiplicatif des
quasi-isomorphismes dans $K(\text{Mod}_{(A, \text{d})})$.

\begin{lemma}
\label{lemma-derived-restriction}
Dans la situation ci-dessus, le foncteur dérivé à droite de $F$
existe. Nous le notons $R\Hom(N, -) : D(B, \text{d}) \to D(A, \text{d})$.
\end{lemma}

\begin{proof}
Nous
utiliserons Catégories dérivées, lemme \ref{derived-lemma-find-existence-computes} pour
le prouver. Comme collection $\mathcal{I}$ d'objets,
nous prendrons les objets ayant la propriété (I). La propriété
(1) a été démontrée dans le lemme \ref{lemma-right-resolution}. La
propriété (2) est vraie parce que si $s : I \to I'$ est un quasi-isomorphisme de
modules avec la propriété (I), alors $s$ est une équivalence d'homotopie selon
le lemme \ref{lemma-hom-derived}.
\end{proof}

\begin{lemma}
\label{lemma-functoriality-derived-restriction}
Soit $R$ un anneau. Soient $(A, \text{d})$ et $(B, \text{d})$ des
$R$-algèbres différentielles graduées. Soit $f : N \to N'$ un
homomorphisme de $(A, B)$-bimodules différentiels gradués.
Alors $f$ induit un morphisme de foncteurs
$$
- \circ f : R\Hom(N', -) \longrightarrow R\Hom(N, -)
$$
Si $f$ est un quasi-isomorphisme, alors $f \circ -$ est un isomorphisme de
foncteurs.
\end{lemma}

\begin{proof}
Notons $\mathcal{B} = \text{Mod}^{dg}_{(B, \text{d})}$ la
catégorie différentielle graduée des $B$-modules différentiels gradués, voir
l'exemple \ref{example-dgm-dg-cat}. Soit $I$
un $B$-module différentiel gradué avec la propriété (I). Alors $f \circ
- : \Hom_\mathcal{B}(N', I) \to \Hom_\mathcal{B}(N, I)$ est une application de
$A$-modules différentiels gradués. De plus, ceci est fonctoriel par rapport à $I$.
Puisque les foncteurs $ R\Hom(N', -)$ et
$R\Hom(N, -)$ sont calculés en prenant pour
deuxième argument de $\Hom_\mathcal{B}$ des objets possédant la propriété (I)
(lemme \ref{lemma-derived-restriction}), nous obtenons une transformation de foncteurs comme
indiqué.

\medskip\noindent
Supposons que $f$ soit un quasi-isomorphisme. Soit $F_\bullet$ la
filtration donnée sur $I$. Puisque $I = \lim I/F_pI$, nous voyons que
$\Hom_\mathcal{B}(N', I) = \lim \Hom_\mathcal{B}(N', I/F_pI)$ et
$\Hom_\mathcal{B}(N, I) = \lim \Hom_\mathcal{B}(N, I/F_pI)$. Puisque les
applications de transition dans le système $I/F_pI$ sont scindées comme
morphismes de modules gradués, nous voyons que les applications de
transition dans les systèmes $\Hom_\mathcal{B}(N', I/F_pI)$ et $\Hom_\mathcal{B}(N,
I/F_pI)$ sont surjectives. Ainsi, $\Hom_\mathcal{B}(N', I)$, resp. $\Hom_\mathcal{B}(N, I)$
vu comme un complexe de groupes abéliens calcule $R\lim$ du système de
complexes
$\Hom_\mathcal{B}(N', I/F_pI)$, resp. $\Hom_\mathcal{B}(N, I/F_pI)$. Voir plus sur
l'algèbre, lemme \ref{more-algebra-lemma-compute-Rlim}. Ainsi, il suffit de
prouver que chaque morphisme
$$
\Hom_\mathcal{B}(N', I/F_pI) \to \Hom_\mathcal{B}(N, I/F_pI)
$$
est un quasi-isomorphisme. Puisque les surjections $I/F_{p + 1}I \to I/F_pI$ sont
scindées en tant que morphismes de $B$-modules gradués, nous voyons que
$$
0 \to \Hom_\mathcal{B}(N', F_pI/F_{p + 1}I) \to
\Hom_\mathcal{B}(N', I/F_{p + 1}I) \to
\Hom_\mathcal{B}(N', I/F_pI) \to 0
$$
est une suite exacte courte de $A$-modules différentiels gradués. Il
existe une suite analogue pour $N$ et $f$ induit un morphisme
de suites exactes courtes. Ainsi, par récurrence sur $p$ (en commençant par $p = 0$
lorsque $I/F_0I = 0$), nous concluons qu'il suffit de montrer que
l'application
$\Hom_\mathcal{B}(N', F_pI/F_{p + 1}I) \to \Hom_\mathcal{B}(N, F_pI/F_{p + 1}I)$ est un
quasi-isomorphisme. Puisque $F_pI/F_{p + 1}I$ est un produit de décalages de
$A^\vee$, il suffit de prouver que
$\Hom_\mathcal{B}(N', B^\vee[k]) \to \Hom_\mathcal{B}(N, B^\vee[k])$ est un
quasi-isomorphisme. D'après le lemme \ref{lemma-hom-into-shift-dual-free}, il
suffit de montrer que $(N')^\vee \to N^\vee$ est un quasi-isomorphisme. Ceci est
vrai parce que $f$ est un quasi-isomorphisme et $(\ )^\vee$ est un
foncteur exact.
\end{proof}

\begin{lemma}
\label{lemma-derived-restriction-exts}
Soient $(A, \text{d})$ et $(B, \text{d})$ des algèbres différentielles
graduées sur un anneau $R$. Soit $N$ un $(A, B)$-bimodule différentiel gradué.
Alors, pour tout $n \in \mathbf{Z}$, il existe des isomorphismes
$$
H^n(R\Hom(N, M)) = \Ext^n_{D(B, \text{d})}(N, M)
$$
de $R$-modules, fonctoriels en $M$. Ils sont également fonctoriels
en $N$ par rapport à l'opération décrite dans
le lemme \ref{lemma-functoriality-derived-restriction}.
\end{lemma}

\begin{proof}
Dans la démonstration du lemme
\ref{lemma-derived-restriction}, nous avons vu
$$
R\Hom(N, M) =
\Hom_{\text{Mod}^{dg}_{(B, \text{d})}}(N, I)
$$
comme $A$-module différentiel gradué
où $M \to I$ est un quasi-isomorphisme de $M$ vers un $B$-module
différentiel gradué ayant la propriété (I). Par conséquent, ce complexe a les
bons modules de cohomologie selon le lemme \ref{lemma-hom-derived}.
Nous omettons une discussion de la nature fonctorielle de ces
identifications.
\end{proof}

\begin{lemma}
\label{lemma-compute-derived-restriction}
Soit $R$ un anneau. Soient $(A, \text{d})$ et $(B, \text{d})$ des
$R$-algèbres différentielles graduées. Soit $N$ un $(A,
B)$-bimodule différentiel gradué. Si
$\Hom_{D(B, \text{d})}(N, N') = \Hom_{K(\text{Mod}_{(B, \text{d})})}(N, N')$ pour
tout $N' \in K(B, \text{d})$, par exemple si $N$ a la
propriété (P) en tant que $B$-module différentiel gradué, alors
$$
R\Hom(N, M) = \Hom_{\text{Mod}^{dg}_{(B, \text{d})}}(N, M)
$$
de manière fonctorielle en $M$ dans $D(B, \text{d})$.
\end{lemma}

\begin{proof}
Par construction (lemme \ref{lemma-derived-restriction}) pour
trouver $R\Hom(N, M)$ nous choisissons un
quasi-isomorphisme $M \to I$ où $I$ est un $B$-module différentiel gradué avec
la propriété (I) et nous posons
$R\Hom(N, M) = \Hom_{\text{Mod}^{dg}_{(B, \text{d})}}(N, I)$. Par
hypothèse, l'application
$$
\Hom_{\text{Mod}^{dg}_{(B, \text{d})}}(N, M) \longrightarrow
\Hom_{\text{Mod}^{dg}_{(B, \text{d})}}(N, I)
$$
induite par $M \to I$ est un quasi-isomorphisme, voir discussion dans
l'exemple \ref{example-dgm-dg-cat}. Cela prouve le lemme. Si
$N$ a la propriété (P) en tant que $B$-module, alors nous voyons que
l'hypothèse est satisfaite par le lemme \ref{lemma-hom-derived}.
\end{proof}




\section{Variante du Hom dérivé}
\label{section-variant}

\noindent
Soit $\mathcal{A}$ une catégorie abélienne. Considérons la catégorie
différentielle graduée $\text{Comp}^{dg}(\mathcal{A})$ des complexes de $\mathcal{A}$, voir
l'exemple \ref{example-category-complexes}.
Soit $K^\bullet$ un complexe de $\mathcal{A}$. Posons
$$
(E, \text{d}) = \Hom_{\text{Comp}^{dg}(\mathcal{A})}(K^\bullet, K^\bullet)
$$
et considérons le foncteur de catégories différentielles graduées
$$
\text{Comp}^{dg}(\mathcal{A}) \longrightarrow \text{Mod}^{dg}_{(E, \text{d})},
\quad
X^\bullet
\longmapsto
\Hom_{\text{Comp}^{dg}(\mathcal{A})}(K^\bullet, X^\bullet)
$$
du lemme \ref{lemma-construction}.

\begin{lemma}
\label{lemma-existence-of-derived}
Dans la situation ci-dessus, si le foncteur dérivé à droite $R\Hom(K^\bullet,
-)$ de $\Hom(K^\bullet, -) : K(\mathcal{A}) \to D(\textit{Ab})$
est défini partout sur $D(\mathcal{A})$, alors nous obtenons un foncteur exact
canonique
$$
R\Hom(K^\bullet, -) : D(\mathcal{A}) \longrightarrow D(E, \text{d})
$$
de catégories triangulées qui se réduit au foncteur usuel lorsqu'on prend
les complexes associés de groupes abéliens.
\end{lemma}

\begin{proof}
Notons que nous avons un foncteur associé
$K(\mathcal{A}) \to K(\text{Mod}_{(E, \text{d})})$ d'après
le lemme \ref{lemma-construction}.
Nous affirmons que ce foncteur est un foncteur exact de catégories
triangulées. En effet, soit $f : A^\bullet \to B^\bullet$ un morphisme de
complexes de $\mathcal{A}$. Alors un calcul montre que
$$
\Hom_{\text{Comp}^{dg}(\mathcal{A})}(K^\bullet, C(f)^\bullet)
=
C\left(
\Hom_{\text{Comp}^{dg}(\mathcal{A})}(K^\bullet, A^\bullet) \to
\Hom_{\text{Comp}^{dg}(\mathcal{A})}(K^\bullet, B^\bullet)
\right)
$$
où le membre de droite est le cône dans $\text{Mod}_{(E, \text{d})}$
défini plus tôt dans ce chapitre. Cela
montre que notre foncteur est compatible avec les cônes, donc avec
les triangles distingués. Soit $X^\bullet$ un objet de $K(\mathcal{A})$.
Considérons la catégorie des quasi-isomorphismes $s : X^\bullet \to Y^\bullet$. Par
hypothèse, le foncteur $(s :
X^\bullet \to Y^\bullet) \mapsto \Hom_\mathcal{A}(K^\bullet, Y^\bullet)$ est
essentiellement constant lorsqu'il est vu dans $D(\textit{Ab})$. Mais
puisque le foncteur d'oubli $D(E, \text{d}) \to D(\textit{Ab})$ est
compatible avec la prise de cohomologie, la même chose est vraie dans $D(E,
\text{d})$. Cela prouve le lemme.
\end{proof}

\noindent
{\bf Avertissement :} Bien que le lemme soit vrai et puisse être
utile sous cette forme, l'algèbre différentielle $E$ n'est la « bonne » que si
$H^n(E) = \Ext^n_{D(\mathcal{A})}(K^\bullet, K^\bullet)$
pour tout $n \in \mathbf{Z}$.





\section{Produit tensoriel dérivé}
\label{section-base-change}

\noindent
Cette section est analogue à Compléments d'algèbre,
section \ref{more-algebra-section-derived-base-change}.

\medskip\noindent
Soit $R$ un anneau. Soient $(A, \text{d})$ et $(B, \text{d})$ des
algèbres différentielles graduées sur $R$. Soit $N$ un
$(A, B)$-bimodule différentiel gradué. Considérons le foncteur
\begin{equation}
\label{equation-bc}
\text{Mod}_{(A, \text{d})}
\longrightarrow
\text{Mod}_{(B, \text{d})},\quad M
\longmapsto M \otimes_A N
\end{equation}
défini dans la section \ref{section-bimodules-tensor}.

\begin{lemma}
\label{lemma-bc-homotopy}
Le foncteur (\ref{equation-bc}) définit un foncteur
exact de catégories triangulées
$K(\text{Mod}_{(A, \text{d})}) \to K(\text{Mod}_{(B, \text{d})})$.
\end{lemma}

\begin{proof}
Via le lemme \ref{lemma-tensor}
et la remarque
\ref{remark-shift-tensor-no-sign}, cela découle du principe général du
lemme \ref{lemma-functor-between-ABC}.
\end{proof}

\noindent
À ce stade, nous pouvons considérer le diagramme
$$
\xymatrix{
K(\text{Mod}_{(A, \text{d})}) \ar[d] \ar[rr]_{- \otimes_A N} \ar[rrd]_F & &
K(\text{Mod}_{(B, \text{d})}) \ar[d] \\
D(A, \text{d}) \ar@{..>}[rr] & &
D(B, \text{d})
}
$$
La flèche en pointillés que nous allons construire ci-dessous
sera le {\it foncteur dérivé à gauche} de la composition
$F$. ({\it Attention} : le diagramme ne commute
pas.) Plus précisément, dans le cadre
général de Catégories dérivées, section
\ref{derived-section-derived-functors}, nous voulons
calculer le foncteur dérivé à gauche de $F$ par rapport au système multiplicatif de
quasi-isomorphismes dans $K(\text{Mod}_{(A, \text{d})})$.

\begin{lemma}
\label{lemma-derived-bc}
Dans la situation ci-dessus, le foncteur dérivé à gauche de $F$
existe. Nous le
notons $- \otimes_A^\mathbf{L} N : D(A, \text{d}) \to D(B, \text{d})$.
\end{lemma}

\begin{proof}
Nous
utiliserons Catégories dérivées, lemme \ref{derived-lemma-find-existence-computes} pour
le prouver. Comme collection $\mathcal{P}$ d'objets,
nous prendrons les objets ayant la propriété (P). La propriété
(1) a été démontrée dans le lemme \ref{lemma-resolve}. La
propriété (2) est vraie parce que si $s : P \to P'$ est un quasi-isomorphisme de
modules avec la propriété (P), alors $s$ est une équivalence d'homotopie selon
le lemme \ref{lemma-hom-derived}.
\end{proof}

\begin{lemma}
\label{lemma-functoriality-bc}
Soit $R$ un anneau. Soient $(A, \text{d})$ et $(B, \text{d})$ des
$R$-algèbres différentielles graduées. Soit $f : N \to N'$ un
homomorphisme de bimodules différentiels gradués sur $(A, B)$.
Alors $f$ induit un morphisme de foncteurs
$$
1\otimes f :
- \otimes_A^\mathbf{L} N
\longrightarrow -
\otimes_A^\mathbf{L} N'
$$
Si $f$ est un quasi-isomorphisme, alors $1 \otimes f$ est un isomorphisme de
foncteurs.
\end{lemma}

\begin{proof}
Soit $M$ un $A$-module différentiel gradué avec la propriété (P).
Alors $1 \otimes f : M \otimes_A N \to M \otimes_A N'$ est une
application de $B$-modules différentiels gradués. De plus, cela est fonctoriel par
rapport à $M$. Étant donné que les foncteurs $-
\otimes_A^\mathbf{L} N$ et $- \otimes_A^\mathbf{L} N'$ sont calculés par
le produit tensoriel des objets avec la propriété (P)
(lemme \ref{lemma-derived-bc}), nous obtenons une transformation de foncteurs comme
indiqué.

\medskip\noindent
Supposons que $f$ soit un quasi-isomorphisme. Soit $F_\bullet$ la
filtration donnée sur $M$. Observons que $M
\otimes_A N = \colim F_i(M) \otimes_A N$ et $M \otimes_A N'
= \colim F_i(M) \otimes_A N'$. Par conséquent, il
suffit de montrer que $F_n(M) \otimes_A
N \to F_n(M) \otimes_A N'$ est un
quasi-isomorphisme (les limites inductives filtrantes sont exactes, voir Algèbre,
lemme \ref{algebra-lemma-directed-colimit-exact}). Comme les
inclusions $F_n(M) \to F_{n + 1}(M)$ sont scindées en tant que
morphismes de $A$-modules gradués, nous voyons que
$$
0 \to F_n(M) \otimes_A N \to F_{n + 1}(M) \otimes_A N \to
F_{n + 1}(M)/F_n(M) \otimes_A N \to 0
$$
est une suite exacte courte de $B$-modules différentiels gradués. Il
existe une suite analogue pour $N'$ et $f$ induit un morphisme de
suites exactes courtes. Ainsi, par récurrence sur $n$ (en commençant par $n = -1$
lorsque $F_{-1}(M) = 0$), nous concluons qu'il suffit de montrer que
l'application $F_{n + 1}(M)/F_n(M) \otimes_A N \to F_{n + 1}(M)/F_n(M) \otimes_A N'$ est un
quasi-isomorphisme. C'est vrai parce que $F_{n + 1}(M)/F_n(M)$ est une
somme directe de décalages de $A$ et le résultat est vrai pour $A[k]$ puisque
$f : N \to N'$ est un quasi-isomorphisme.
\end{proof}

\begin{lemma}
\label{lemma-compute-bc}
Soit $R$ un anneau. Soient $(A, \text{d})$ et $(B, \text{d})$ des
$R$-algèbres différentielles graduées. Soit $N$ un $(A, B)$-bimodule
différentiel gradué qui possède la propriété (P) en tant que $A$-module différentiel
gradué à gauche. Alors $M \otimes_A^\mathbf{L} N$ est calculé
par $M \otimes_A N$ pour tout $A$-module différentiel gradué $M$.
\end{lemma}

\begin{proof}
Soit $f : M \to M'$ un homomorphisme de $A$-modules différentiels gradués qui
est un quasi-isomorphisme. Nous affirmons que $f \otimes \text{id} : M
\otimes_A N \to M' \otimes_A N$ est un quasi-isomorphisme. Si cela est
vrai, alors par la construction du produit tensoriel dérivé dans la
démonstration du lemme \ref{lemma-derived-bc}, nous obtenons le résultat souhaité. La
construction de l'application $f \otimes \text{id}$ ne dépend que de la
structure de $A$-module différentiel gradué à gauche de $N$. De
plus, nous avons $M \otimes_A N = N \otimes_{A^{opp}} M = N
\otimes_{A^{opp}}^\mathbf{L} M$ parce que $N$ possède la propriété (P) comme
$A^{opp}$-module différentiel gradué. Par conséquent, l'énoncé découle du
lemme \ref{lemma-functoriality-bc}.
\end{proof}

\begin{lemma}
\label{lemma-tensor-hom-adjoint}
Soit $R$ un anneau.
Soient $(A, \text{d})$ et $(B, \text{d})$ des $R$-algèbres différentielles
graduées. Soit $N$ un $(A, B)$-bimodule différentiel gradué.
Alors le foncteur
$$
- \otimes_A^\mathbf{L} N : D(A, \text{d}) \longrightarrow D(B, \text{d})
$$
du lemme \ref{lemma-derived-bc} est un adjoint à gauche du foncteur
$$
R\Hom(N, -) : D(B, \text{d}) \longrightarrow D(A, \text{d})
$$
du lemme \ref{lemma-derived-restriction}.
\end{lemma}

\begin{proof}
Cela découle de Catégories dérivées, lemme
\ref{derived-lemma-pre-derived-adjoint-functors-general} et
du fait que $- \otimes_A N$ et
$\Hom_{\text{Mod}^{dg}_{(B, \text{d})}}(N, -)$ sont adjoints d'après le
lemme \ref{lemma-tensor-hom-adjunction}.
\end{proof}

\begin{example}
\label{example-map-hom-tensor}
Soit $R$ un anneau. Soit $(A, \text{d}) \to (B, \text{d})$ un
homomorphisme d'algèbres différentielles graduées sur $R$. Alors nous pouvons
considérer $B$ comme un $(A, B)$-bimodule différentiel gradué et nous obtenons un foncteur
$$
- \otimes_A B : D(A, \text{d}) \longrightarrow D(B, \text{d})
$$
D'après le lemme \ref{lemma-tensor-hom-adjoint}, le foncteur adjoint à
gauche de celui-ci est le foncteur $R\Hom(B, -)$. Pour un $B$-module
différentiel gradué, notons $N_A$ le $A$-module différentiel gradué obtenu à partir
de $N$ par restriction via $A \to B$. Alors nous avons clairement
un isomorphisme canonique
$$
\Hom_{\text{Mod}^{dg}_{(B, \text{d})}}(B, N) \longrightarrow N_A,\quad
f \longmapsto f(1)
$$
fonctoriel par rapport au $B$-module $N$. Ainsi, nous voyons
que $R\Hom(B, -)$ est le foncteur de restriction et nous obtenons
$$
\Hom_{D(A, \text{d})}(M, N_A) =
\Hom_{D(B, \text{d})}(M \otimes^\mathbf{L}_A B, N)
$$
bifonctoriellement en $M$ et $N$ exactement comme dans le cas des anneaux
commutatifs. Enfin, observons que la restriction est également un foncteur
tensoriel, puisque $N_A = N \otimes_B {}_BB_A = N \otimes_B^\mathbf{L}
{}_BB_A$ où ${}_BB_A$ est $B$ vu comme un $(B, A)$-bimodule différentiel gradué.
\end{example}

\begin{lemma}
\label{lemma-tensor-with-compact-fully-faithful}
Avec les notations et sous les hypothèses du lemme \ref{lemma-tensor-hom-adjoint}, faisons les hypothèses
suivantes :
\begin{enumerate}
\item $N$ définit un objet compact de $D(B, \text{d})$, et
\item l'application $H^k(A) \to \Hom_{D(B, \text{d})}(N, N[k])$ est
un isomorphisme pour tout $k \in \mathbf{Z}$.
\end{enumerate}
Alors le foncteur $-\otimes_A^\mathbf{L} N$ est pleinement fidèle.
\end{lemma}

\begin{proof}
Notre foncteur a un adjoint à gauche donné
par $R\Hom(N, -)$ d'après le lemme
\ref{lemma-tensor-hom-adjoint}. D'après Catégories, lemme \ref{categories-lemma-adjoint-fully-faithful},
il suffit de montrer que pour un $A$-module différentiel gradué $M$,
l'application
$$
M \longrightarrow R\Hom(N, M \otimes_A^\mathbf{L} N)
$$
est un isomorphisme dans $D(A, \text{d})$. Pour cela, il suffit de montrer que
$$
H^n(M) \longrightarrow
\text{Ext}^n_{D(B, \text{d})}(N, M \otimes_A^\mathbf{L} N)
$$
est un isomorphisme, voir le lemme \ref{lemma-derived-restriction-exts}.
Puisque $N$ est un objet compact, le membre de droite commute avec les
sommes directes. Ainsi, d'après la remarque \ref{remark-P-resolution}, il
suffit de prouver que cette application est un isomorphisme pour $M = A[k]$.
Puisque $(A[k] \otimes_A^\mathbf{L} N) = N[k]$ d'après la
remarque \ref{remark-shift-tensor-no-sign},
l'hypothèse (2) sur $N$ affirme précisément le résultat dans ces cas.
\end{proof}

\begin{lemma}
\label{lemma-base-change-K-flat}
Soit $R \to R'$ un homomorphisme d'anneaux. Soit $(A, \text{d})$ une $R$-algèbre
différentielle graduée. Soit $(A', \text{d})$ le changement de base,
c'est-à-dire $A' = A \otimes_R R'$. Si $A$ est K-plat en tant que complexe de $R$-modules,
alors
\begin{enumerate}
\item $- \otimes_A^\mathbf{L} A' : D(A, \text{d}) \to D(A', \text{d})$
est égal au foncteur dérivé à droite de
$$
K(A, \text{d}) \longrightarrow K(A', \text{d}),\quad
M \longmapsto M \otimes_R R'
$$
\item le diagramme
$$
\xymatrix{
D(A, \text{d}) \ar[r]_{- \otimes_A^\mathbf{L} A'} \ar[d]_{restriction} &
D(A', \text{d}) \ar[d]^{restriction} \\
D(R) \ar[r]^{- \otimes_R^\mathbf{L} R'} & D(R')
}
$$
commute, et
\item si $M$ est K-plat en tant que complexe de $R$-modules,
alors le $A'$-module différentiel gradué $M \otimes_R R'$
représente $M \otimes_A^\mathbf{L} A'$.
\end{enumerate}
\end{lemma}

\begin{proof}
Pour tout $A$-module différentiel gradué $M$, il existe une application canonique
$$
c_M : M \otimes_R R' \longrightarrow M \otimes_A A'
$$
Soit $P$ un $A$-module différentiel gradué avec la propriété (P).
Nous affirmons que $c_P$ est un isomorphisme et que $P$ est K-plat en
tant que complexe de $R$-modules. Cela prouvera tous les résultats
énoncés dans le lemme par des arguments formels utilisant la
définition du produit tensoriel dérivé dans le lemme \ref{lemma-derived-bc}
et Compléments d'algèbre, section \ref{more-algebra-section-derived-tensor-product}.

\medskip\noindent
Soit $F_\bullet$ la filtration sur $P$ montrant que $P$ a la propriété (P). Notons que
$c_A$ est un isomorphisme et que $A$ est K-plat en tant que complexe de
$R$-modules par hypothèse. Par conséquent, il en va de même pour les sommes
directes de décalages de $A$ (vous pouvez utiliser
Compléments d'algèbre, lemme \ref{more-algebra-lemma-colimit-K-flat} pour
traiter les sommes directes si vous le souhaitez).
Par conséquent, cela vaut pour les complexes $F_{p + 1}P/F_pP$.
Comme les suites exactes courtes
$$
0 \to F_pP \to F_{p + 1}P \to F_{p + 1}P/F_pP \to 0
$$
sont scindées comme suites de modules gradués, nous pouvons montrer
par récurrence que $c_{F_pP}$ est un isomorphisme pour tout $p$
et que $F_pP$ est K-plat en tant que complexe de $R$-modules (voir
Compléments d'algèbre, lemme \ref{more-algebra-lemma-K-flat-two-out-of-three}).
Enfin, en utilisant que $P = \colim F_pP$, nous concluons que
$c_P$ est un isomorphisme et que $P$ est K-plat en tant
que complexe de $R$-modules (voir
Compléments d'algèbre, lemme \ref{more-algebra-lemma-colimit-K-flat}).
\end{proof}

\begin{lemma}
\label{lemma-RHom-is-tensor}
Soit $R$ un anneau.
Soient $(A, \text{d})$ et $(B, \text{d})$ des $R$-algèbres différentielles graduées. Soit $T$ un
$(A, B)$-bimodule différentiel gradué. Faisons les hypothèses
suivantes :
\begin{enumerate}
\item $T$ définit un objet compact de $D(B, \text{d})$, et
\item $S = \Hom_{\text{Mod}^{dg}_{(B, \text{d})}}(T, B)$
représente $R\Hom(T, B)$ dans $D(A, \text{d})$.
\end{enumerate}
Alors $S$ a une structure de $(B, A)$-bimodule différentiel gradué
et il existe un isomorphisme
$$
N \otimes_B^\mathbf{L} S \longrightarrow R\Hom(T, N)
$$
fonctorielle en $N$ dans $D(B, \text{d})$.
\end{lemma}

\begin{proof}
Notons $\mathcal{B} = \text{Mod}^{dg}_{(B, \text{d})}$. La
structure de $A$-module à droite de $S$ provient de l'application
$A \to \Hom_\mathcal{B}(T, T)$ et de la composition
$\Hom_\mathcal{B}(T, B) \otimes \Hom_\mathcal{B}(T, T) \to
\Hom_\mathcal{B}(T, B)$ définie dans l'exemple \ref{example-dgm-dg-cat}. En utilisant cette
multiplication une seconde fois, on obtient une application
$$
c_N :
N \otimes_B S = \Hom_\mathcal{B}(B, N) \otimes_B \Hom_\mathcal{B}(T, B)
\longrightarrow
\Hom_\mathcal{B}(T, N)
$$
fonctoriel en $N$. Étant donné $N$, nous pouvons choisir des
quasi-isomorphismes $P \to N \to I$ où $P$, resp.\ $I$ est un $B$-module différentiel gradué avec
la propriété (P), resp.\ (I). Ensuite, en utilisant $c_N$, nous obtenons
une application $P \otimes_B S \to \Hom_\mathcal{B}(T, I)$ entre les
objets représentant $S \otimes_B^\mathbf{L} N$ et $R\Hom(T, N)$. De toute
évidence, cela définit une transformation de foncteurs $c$ comme dans le lemme.

\medskip\noindent
Pour prouver que $c$ est un isomorphisme de foncteurs, nous
pouvons supposer que $N$ est un $B$-module différentiel gradué qui
possède la propriété (P). Puisque $T$ définit un objet compact
dans $D(B, \text{d})$ et puisque les deux côtés de la flèche définissent
des foncteurs exacts de catégories triangulées, nous nous ramenons, à
l'aide du lemme
\ref{lemma-property-P-sequence}, au cas où $N$ possède une filtration finie dont les
quotients successifs sont des sommes directes de $B[k]$. En utilisant à
nouveau que les deux côtés de la flèche sont des foncteurs exacts de
catégories triangulées et la compacité de $T$, nous nous
ramenons au cas $N = B[k]$. L'hypothèse (2) est exactement
l'hypothèse que $c$ est un isomorphisme dans ce cas.
\end{proof}






\section{Composition des produits tensoriels dérivés}
\label{section-compose-tensor-functors}

\noindent
Nous encourageons le lecteur à passer cette section.

\medskip\noindent
Soit $R$ un anneau. Soient $(A, \text{d})$, $(B, \text{d})$, et
$(C, \text{d})$ des $R$-algèbres différentielles graduées.
Soit $N$ un $(A, B)$-bimodule différentiel gradué. Soit $N'$
un $(B, C)$-module différentiel gradué. Nous notons
$N_B$ le bimodule $N$ considéré comme un $B$-module différentiel gradué (en
oubliant la structure de module sur $A$). Il existe une application canonique
\begin{equation}
\label{equation-plain-versus-derived}
N_B \otimes_B^\mathbf{L} N'
\longrightarrow
(N \otimes_B N')_C
\end{equation}
dans $D(C, \text{d})$. Ici, $(N \otimes_B N')_C$ désigne le $(A,
C)$-bimodule $N \otimes_B N'$ considéré comme un $C$-module
différentiel gradué. En effet, cette application découle du
fait qu'il existe toujours un morphisme du produit tensoriel dérivé vers le produit
tensoriel ordinaire (comme il s'agit d'un foncteur dérivé à gauche).

\begin{lemma}
\label{lemma-compose-tensor-functors-general}
Soit $R$ un anneau. Soient $(A, \text{d})$, $(B, \text{d})$ et
$(C, \text{d})$ des algèbres différentielles graduées sur
$R$. Soit $N$ un $(A, B)$-bimodule différentiel gradué.
Soit $N'$ un module différentiel gradué $(B, C)$.
Supposons que (\ref{equation-plain-versus-derived}) soit un isomorphisme.
Alors la composition
$$
\xymatrix{
D(A, \text{d}) \ar[rr]^{- \otimes_A^\mathbf{L} N} & &
D(B, \text{d}) \ar[rr]^{- \otimes_B^\mathbf{L} N'} & &
D(C, \text{d})
}
$$
est isomorphe à $- \otimes_A^\mathbf{L} N''$ avec
$N'' = N \otimes_B N'$ considéré comme $(A, C)$-bimodule.
\end{lemma}

\begin{proof}
Définissons une transformation de foncteurs
$$
(- \otimes_A^\mathbf{L} N) \otimes_B^\mathbf{L} N'
\longrightarrow
- \otimes_A^\mathbf{L} N''
$$
Pour ce faire,
soit $M$ un $A$-module différentiel gradué avec la propriété
(P). Selon la construction du foncteur $- \otimes_A^\mathbf{L} N''$ de la
démonstration du lemme \ref{lemma-derived-bc}, le produit tensoriel
ordinaire $M \otimes_A N''$ représente $M \otimes_A^\mathbf{L} N''$ dans $D(C,
\text{d})$. Ensuite, nous écrivons
$$
M \otimes_A N'' =
M \otimes_A (N \otimes_B N') =
(M \otimes_A N) \otimes_B N'
$$
Le module $M \otimes_A N$ représente $M \otimes_A^\mathbf{L} N$ dans
$D(B, \text{d})$. Choisissons un quasi-isomorphisme $Q \to M \otimes_A N$ où $Q$
est un $B$-module différentiel gradué avec la propriété (P). Alors $Q
\otimes_B N'$ représente $(M
\otimes_A^\mathbf{L} N) \otimes_B^\mathbf{L} N'$ dans $D(C, \text{d})$.
Ainsi, nous pouvons
définir notre application via
$$
(M \otimes_A^\mathbf{L} N) \otimes_B^\mathbf{L} N' =
Q \otimes_B N' \to
M \otimes_A N \otimes_B N' = M
\otimes_A^\mathbf{L} N''
$$
La construction de cette application est fonctorielle en $M$ et compatible
avec les triangles distingués et les sommes directes ; nous omettons les
détails. Considérons la propriété $T$ des objets $M$ de $D(A, \text{d})$
exprimant que cette application est un isomorphisme. Alors
\begin{enumerate}
\item si $T$ est vraie pour $M_i$, alors $T$ est vraie pour $\bigoplus M_i$,
\item si $T$ est vraie pour $2$ des $3$ objets d'un triangle
distingué, alors elle est vraie pour le troisième, et
\item $T$ vaut pour $A[k]$ car ici nous obtenons un décalage de
l'application (\ref{equation-plain-versus-derived}) que nous
avons supposée être un isomorphisme.
\end{enumerate}
Ainsi, d'après la remarque \ref{remark-P-resolution}, la propriété $T$
est toujours vraie et la démonstration est complète.
\end{proof}

\noindent
Soit $R$ un anneau. Soient $(A, \text{d})$, $(B, \text{d})$ et
$(C, \text{d})$ des $R$-algèbres différentielles graduées.
Nous notons temporairement $(A \otimes_R B)_B$ l'algèbre
différentielle graduée $A \otimes_R B$ vue comme un module différentiel gradué à
droite sur $B$, et ${}_B(B \otimes_R C)_C$ l'algèbre différentielle
graduée $B \otimes_R C$ vue comme un bimodule différentiel gradué sur
$(B, C)$. Alors il existe une application canonique
\begin{equation}
\label{equation-plain-versus-derived-algebras}
(A \otimes_R B)_B \otimes_B^\mathbf{L} {}_B(B \otimes_R C)_C
\longrightarrow
(A \otimes_R B \otimes_R C)_C
\end{equation}
dans $D(C, \text{d})$ où $(A \otimes_R B \otimes_R C)_C$
désigne la $R$-algèbre
différentielle graduée $A \otimes_R B \otimes_R C$ considérée comme un
$C$-module différentiel gradué (à droite). En effet, cette
application provient de l'identification
$$
(A \otimes_R B)_B \otimes_B {}_B(B \otimes_R C)_C =
(A \otimes_R B \otimes_R C)_C
$$
et du fait que le produit tensoriel dérivé s'envoie toujours dans le
produit tensoriel ordinaire (puisque c'est un foncteur dérivé à gauche).

\begin{lemma}
\label{lemma-compose-tensor-functors-general-algebra}
Soit $R$ un anneau. Soient $(A, \text{d})$, $(B, \text{d})$ et
$(C, \text{d})$ des $R$-algèbres différentielles graduées.
Supposons que (\ref{equation-plain-versus-derived-algebras}) soit un
isomorphisme. Soit $N$ un $(A, B)$-bimodule différentiel gradué.
Soit $N'$ un $(B, C)$-bimodule différentiel gradué.
Alors la composition
$$
\xymatrix{
D(A, \text{d}) \ar[rr]^{- \otimes_A^\mathbf{L} N} & &
D(B, \text{d}) \ar[rr]^{- \otimes_B^\mathbf{L} N'} & &
D(C, \text{d})
}
$$
est isomorphe à $- \otimes_A^\mathbf{L} N''$ pour un $(A, C)$-bimodule
différentiel gradué $N''$ décrit dans la démonstration.
\end{lemma}

\begin{proof}
Selon le lemme \ref{lemma-functoriality-bc}, nous pouvons remplacer $N$ et $N'$ par
des bimodules quasi-isomorphes. Ainsi, nous pouvons supposer que $N$, resp.\ $N'$,
possède la propriété (P) en tant que bimodule différentiel gradué $(A, B)$,
resp.\ $(B, C)$,
voir le lemme \ref{lemma-bimodule-resolve}. Nous affirmons que le
lemme est valable avec le $(A, C)$-bimodule $N'' = N \otimes_B
N'$. Pour le prouver, il suffit de montrer que
$$
N_B \otimes_B^\mathbf{L} N' \longrightarrow (N \otimes_B N')_C
$$
est un isomorphisme dans $D(C, \text{d})$,
voir le lemme \ref{lemma-compose-tensor-functors-general}.

\medskip\noindent
Soit $F_\bullet$ la filtration sur $N$ comme dans la propriété (P) pour les bimodules.
D'après le lemme
\ref{lemma-bimodule-property-P-sequence}, il existe une suite exacte courte
$$
0 \to
\bigoplus\nolimits F_iN \to
\bigoplus\nolimits F_iN \to N \to 0
$$
de $(A, B)$-bimodules différentiels gradués, qui est scindée comme suite
de $(A, B)$-bimodules gradués. A fortiori, il s'agit d'une suite exacte
courte admissible de $B$-modules différentiels gradués et cela produit un triangle
distingué
$$
\bigoplus\nolimits F_iN_B \to
\bigoplus\nolimits F_iN_B \to N_B \to
\bigoplus\nolimits F_iN_B[1]
$$
dans $D(B, \text{d})$. En utilisant que $- \otimes_B^\mathbf{L} N'$
est un foncteur exact de catégories triangulées et commute avec les
sommes directes et en utilisant que $- \otimes_B N'$ transforme les
suites exactes admissibles en suites exactes admissibles et commute
avec les sommes directes, nous nous ramenons à prouver
que
$$
(F_pN)_B \otimes_B^\mathbf{L} N' \longrightarrow (F_pN)_B \otimes_B N'
$$
est un quasi-isomorphisme pour tout $p$. En répétant
l'argument avec les suites exactes courtes de $(A, B)$-bimodules
$$
0 \to F_pN \to F_{p + 1}N \to F_{p + 1}N/F_pN \to 0
$$
qui sont scindées comme suites de $(A, B)$-bimodules
gradués, nous nous ramenons à montrer la même assertion pour $F_{p +
1}N/F_pN$. Comme ces modules sont des sommes directes de décalages de $(A \otimes_R B)_B$,
nous nous ramenons à montrer que
$$
(A \otimes_R B)_B \otimes_B^\mathbf{L} N'
\longrightarrow
(A \otimes_R B)_B \otimes_B N'
$$
est un quasi-isomorphisme.

\medskip\noindent
Choisissons une filtration $F_\bullet$ sur $N'$ comme dans la propriété (P) pour
les bimodules. Choisissons un quasi-isomorphisme $P \to (A
\otimes_R B)_B$ de $B$-modules différentiels gradués où $P$ a la propriété
(P). Nous devons montrer
que $P \otimes_B N' \to (A \otimes_R B)_B \otimes_B N'$ est un
quasi-isomorphisme parce que $P \otimes_B N'$ représente $(A
\otimes_R B)_B \otimes_B^\mathbf{L} N'$ dans $D(C, \text{d})$ par la
construction dans le lemme \ref{lemma-derived-bc}. Puisque $N' =
\colim F_pN'$, nous constatons qu'il
suffit de montrer que $P \otimes_B
F_pN' \to (A \otimes_R B)_B \otimes_B F_pN'$ est un
quasi-isomorphisme. En utilisant les suites exactes courtes $0 \to F_pN' \to
F_{p + 1}N' \to F_{p + 1}N'/F_pN' \to 0$ qui sont scindées comme
suites de $(B, C)$-bimodules gradués, nous nous ramenons à montrer que $P
\otimes_B F_{p + 1}N'/F_pN' \to (A
\otimes_R B)_B \otimes_B F_{p + 1}N'/F_pN'$ est un
quasi-isomorphisme pour tout $p$. Ensuite, en
utilisant enfin que $F_{p + 1}N'/F_pN'$ est une somme
directe de décalages de ${}_B(B \otimes_R C)_C$, nous
concluons qu'il suffit de montrer que
$$
P \otimes_B {}_B(B \otimes_R C)_C \to
(A \otimes_R B)_B \otimes_B {}_B(B \otimes_R C)_C
$$
est un quasi-isomorphisme. Puisque $P \to (A \otimes_R B)_B$
est une résolution par un module satisfaisant à la propriété
(P), cette application de $C$-modules différentiels
gradués représente le morphisme (\ref{equation-plain-versus-derived-algebras}) dans
$D(C, \text{d})$ et la démonstration est terminée.
\end{proof}

\begin{lemma}
\label{lemma-compose-tensor-functors}
Soit $R$ un anneau. Soient $(A, \text{d})$, $(B, \text{d})$ et
$(C, \text{d})$ des $R$-algèbres différentielles
graduées. Si $C$ est K-plat comme complexe de $R$-modules,
alors (\ref{equation-plain-versus-derived-algebras})
est un isomorphisme et la conclusion du
lemme \ref{lemma-compose-tensor-functors-general-algebra} est valable.
\end{lemma}

\begin{proof}
Choisissons un quasi-isomorphisme $P \to (A \otimes_R B)_B$ de $B$-modules
différentiels gradués, où $P$ a la propriété (P). Ensuite, nous devons montrer
que
$$
P \otimes_B (B \otimes_R C) \longrightarrow
(A \otimes_R B) \otimes_B (B \otimes_R C)
$$
est un quasi-isomorphisme. De manière équivalente, il s'agit du morphisme
$$
P \otimes_R C \longrightarrow
A \otimes_R B \otimes_R C
$$
C'est un quasi-isomorphisme si $C$ est K-plat en tant que complexe de $R$-modules
selon Compléments d'algèbre, lemme \ref{more-algebra-lemma-K-flat-quasi-isomorphism}.
\end{proof}





\section{Variante du produit tensoriel dérivé}
\label{section-variant-base-change}

\noindent
Soit $(\mathcal{C}, \mathcal{O})$ un site annelé. Alors nous avons les foncteurs
$$
\text{Comp}(\mathcal{O}) \to K(\mathcal{O}) \to D(\mathcal{O})
$$
et comme nous l'avons vu ci-dessus, nous avons un enrichissement
différentiel gradué $\text{Comp}^{dg}(\mathcal{O})$. Plus précisément, il s'agit
de la catégorie différentielle graduée de l'exemple
\ref{example-category-complexes} associée à la catégorie abélienne
$\textit{Mod}(\mathcal{O})$. Soit $K^\bullet$ un complexe de $\mathcal{O}$-modules, en d'autres
termes, un objet de $\text{Comp}^{dg}(\mathcal{O})$. Posons
$$
(E, \text{d}) =
\Hom_{\text{Comp}^{dg}(\mathcal{O})}(K^\bullet, K^\bullet)
$$
Ceci est une $\mathbf{Z}$-algèbre différentielle graduée. Nous affirmons qu'il
existe un analogue du changement de base dérivé dans cette situation.

\begin{lemma}
\label{lemma-tensor-with-complex}
Dans la situation ci-dessus, il y a un foncteur
$$
- \otimes_E K^\bullet :
\text{Mod}^{dg}_{(E, \text{d})}
\longrightarrow
\text{Comp}^{dg}(\mathcal{O})
$$
de catégories différentielles graduées. Ce foncteur envoie $E$ sur $K^\bullet$ et
commute avec les sommes directes.
\end{lemma}

\begin{proof}
Soit $M$ un $E$-module différentiel gradué. Pour chaque objet $U$ de
$\mathcal{C}$, le complexe $K^\bullet(U)$ est un module différentiel gradué à
gauche sur $E$, ainsi qu'un module à droite sur $\mathcal{O}(U)$. Les
actions commutent, donc nous avons un bimodule.
Ainsi, d'après les constructions des
sections \ref{section-tensor-product} et \ref{section-bimodules}, nous
pouvons former le produit tensoriel
$$
M \otimes_E K^\bullet(U)
$$
qui est un $\mathcal{O}(U)$-module différentiel gradué, c'est-à-dire un complexe
de $\mathcal{O}(U)$-modules. Cette construction est fonctorielle par rapport à $U$,
nous pouvons donc faisceautiser pour obtenir un complexe de $\mathcal{O}$-modules
que nous notons
$$
M \otimes_E K^\bullet
$$
De plus, pour chaque $U$, la construction détermine un foncteur
$\text{Mod}^{dg}_{(E, \text{d})} \to \text{Comp}^{dg}(\mathcal{O}(U))$ de
catégories différentielles graduées selon le lemme \ref{lemma-tensor}.
Il est donc clair que nous obtenons un foncteur tel qu'énoncé dans le lemme.
\end{proof}

\begin{lemma}
\label{lemma-tensor-with-complex-homotopy}
Le foncteur du lemme \ref{lemma-tensor-with-complex} définit un foncteur exact
de catégories triangulées
$K(\text{Mod}_{(E, \text{d})}) \to K(\mathcal{O})$.
\end{lemma}

\begin{proof}
Le foncteur induit un foncteur entre les catégories d'homotopie
selon le lemme
\ref{lemma-functorial}. Nous devons montrer que $- \otimes_E K^\bullet$ transforme les triangles
distingués en triangles distingués. Supposons
que $0 \to K \to L \to M \to 0$ soit une suite exacte courte admissible de
$E$-modules différentiels gradués. Soit $s : M \to L$ un homomorphisme de
$E$-modules gradués qui est l'inverse à gauche de $L \to M$. Alors $s$ définit une
application $M \otimes_E K^\bullet \to L \otimes_E K^\bullet$ de
$\mathcal{O}$-modules gradués (c'est-à-dire respectant la structure de $\mathcal{O}$-module et
la graduation, mais pas les différentielles) qui est
l'inverse à gauche de $L \otimes_E K^\bullet \to M \otimes_E K^\bullet$. Ainsi nous
voyons que
$$
0 \to K \otimes_E K^\bullet \to L \otimes_E K^\bullet \to
M \otimes_E K^\bullet \to 0
$$
est une suite exacte courte scindée terme à terme de complexes,
c'est-à-dire qu'elle définit un triangle distingué dans $K(\mathcal{O})$.
\end{proof}

\begin{lemma}
\label{lemma-tensor-with-complex-derived}
Le foncteur $K(\text{Mod}_{(E, \text{d})}) \to K(\mathcal{O})$ du
lemme \ref{lemma-tensor-with-complex-homotopy} a une version dérivée à
gauche définie sur l'ensemble de $D(E, \text{d})$. Nous la
notons $- \otimes_E^\mathbf{L} K^\bullet : D(E, \text{d}) \to D(\mathcal{O})$.
\end{lemma}

\begin{proof}
Nous
utiliserons Catégories dérivées, lemme \ref{derived-lemma-find-existence-computes} pour
le prouver. Comme collection $\mathcal{P}$ d'objets,
nous prendrons les objets avec la propriété (P). La propriété (1)
a été démontrée dans le lemme \ref{lemma-resolve}. La
propriété (2) est vraie parce que si $s : P \to P'$ est un quasi-isomorphisme de
modules avec la propriété (P), alors $s$ est une équivalence d'homotopie selon
le lemme \ref{lemma-hom-derived}.
\end{proof}

\begin{lemma}
\label{lemma-upgrade-tensor-with-complex-derived}
Soit $R$ un anneau. Soit $\mathcal{C}$ un site. Soit $\mathcal{O}$ un
faisceau de $R$-algèbres commutatives. Soit $K^\bullet$ un
complexe de $\mathcal{O}$-modules. Le
foncteur du lemme
\ref{lemma-tensor-with-complex-derived} possède la propriété suivante :
Pour chaque $M$, $N$ dans $D(E, \text{d})$, il existe une flèche
canonique
$$
R\Hom(M, N)
\longrightarrow
R\Hom_\mathcal{O}(M \otimes_E^\mathbf{L} K^\bullet, N
\otimes_E^\mathbf{L} K^\bullet)
$$
dans $D(R)$ qui sur les modules de cohomologie donne les flèches
$$
\Ext^n_{D(E, \text{d})}(M, N) \to
\Ext^n_{D(\mathcal{O})}
(M \otimes_E^\mathbf{L} K^\bullet, N \otimes_E^\mathbf{L} K^\bullet)
$$
induites par le foncteur $- \otimes_E^\mathbf{L} K^\bullet$.
\end{lemma}

\begin{proof}
Le membre de droite est le Hom dérivé global introduit dans Cohomologie
sur les sites, section \ref{sites-cohomology-section-global-RHom} qui a les bons
modules de cohomologie. Pour le membre de
gauche, nous considérons $M$ comme un $(R, A)$-bimodule et nous avons le
$\Hom$ dérivé introduit dans la section \ref{section-restriction}, qui possède
également les modules de cohomologie corrects. Pour
prouver le lemme, nous pouvons supposer que $M$ et $N$ sont des
$E$-modules différentiels gradués avec la propriété (P) ; cela ne change pas le
membre de gauche selon
le lemme \ref{lemma-functoriality-derived-restriction}.
D'après le lemme
\ref{lemma-compute-derived-restriction}, cela signifie que le membre de gauche devient
$\Hom_{\text{Mod}^{dg}_{(B, \text{d})}}(M, N)$. Dans les
lemmes \ref{lemma-tensor-with-complex},
\ref{lemma-tensor-with-complex-homotopy}, et
\ref{lemma-tensor-with-complex-derived} nous avons
construit un foncteur
$$
- \otimes_E K^\bullet :
\text{Mod}^{dg}_{(E, \text{d})}
\longrightarrow
\text{Comp}^{dg}(\mathcal{O})
$$
de catégories différentielles graduées
et nous avons montré que $- \otimes_E^\mathbf{L} K^\bullet$ est calculé en évaluant
ce foncteur sur des $E$-modules
différentiels gradués ayant la propriété (P). Par conséquent,
nous obtenons un morphisme de complexes de $R$-modules
$$
\Hom_{\text{Mod}^{dg}_{(B, \text{d})}}(M, N)
\longrightarrow
\Hom_{\text{Comp}^{dg}(\mathcal{O})}
(M \otimes_E K^\bullet, N \otimes_E K^\bullet)
$$
Pour tous les complexes de $\mathcal{O}$-modules
$\mathcal{F}^\bullet$ et $\mathcal{G}^\bullet$, il existe une
application canonique
$$
\Hom_{\text{Comp}^{dg}(\mathcal{O})}
(\mathcal{F}^\bullet, \mathcal{G}^\bullet) =
\Gamma(\mathcal{C},
\SheafHom^\bullet(\mathcal{F}^\bullet, \mathcal{G}^\bullet))
\longrightarrow
R\Hom_\mathcal{O}(\mathcal{F}^\bullet, \mathcal{G}^\bullet).
$$
En combinant ces
applications, nous obtenons l'application désirée du lemme.
\end{proof}

\begin{lemma}
\label{lemma-tensor-with-complex-hom-adjoint}
Soit $(\mathcal{C}, \mathcal{O})$ un site annelé.
Soit $K^\bullet$ un complexe de $\mathcal{O}$-modules.
Alors le foncteur
$$
- \otimes_E^\mathbf{L} K^\bullet :
D(E, \text{d})
\longrightarrow
D(\mathcal{O})
$$
du lemme \ref{lemma-tensor-with-complex-derived} est un adjoint à
gauche du foncteur
$$
R\Hom(K^\bullet, -) : D(\mathcal{O}) \longrightarrow D(E, \text{d})
$$
du lemme \ref{lemma-existence-of-derived}.
\end{lemma}

\begin{proof}
L'énoncé signifie que nous avons
$$
\Hom_{D(E, \text{d})}(M, R\Hom(K^\bullet, L^\bullet)) =
\Hom_{D(\mathcal{O})}(M \otimes^\mathbf{L}_E K^\bullet, L^\bullet)
$$
bifonctoriellement en $M$ et $L^\bullet$. Pour voir cela, nous pouvons
remplacer $M$ par un $E$-module différentiel gradué $P$ ayant la
propriété (P). Nous pouvons également remplacer $L^\bullet$ par un
complexe K-injectif de $\mathcal{O}$-modules $I^\bullet$.
Le calcul des foncteurs dérivés donné dans les lemmes mentionnés dans l'énoncé,
combiné avec le lemme \ref{lemma-hom-derived}, traduit ce qui précède en
$$
\Hom_{K(\text{Mod}_{(E, \text{d})})}
(P, \Hom_\mathcal{B}(K^\bullet, I^\bullet)) =
\Hom_{K(\mathcal{O})}(P \otimes_E K^\bullet, I^\bullet)
$$
où $\mathcal{B} = \text{Comp}^{dg}(\mathcal{O})$. Il existe
une application d'évaluation de droite à gauche fonctorielle en
$P$ et $I^\bullet$ (détails omis). Choisissons
une filtration $F_\bullet$ sur $P$ comme dans la définition de la propriété (P). D'après le
lemme \ref{lemma-property-P-sequence} et le fait que les deux
membres de l'égalité sont des foncteurs homologiques en $P$ sur
$K(\text{Mod}_{(E, \text{d})})$, nous nous
ramenons au cas où $P$ est remplacé par le $E$-module
différentiel gradué $\bigoplus F_iP$. Puisque les deux membres
transforment les sommes directes dans la variable $P$ en
produits directs, nous nous ramenons au cas où $P$ est l'un des $E$-modules
différentiels gradués $F_iP$. Puisque chaque $F_iP$
possède une filtration finie (donnée par des monomorphismes admissibles)
dont les quotients successifs sont des $E$-modules projectifs gradués, nous
nous ramenons au cas où $P$ est un $E$-module projectif gradué. Dans ce cas,
nous avons clairement
$$
\Hom_{\text{Mod}^{dg}_{(E, \text{d})}}
(P, \Hom_\mathcal{B}(K^\bullet, I^\bullet)) =
\Hom_{\text{Comp}^{dg}(\mathcal{O})}(P \otimes_E K^\bullet, I^\bullet)
$$
comme $\mathbf{Z}$-modules gradués (car cette affirmation se réduit au cas $P =
E[k]$ où c'est évident). Comme l'isomorphisme est compatible avec les différentielles,
nous obtenons la conclusion.
\end{proof}

\begin{lemma}
\label{lemma-fully-faithful-in-compact-case}
Soit $(\mathcal{C}, \mathcal{O})$ un site annelé. Soit
$K^\bullet$ un complexe de $\mathcal{O}$-modules. Faisons les hypothèses
suivantes :
\begin{enumerate}
\item $K^\bullet$ représente un objet compact de $D(\mathcal{O})$, et
\item $E = \Hom_{\text{Comp}^{dg}(\mathcal{O})}(K^\bullet, K^\bullet)$
calcule les groupes Ext de $K^\bullet$ dans $D(\mathcal{O})$.
\end{enumerate}
Alors le foncteur
$$
- \otimes_E^\mathbf{L} K^\bullet :
D(E, \text{d})
\longrightarrow
D(\mathcal{O})
$$
du lemme \ref{lemma-tensor-with-complex-derived} est pleinement fidèle.
\end{lemma}

\begin{proof}
Comme notre foncteur a un adjoint à gauche donné par
$R\Hom(K^\bullet, -)$ selon le lemme \ref{lemma-tensor-with-complex-hom-adjoint},
il suffit de montrer pour un $E$-module différentiel gradué $M$ que l'application
$$
H^0(M) \longrightarrow
\Hom_{D(\mathcal{O})}(K^\bullet, M \otimes_E^\mathbf{L} K^\bullet)
$$
est un isomorphisme. Nous pouvons supposer que $M = P$ est un $E$-module
différentiel gradué qui a la propriété (P). Puisque $K^\bullet$ définit un objet
compact, nous nous ramenons, à l'aide du
lemme \ref{lemma-property-P-sequence}, au
cas où $P$ a une filtration finie dont les quotients successifs sont des sommes
directes de $E[k]$. Encore une fois, en utilisant la compacité, nous
nous ramenons au cas $P = E[k]$. L'hypothèse sur $K^\bullet$ est que le
résultat est vrai dans ces cas.
\end{proof}











\section{Caractérisation des objets compacts}
\label{section-compact}

\noindent
Les objets compacts des catégories additives sont définis
dans Catégories dérivées, définition \ref{derived-definition-compact-object}. Dans
cette section, nous caractérisons les objets compacts de la catégorie
dérivée d'une algèbre différentielle graduée.

\begin{remark}
\label{remark-source-graded-projective}
Soit $(A, \text{d})$ une algèbre différentielle graduée. Existe-t-il
une caractérisation des $A$-modules différentiels gradués $P$ pour
lesquels nous avons
$$
\Hom_{K(A, \text{d})}(P, M) = \Hom_{D(A, \text{d})}(P, M)
$$
pour tous les $A$-modules différentiels gradués $M$ ?
Soit $\mathcal{D} \subset K(A, \text{d})$ la sous-catégorie pleine
dont les objets sont les objets $P$ satisfaisant ce qui précède. Alors
$\mathcal{D}$ est une sous-catégorie triangulée strictement pleine et saturée de $K(A,
\text{d})$. Si $P$ est projectif en tant que $A$-module gradué, alors pour voir si
$P$ est un objet de $\mathcal{D}$, il suffit de vérifier que
$\Hom_{K(A, \text{d})}(P, M) = 0$ chaque fois que $M$ est
acyclique. Cependant, en général, il ne suffit pas de supposer que $P$ est projectif en
tant que $A$-module gradué. Exemple : prenons $A = R = k[\epsilon]$ où $k$ est
un corps et $k[\epsilon] = k[x]/(x^2)$ est l'anneau des nombres duaux. Soit
$P$ l'objet avec $P^n = R$ pour tous les $n \in \mathbf{Z}$ et de
différentielle donnée par multiplication par $\epsilon$. Alors
$\text{id}_P \in \Hom_{K(A, \text{d})}(P, P)$ est un élément non nul mais $P$ est
acyclique.
\end{remark}

\begin{remark}
\label{remark-graded-projective-is-compact}
Soit $(A, \text{d})$ une algèbre différentielle graduée. Disons qu'un
$A$-module différentiel gradué $M$ est {\it de type fini} si $M$ est engendré, en tant que
module à droite sur $A$, par un nombre fini d'éléments. Si $P$ est un
$A$-module différentiel gradué qui est projectif gradué de type fini, alors nous
pouvons nous demander : Est-ce que $P$ donne un objet compact de $D(A, \text{d})$ ? Nous
nous attendons à ce que ce ne soit pas vrai en général, mais nous ne
connaissons pas de contre-exemple. Cependant, si $P$ est également un objet de la
catégorie $\mathcal{D}$ de la remarque \ref{remark-source-graded-projective},
alors c'est le cas (cela découle du fait que les sommes directes dans $D(A,
\text{d})$ sont données par des sommes directes de modules ; les détails sont omis).
\end{remark}

\begin{lemma}
\label{lemma-factor-through-nicer}
Soit $(A, \text{d})$ une algèbre différentielle graduée. Soit $E$ un objet
compact de $D(A, \text{d})$. Soit $P$ un $A$-module différentiel gradué qui
possède une filtration finie
$$
0 = F_{-1}P \subset F_0P \subset F_1P \subset \ldots \subset F_nP = P
$$
par des sous-modules différentiels gradués, telle que
$$
F_{i + 1}P/F_iP \cong \bigoplus\nolimits_{j \in J_i} A[k_{i, j}]
$$
comme $A$-modules différentiels gradués pour certains ensembles $J_i$ et entiers
$k_{i, j}$. Soit $E \to P$ un morphisme de $D(A,
\text{d})$. Alors il existe un sous-module différentiel gradué $P' \subset P$ tel que
$F_{i + 1}P \cap P'/(F_iP \cap P')$ soit égal à
$\bigoplus_{j \in J'_i} A[k_{i, j}]$ pour certains sous-ensembles
finis $J'_i \subset J_i$ et tel que $E \to P$ se factorise par $P'$.
\end{lemma}

\begin{proof}
Nous allons démontrer par récurrence sur $-1 \leq m \leq n$ qu'il
existe un sous-module différentiel gradué $P' \subset P$ tel que
\begin{enumerate}
\item $F_mP \subset P'$,
\item pour $i \geq m$ le quotient $F_{i + 1}P \cap P'/(F_iP \cap P')$ est
isomorphe à $\bigoplus_{j \in J'_i} A[k_{i, j}]$ pour certains sous-ensembles
finis $J'_i \subset J_i$, et
\item $E \to P$ se factorise par $P'$.
\end{enumerate}
Le cas de base est $m = n$ où nous pouvons prendre $P' = P$.

\medskip\noindent
Étape de récurrence. Supposons que $P'$
convienne pour $m$. Pour $i \geq m$ et $j \in J'_i$, soit $x_{i, j} \in F_{i +
1}P \cap P'$ un élément homogène de degré $k_{i, j}$ dont l'image
dans $F_{i + 1}P \cap P'/(F_iP \cap P')$ est le générateur du
facteur correspondant à $j \in J_i$. Les $x_{i, j}$
engendrent $P'/F_mP$ en tant que $A$-module. Écrivons
$$
\text{d}(x_{i, j}) = \sum x_{i', j'} a_{i, j}^{i', j'} + y_{i, j}
$$
avec $y_{i, j} \in F_mP$ et $a_{i, j}^{i', j'} \in A$. Il
existe un sous-ensemble fini
$J'_{m - 1} \subset J_{m - 1}$ tel que l'image de chaque $y_{i, j}$
appartienne au sous-module $\bigoplus_{j \in J'_{m - 1}} A[k_{m - 1, j}]$ de
$F_mP/F_{m - 1}P$. Soit $P'' \subset F_mP$ l'image inverse de
$\bigoplus_{j \in J'_{m - 1}} A[k_{m - 1, j}]$ sous l'application
$F_mP \to F_mP/F_{m - 1}P$. Alors nous voyons que le $A$-sous-module
$$
P'' + \sum x_{i, j}A
$$
est un sous-module différentiel gradué du type que nous recherchons. De plus
$$
P'/(P'' + \sum x_{i, j}A) =
\bigoplus\nolimits_{j \in J_{m - 1} \setminus J'_{m - 1}} A[k_{m - 1, j}]
$$
Puisque $E$ est compact, la composition de l'application donnée $E \to P'$ avec
l'application quotient se factorise par la somme directe d'un nombre fini des
facteurs du module ci-dessus. Par conséquent, après avoir agrandi $J'_{m - 1}$,
nous pouvons supposer que $E \to P'$ se factorise
par $P'' + \sum x_{i, j}A$ comme souhaité.
\end{proof}

\noindent
Il n'est pas vrai que tout objet compact de $D(A, \text{d})$ provienne
d'un $A$-module différentiel gradué qui soit projectif gradué de type
fini, voir Exemples, section \ref{examples-section-interesting-compact}.

\begin{proposition}
\label{proposition-compact}
Soit $(A, \text{d})$ une algèbre différentielle graduée. Soit $E$ un objet
de $D(A, \text{d})$. Alors les assertions suivantes sont équivalentes
\begin{enumerate}
\item $E$ est un objet compact,
\item $E$ est un facteur direct d'un objet de $D(A, \text{d})$ qui est
représenté par un module différentiel gradué $P$ qui possède une
filtration finie $F_\bullet$ par des sous-modules différentiels gradués, telle que
les $F_iP/F_{i - 1}P$ soient des sommes directes finies de décalages de $A$.
\end{enumerate}
\end{proposition}

\begin{proof}
Supposons que $E$ soit compact. D'après le lemme \ref{lemma-resolve}, nous pouvons
supposer que $E$ est représenté par un $A$-module différentiel gradué $P$ avec la
propriété (P). Considérons le triangle distingué
$$
\bigoplus F_iP \to \bigoplus F_iP \to P
\xrightarrow{\delta} \bigoplus F_iP[1]
$$
provenant de la suite exacte courte admissible du lemme
\ref{lemma-property-P-sequence}. Puisque $E$ est compact, nous avons
$\delta = \sum_{i = 1, \ldots, n} \delta_i$ pour certains
$\delta_i : P \to F_iP[1]$. Puisque la composition de $\delta$ avec
l'application $\bigoplus F_iP[1] \to \bigoplus F_iP[1]$ est nulle (Catégories
dérivées, lemme \ref{derived-lemma-composition-zero}), il s'ensuit que
$\delta = 0$ (car $\bigoplus F_iP \to \bigoplus F_iP$ envoie le facteur direct $F_iP$
via la différence de $\text{id}$ et de l'application d'inclusion dans $F_{i - 1}P$).
Ainsi, nous voyons que
l'identité sur $E$ se factorise à travers $\bigoplus F_iP$ dans
$D(A, \text{d})$ (d'après Catégories
dérivées, lemme \ref{derived-lemma-split}). Ensuite, nous utilisons
à nouveau la compacité de $P$ pour voir que l'application $E \to
\bigoplus F_iP$ se factorise par $\bigoplus_{i = 1, \ldots, n} F_iP$ pour un
certain $n$. En d'autres termes, l'identité sur $E$ se factorise par
$\bigoplus_{i = 1, \ldots, n} F_iP$. Par le lemme
\ref{lemma-factor-through-nicer}, nous voyons
que l'identité de $E$ se factorise comme $E \to P \to E$ où $P$ est
comme dans la partie (2) de l'énoncé du lemme. En d'autres termes,
nous avons prouvé que (1) implique (2).

\medskip\noindent
Supposons (2).
D'après Catégories dérivées, lemme
\ref{derived-lemma-compact-objects-subcategory}, il suffit de montrer que $P$ donne un objet compact. Remarquons que
$P$ a la propriété (P), donc nous avons
$$
\Hom_{D(A, \text{d})}(P, M) = \Hom_{K(A, \text{d})}(P, M)
$$
pour tout module différentiel gradué $M$ d'après le lemme \ref{lemma-hom-derived}.
Comme les sommes directes dans $D(A, \text{d})$ sont données par des sommes
directes de modules gradués (lemme \ref{lemma-derived-products}), nous
nous ramenons à montrer que $\Hom_{K(A, \text{d})}(P, M)$ commute avec les sommes
directes. En utilisant le fait que $K(A, \text{d})$ est une catégorie
triangulée, que $\Hom$ est un foncteur cohomologique dans la
première variable, et la filtration sur $P$, nous nous ramenons au cas où $P$
est une somme directe finie de décalages de $A$. Ainsi, nous nous ramenons
au cas $P = A[k]$, ce qui est clair.
\end{proof}

\begin{lemma}
\label{lemma-compact-implies-bounded}
Soit $(A, \text{d})$ une algèbre différentielle
graduée. Pour tout objet compact $E$ de $D(A, \text{d})$, il
existe des entiers $a \leq b$ tels que $\Hom_{D(A, \text{d})}(E, M) = 0$
si $H^i(M) = 0$ pour $i \in [a, b]$.
\end{lemma}

\begin{proof}
Observons que la collection d'objets de $D(A, \text{d})$ pour lesquels un
tel couple d'entiers existe est une sous-catégorie triangulée saturée et strictement
pleine de $D(A, \text{d})$. Ainsi,
d'après la proposition \ref{proposition-compact}, il suffit de prouver cela
lorsque $E$ est représenté par un module différentiel gradué $P$ qui possède
une filtration finie $F_\bullet$ par des sous-modules différentiels gradués tels
que $F_iP/F_{i - 1}P$ soient des sommes directes finies de décalages de $A$.
En utilisant la compatibilité avec les triangles, nous voyons qu'il
suffit de le prouver pour $P = A$. Dans ce cas, $\Hom_{D(A, \text{d})}(A, M) = H^0(M)$ et
le résultat vaut pour $a = b = 0$.
\end{proof}

\noindent
Si $(A, \text{d})$ est simplement une algèbre placée en degré $0$
avec une différentielle nulle ou, plus généralement,
est concentrée en un nombre fini de degrés, alors nous obtenons
la description plus précise des objets compacts.

\begin{lemma}
\label{lemma-compact}
Soit $(A, \text{d})$ une algèbre différentielle graduée. Supposons que $A^n = 0$
pour $|n| \gg 0$. Soit $E$ un objet de $D(A, \text{d})$. Les
assertions suivantes sont équivalentes
\begin{enumerate}
\item $E$ est un objet compact, et
\item $E$ peut être représenté par un $A$-module différentiel gradué $P$
qui est projectif de type fini en tant que $A$-module gradué et
satisfait $\Hom_{K(A, \text{d})}(P, M) = \Hom_{D(A, \text{d})}(P,
M)$ pour tout $A$-module différentiel gradué $M$.
\end{enumerate}
\end{lemma}

\begin{proof}
Soit $\mathcal{D} \subset K(A, \text{d})$ la sous-catégorie triangulée
discutée dans la remarque \ref{remark-source-graded-projective}.
Soit $P$ un objet de $\mathcal{D}$ qui est projectif de type fini comme
$A$-module gradué. Alors $P$ représente un objet compact de $D(A,
\text{d})$ d'après la remarque \ref{remark-graded-projective-is-compact}.

\medskip\noindent
Pour prouver la réciproque, soit $E$ un objet compact de $D(A, \text{d})$.
Fixons $a \leq b$ comme dans le lemme
\ref{lemma-compact-implies-bounded}. Après avoir diminué $a$ et augmenté $b$ si nécessaire, nous pouvons
également supposer que $H^i(E) = 0$ pour $i \not \in [a, b]$ (cela
découle de la proposition \ref{proposition-compact} et de notre hypothèse sur
$A$). De plus, fixons un entier $c > 0$ tel que $A^n = 0$ si $|n| \geq c$.

\medskip\noindent
Par la proposition \ref{proposition-compact} nous voyons que $E$ est un
facteur direct, dans $D(A, \text{d})$, d'un $A$-module différentiel gradué $P$ qui
possède une filtration finie $F_\bullet$ par des sous-modules
différentiels gradués telle que les $F_iP/F_{i - 1}P$ soient des sommes directes
finies de décalages de $A$. En particulier, $P$ a la propriété (P) et nous
avons $\Hom_{D(A, \text{d})}(P, M) = \Hom_{K(A, \text{d})}(P, M)$ pour tout
module différentiel gradué $M$ selon le lemme \ref{lemma-hom-derived}. En
d'autres termes, $P$ est un objet de la sous-catégorie
triangulée $\mathcal{D} \subset K(A, \text{d})$ discutée dans la remarque
\ref{remark-source-graded-projective}. Notons que
$P$ est libre de type fini comme $A$-module gradué.

\medskip\noindent
Choisissons $n > 0$ tel que $b + 4c - n <
a$. Représentons le projecteur sur $E$ par un endomorphisme $\varphi : P \to P$
de $A$-modules différentiels gradués. Considérons le triangle distingué
$$
P \xrightarrow{1 - \varphi} P \to C \to P[1]
$$
dans $K(A, \text{d})$ où $C$ est le cône de la première flèche. Ensuite,
$C$ est un objet de $\mathcal{D}$, nous
avons $C \cong E \oplus E[1]$ dans $D(A, \text{d})$, et $C$
est un $A$-module gradué libre de type fini.
Ensuite, considérons un triangle distingué
$$
C[1] \to C \to C' \to C[2]
$$
dans $K(A, \text{d})$ où $C'$ est le cône sur un morphisme $C[1] \to
C$ représentant la composition
$$
C[1] \cong E[1] \oplus E[2] \to E[1] \to E \oplus E[1] \cong C
$$
dans $D(A, \text{d})$. Ensuite, nous voyons que $C'$ représente $E \oplus E[2]$.
En continuant de cette manière, nous voyons que nous pouvons trouver un
$A$-module différentiel gradué $P$ qui est un objet de $\mathcal{D}$,
est libre de type fini en tant que $A$-module gradué, et représente $E \oplus E[n]$.

\medskip\noindent
Choisissons une base $x_i$, $i \in I$ d'éléments homogènes pour $P$ en
tant que $A$-module. Posons $d_i =
\deg(x_i)$. Soit $P_1$ le $A$-sous-module de $P$ engendré par $x_i$
et $\text{d}(x_i)$ pour $d_i \leq a - c - 1$.
Soit $P_2$ le $A$-sous-module de $P$ engendré par $x_i$ et
$\text{d}(x_i)$ pour $d_i \geq b - n + c$. Nous
observons
\begin{enumerate}
\item $P_1$ et $P_2$ sont des sous-modules différentiels gradués de $P$,
\item $P_1^t = 0$ pour $t \geq a$,
\item $P_1^t = P^t$ pour $t \leq a - 2c$,
\item $P_2^t = 0$ pour $t \leq b - n$,
\item $P_2^t = P^t$ pour $t \geq b - n + 2c$.
\end{enumerate}
Comme $b - n + 2c \geq a - 2c$ par notre choix de $n$
nous obtenons une suite exacte courte de $A$-modules différentiels gradués
$$
0 \to P_1 \cap P_2 \to P_1 \oplus P_2 \xrightarrow{\pi} P \to 0
$$
Puisque $P$ est projectif en tant que $A$-module gradué, il s'agit d'une suite
exacte courte admissible (lemme \ref{lemma-target-graded-projective}).
Ainsi, nous obtenons un morphisme de
liaison $\delta : P \to (P_1 \cap P_2)[1]$ dans $K(A, \text{d})$, voir
le lemme \ref{lemma-admissible-ses}.
Puisque $P = E \oplus E[n]$ et puisque $P_1 \cap P_2$ est concentré en
degrés $(b - n, a)$, nous constatons que
$\Hom_{D(A, \text{d})}(E \oplus E[n], (P_1 \cap P_2)[1])$ est nul. Par
conséquent, $\delta = 0$ en tant que morphisme dans $K(A, \text{d})$
puisque $P$ est un objet de $\mathcal{D}$. Par
Catégories dérivées, le lemme \ref{derived-lemma-split} nous
permet de trouver une application $s : P \to P_1 \oplus P_2$
telle que $\pi \circ s = \text{id}_P + \text{d}h + h\text{d}$ pour un certain $h : P \to
P$ de degré $-1$. Puisque $P_1 \oplus P_2 \to P$ est surjectif et que $P$ est
projectif comme $A$-module gradué, nous pouvons choisir un relèvement homogène
$\tilde h : P \to P_1 \oplus P_2$ de $h$. Puis nous changeons $s$ en $s +
\text{d} \tilde h + \tilde h \text{d}$ pour obtenir $\pi \circ s =
\text{id}_P$. Cela signifie que nous obtenons une décomposition en
somme directe $P = s^{-1}(P_1) \oplus s^{-1}(P_2)$. Puisque
$s^{-1}(P_2)$ est égal à $P$ en degrés $\geq b - n + 2c$, nous voyons que
$s^{-1}(P_2) \to P \to E$ est un quasi-isomorphisme, c'est-à-dire un
isomorphisme dans $D(A, \text{d})$. Cela termine la démonstration.
\end{proof}








\section{Équivalences de catégories dérivées}
\label{section-equivalence}

\noindent
Soit $R$ un anneau. Soient $(A, \text{d})$ et $(B, \text{d})$ des $R$-algèbres
différentielles graduées. Une question naturelle est de savoir ce que
signifie l'équivalence de $D(A, \text{d})$ avec $D(B, \text{d})$ en tant que
catégorie triangulée linéaire sur $R$. C'est une question plutôt subtile et il
s'avérera qu'elle n'est pas toujours la bonne question à poser.
Néanmoins, dans cette section nous recueillons certaines conditions qui
garantissent que c'est le cas.

\medskip\noindent
Nous recommandons vivement au lecteur de consulter l'article
fondateur \cite{Rickard} sur ce sujet.

\begin{lemma}
\label{lemma-qis-equivalence}
Soit $R$ un anneau. Soit $(A, \text{d}) \to (B, \text{d})$ un
homomorphisme d'algèbres différentielles graduées sur $R$, qui induit un
isomorphisme sur les algèbres de cohomologie. Alors
$$
- \otimes_A^\mathbf{L} B : D(A, \text{d}) \to D(B, \text{d})
$$
donne une équivalence $R$-linéaire de catégories triangulées avec
quasi-inverse le foncteur de restriction $N \mapsto N_A$.
\end{lemma}

\begin{proof}
D'après le lemme
\ref{lemma-tensor-with-compact-fully-faithful}, le foncteur $M \longmapsto M \otimes_A^\mathbf{L} B$ est
pleinement fidèle. D'après le lemme
\ref{lemma-tensor-hom-adjoint}, le foncteur $N \longmapsto R\Hom(B, N) = N_A$ est un adjoint à droite,
voir l'exemple
\ref{example-map-hom-tensor}. Il est clair que le noyau de $R\Hom(B, -)$ est nul. Par
conséquent, le résultat découle de
Catégories dérivées, lemme
\ref{derived-lemma-fully-faithful-adjoint-kernel-zero}.
\end{proof}

\noindent
Lorsque nous analysons la démonstration ci-dessus, nous voyons que nous obtenons sans
travail supplémentaire la généralisation suivante.

\begin{lemma}
\label{lemma-tilting-equivalence}
Soit $R$ un anneau. Soient $(A, \text{d})$ et $(B, \text{d})$ des algèbres
différentielles graduées sur $R$. Soit $N$ un bimodule
différentiel gradué sur $(A, B)$. Faisons les hypothèses suivantes :
\begin{enumerate}
\item $N$ définit un objet compact de $D(B, \text{d})$,
\item si $N' \in D(B, \text{d})$ et
$\Hom_{D(B, \text{d})}(N, N'[n]) = 0$ pour $n \in \mathbf{Z}$,
alors $N' = 0$, et
\item l'application $H^k(A) \to \Hom_{D(B, \text{d})}(N, N[k])$ est
un isomorphisme pour tout $k \in \mathbf{Z}$.
\end{enumerate}
Alors
$$
- \otimes_A^\mathbf{L} N : D(A, \text{d}) \to D(B, \text{d})
$$
donne une équivalence $R$-linéaire de catégories triangulées.
\end{lemma}

\begin{proof}
D'après le lemme
\ref{lemma-tensor-with-compact-fully-faithful} le foncteur $M \longmapsto M \otimes_A^\mathbf{L} N$ est
pleinement fidèle. D'après le lemme
\ref{lemma-tensor-hom-adjoint} le foncteur $N' \longmapsto R\Hom(N, N')$ est un adjoint à droite.
D'après l'hypothèse (3) le noyau de $R\Hom(N, -)$ est nul. Par
conséquent, le résultat découle
de Catégories dérivées, lemme
\ref{derived-lemma-fully-faithful-adjoint-kernel-zero}.
\end{proof}

\begin{remark}
\label{remark-tilting-equivalence}
Dans le lemme \ref{lemma-tilting-equivalence}, nous pouvons
remplacer la condition (2) par la condition selon laquelle $N$ est un
générateur classique pour $D_{compact}(B,
d)$, voir Catégories dérivées,
proposition \ref{derived-proposition-generator-versus-classical-generator}. De plus,
si nous savions que $R\Hom(N, B)$ est un objet compact de $D(A,
\text{d})$, alors il suffirait de vérifier que $N$ est un générateur
faible pour $D_{compact}(B, \text{d})$. Nous omettons la
démonstration ; nous l'ajouterons ici si nous en avons besoin un
jour dans le projet Champs.
\end{remark}

\noindent
Le $B$-module $P$ du lemme ci-dessous est parfois appelé
« complexe $(A, B)$-basculant ».

\begin{lemma}
\label{lemma-rickard}
Soit $R$ un anneau. Soient $(A, \text{d})$ et $(B, \text{d})$ des
$R$-algèbres différentielles graduées. Supposons que $A = H^0(A)$. Les
assertions suivantes sont équivalentes
\begin{enumerate}
\item $D(A, \text{d})$ et $D(B, \text{d})$ sont équivalentes en tant que
catégories triangulées $R$-linéaires, et
\item il existe un objet $P$ de $D(B, \text{d})$ tel que
\begin{enumerate}
\item $P$ est un objet compact de $D(B, \text{d})$,
\item si $N \in D(B, \text{d})$ avec $\Hom_{D(B, \text{d})}(P, N[i]) = 0$
pour $i \in \mathbf{Z}$, alors $N = 0$,
\item $\Hom_{D(B, \text{d})}(P, P[i]) = 0$ pour $i \not = 0$ et
est égal à $A$ pour $i = 0$.
\end{enumerate}
\end{enumerate}
L'équivalence $D(A, \text{d}) \to D(B, \text{d})$
construite dans (2) envoie $A$ sur $P$.
\end{lemma}

\begin{proof}
Soit $F : D(A, \text{d}) \to D(B, \text{d})$ une équivalence. Alors $F$
envoie les objets compacts vers des objets compacts. Par conséquent, $P = F(A)$ est
compact, c'est-à-dire que (2)(a) est vérifié. Les conditions (2)(b) et (2)(c) découlent
immédiatement du fait que $F$ est une équivalence.

\medskip\noindent
Soit $P$ un objet comme en (2). Représentons $P$ par un
module différentiel gradué avec la propriété (P). Posons
$$
(E, \text{d}) = \Hom_{\text{Mod}^{dg}_{(B, \text{d})}}(P, P)
$$
Alors $H^0(E) = A$ et $H^k(E) = 0$ pour $k \not = 0$ par le
lemme \ref{lemma-hom-derived} et l'hypothèse (2)(c). En
considérant $P$ comme un $(E, B)$-bimodule et en utilisant
le lemme \ref{lemma-tilting-equivalence} et l'hypothèse (2)(b), nous
obtenons une équivalence
$$
D(E, \text{d}) \to D(B, \text{d})
$$
envoyant $E$ sur $P$.
Soit $E' \subset E$ la sous-algèbre différentielle graduée sur $R$
avec
$$
(E')^i = \left\{
\begin{matrix}
E^i & \text{si }i < 0 \\
\Ker(E^0 \to E^1) & \text{si }i = 0 \\
0 & \text{si }i > 0
\end{matrix}
\right.
$$
Il y a alors des quasi-isomorphismes d'algèbres
différentielles graduées $(A, \text{d}) \leftarrow (E', \text{d}) \rightarrow (E, \text{d})$. Ainsi,
nous obtenons des équivalences
$$
D(A, \text{d}) \leftarrow D(E', \text{d}) \rightarrow D(E, \text{d})
\rightarrow D(B, \text{d})
$$
par le lemme \ref{lemma-qis-equivalence}.
\end{proof}

\begin{remark}
\label{remark-lift-equivalence-to-dga}
Soit $R$ un anneau. Soient $(A, \text{d})$ et $(B, \text{d})$ des $R$-algèbres
différentielles graduées. Supposons donnée une équivalence $R$-linéaire
$$
F : D(A, \text{d}) \longrightarrow D(B, \text{d})
$$
de catégories triangulées. Posons $N = F(A)$. Alors $N$ est un $B$-module
différentiel gradué. Puisque $F$ est une équivalence et que $A$ est un objet
compact de $D(A, \text{d})$, nous concluons que $N$ est un objet compact de
$D(B, \text{d})$. Puisque $A$ engendre $D(A, \text{d})$ et que $F$ est
une équivalence, nous voyons que $N$ engendre $D(B, \text{d})$. Enfin,
$H^k(A) = \Hom_{D(A, \text{d})}(A, A[k])$ et comme $F$ est une équivalence, nous voyons
que $F$ induit un isomorphisme $H^k(A) =
\Hom_{D(B, \text{d})}(N, N[k])$ pour tous les $k$. Afin de
conclure qu'il existe une équivalence $D(A, \text{d})
\longrightarrow D(B, \text{d})$ qui découle de la construction dans
le lemme
\ref{lemma-tilting-equivalence} tout ce dont nous avons besoin est
une structure de $A$-module à gauche sur $N$ compatible
avec la différentielle et commutant avec la
structure de $B$-module à droite donnée. En fait, il suffit de
le faire après avoir remplacé $N$ par un $B$-module différentiel
gradué quasi-isomorphe. La structure du
module peut être construite dans certains cas. Par exemple, si
nous supposons que $F$ peut être relevé à un foncteur
différentiel gradué
$$
F^{dg} :
\text{Mod}^{dg}_{(A, \text{d})}
\longrightarrow
\text{Mod}^{dg}_{(B, \text{d})}
$$
(pour la notation voir exemple \ref{example-dgm-dg-cat})
entre les catégories différentielles graduées associées,
alors cela est vrai. Un autre cas est discuté dans la proposition ci-dessous.
\end{remark}

\begin{proposition}
\label{proposition-rickard}
Soit $R$ un anneau. Soient $(A, \text{d})$ et $(B, \text{d})$ des
$R$-algèbres différentielles graduées. Soit $F : D(A, \text{d}) \to D(B, \text{d})$ une
équivalence $R$-linéaire de catégories triangulées. Faisons les hypothèses suivantes :
\begin{enumerate}
\item $A = H^0(A)$, et
\item $B$ est K-plat en tant que complexe de $R$-modules.
\end{enumerate}
Alors il existe un $(A, B)$-bimodule $N$ comme dans
le lemme \ref{lemma-tilting-equivalence}.
\end{proposition}

\begin{proof}
Comme dans la remarque \ref{remark-lift-equivalence-to-dga} ci-dessus, nous posons $N =
F(A)$ dans $D(B, \text{d})$. Nous pouvons supposer que $N$ est un $B$-module
différentiel gradué avec la propriété (P). Posons
$$
(E, \text{d}) = \Hom_{\text{Mod}^{dg}_{(B, \text{d})}}(N, N)
$$
Puis $H^0(E) = A$ et $H^k(E) = 0$ pour $k \not = 0$ par le
lemme \ref{lemma-hom-derived}. De
plus, d'après la discussion dans la remarque \ref{remark-lift-equivalence-to-dga} et
par le lemme \ref{lemma-tilting-equivalence},
nous voyons que $N$ en tant que $(E, B)$-bimodule induit
une équivalence $- \otimes_E^\mathbf{L} N : D(E, \text{d}) \to D(B, \text{d})$. Soit
$E' \subset E$ la sous-algèbre différentielle graduée sur $R$
avec
$$
(E')^i = \left\{
\begin{matrix}
E^i & \text{si }i < 0 \\
\Ker(E^0 \to E^1) & \text{si }i = 0 \\
0 & \text{si }i > 0
\end{matrix}
\right.
$$
Ensuite, il y a des quasi-isomorphismes d'algèbres
différentielles graduées $(A, \text{d}) \leftarrow (E', \text{d}) \rightarrow (E, \text{d})$. Ainsi,
nous obtenons des équivalences
$$
D(A, \text{d}) \leftarrow D(E', \text{d}) \rightarrow D(E, \text{d})
\rightarrow D(B, \text{d})
$$
d'après le lemme
\ref{lemma-qis-equivalence}. Notons que le quasi-inverse $D(A, \text{d}) \to D(E', \text{d})$
de la flèche verticale gauche est
donné par $M \mapsto M \otimes_A^\mathbf{L} A$ où $A$ est considéré comme
un bimodule $(A, E')$, voir l'exemple \ref{example-map-hom-tensor}.
D'autre part, le foncteur $D(E', \text{d}) \to D(B, \text{d})$ est donné par $M
\mapsto M \otimes_{E'}^\mathbf{L} N$ où $N$ est comme ci-dessus. Nous
concluons par le lemme \ref{lemma-compose-tensor-functors}.
\end{proof}

\begin{remark}
\label{remark-rickard}
Soient $A, B, F, N$ comme dans la proposition \ref{proposition-rickard}.
Il n'est pas clair que $F$ et le foncteur
$G(-) = - \otimes_A^\mathbf{L} N$ soient isomorphes. Par
construction, il existe un isomorphisme $N =
G(A) \to F(A)$ dans $D(B, \text{d})$. Il est
simple d'étendre cela à un isomorphisme fonctoriel $G(M) \to F(M)$ lorsque
$M$ est un $A$-module différentiel gradué qui est projectif gradué (par
exemple, une somme de décalages de $A$). On peut alors
conclure que $G(M) \cong F(M)$ lorsque $M$ est un cône d'une application entre
de tels modules. Nous ne savons pas si l'on peut en dire davantage en
général.
\end{remark}

\begin{lemma}
\label{lemma-rickard-rings}
Soit $R$ un anneau.
Soient $A$ et $B$ des $R$-algèbres. Les assertions suivantes sont équivalentes
\begin{enumerate}
\item il existe une équivalence $R$-linéaire $D(A) \to
D(B)$ de catégories triangulées,
\item il existe un objet $P$ de $D(B)$ tel que
\begin{enumerate}
\item $P$ peut être représenté par un complexe borné
de $B$-modules projectifs de type fini,
\item si $K \in D(B)$ avec $\Ext^i_B(P, K) = 0$ pour
$i \in \mathbf{Z}$, alors $K = 0$, et
\item $\Ext^i_B(P, P) = 0$ pour $i \not = 0$ et
est égal à $A$ pour $i= 0$.
\end{enumerate}
\end{enumerate}
De plus, si $B$ est plat en tant que $R$-module, alors cela
équivaut également à
\begin{enumerate}
\item[(3)] il existe un $(A, B)$-bimodule $N$ tel que $-
\otimes_A^\mathbf{L} N : D(A) \to D(B)$ soit une équivalence.
\end{enumerate}
\end{lemma}

\begin{proof}
L'équivalence de (1) et (2) est un cas particulier
du lemme \ref{lemma-rickard} combiné avec le résultat du
lemme \ref{lemma-compact} caractérisant les objets compacts de
$D(B)$ (petit détail omis).
L'équivalence avec (3) si $B$ est $R$-plat découle de la
proposition \ref{proposition-rickard}.
\end{proof}

\begin{remark}
\label{remark-centers}
Soit $R$ un anneau. Soient $A$ et $B$ des $R$-algèbres. Si
$D(A)$ et $D(B)$ sont équivalentes comme catégories triangulées
$R$-linéaires, alors les centres de $A$ et $B$ sont isomorphes comme
$R$-algèbres. En particulier, si $A$ et $B$ sont commutatifs, alors $A \cong B$.
La démonstration plutôt délicate peut être trouvée dans
\cite[Proposition 9.2]{Rickard} ou \cite[Proposition 6.3.2]{KZ}. Une autre approche
pourrait être d'utiliser la cohomologie de Hochschild (voir remarque
ci-dessous).
\end{remark}

\begin{remark}
\label{remark-hochschild-cohomology}
Soit $R$ un anneau. Soient $(A, \text{d})$ et $(B, \text{d})$ des $R$-algèbres
différentielles graduées dont les catégories dérivées sont équivalentes, c'est-à-dire
telles qu'il existe une équivalence $R$-linéaire $D(A, \text{d}) \to D(B,
\text{d})$ de catégories triangulées. Nous souhaitons montrer que certains invariants
de $(A, \text{d})$ et $(B, \text{d})$ coïncident. Dans de nombreuses
situations, on a un meilleur contrôle de la situation. Par exemple, il peut arriver
qu'il existe une équivalence de la forme
$$
- \otimes_A \Omega : D(A, \text{d}) \longrightarrow D(B, \text{d})
$$
pour un certain $(A, B)$-bimodule différentiel gradué
$\Omega$ (cela se produit dans la situation de la
proposition \ref{proposition-rickard} et est souvent vrai si
l'équivalence provient d'une construction géométrique). Supposons de
plus que le quasi-inverse de notre foncteur soit donné par
$$
- \otimes_B^\mathbf{L} \Omega' : D(B, \text{d}) \longrightarrow D(A, \text{d})
$$
pour un $(B, A)$-bimodule différentiel gradué $\Omega'$ (et
comme précédemment un tel module $\Omega'$ existe souvent en pratique).
Dans ce cas, nous pouvons considérer le foncteur
$$
D(A^{opp} \otimes_R A, \text{d})
\longrightarrow
D(B^{opp} \otimes_R B, \text{d}),\quad
M \longmapsto \Omega' \otimes^\mathbf{L}_A M \otimes_A^\mathbf{L} \Omega
$$
sur les catégories dérivées de bimodules
(utilisez le lemme \ref{lemma-bimodule-over-tensor} pour transformer les
bimodules en
modules à droite). Observons que ce foncteur envoie le $(A, A)$-bimodule
$A$ sur le $(B, B)$-bimodule $B$. Sous des conditions adéquates
(par exemple, la platitude de $A$, $B$, $\Omega$ sur $R$,
etc.) ce foncteur sera également une équivalence. Si
tel est le cas, il s'ensuit que nous avons des isomorphismes de groupes
de cohomologie de Hochschild
$$
HH^i(A, \text{d}) =
\Hom_{D(A^{opp} \otimes_R A, \text{d})}(A, A[i])
\longrightarrow
\Hom_{D(B^{opp} \otimes_R B, \text{d})}(B, B[i]) =
HH^i(B, \text{d}).
$$
Par exemple, si $A = H^0(A)$, alors $HH^0(A, \text{d})$ est
égal au centre de $A$, et cela donne une démonstration conceptuelle du
résultat mentionné dans la remarque \ref{remark-centers}. Si jamais nous
avons besoin de cette remarque, nous fournirons ici un énoncé précis avec
une démonstration détaillée.
\end{remark}



\section{Résolutions d'algèbres différentielles graduées}
\label{section-resolution-dgas}

\noindent
Soit $R$ un anneau. Sous nos hypothèses, la $R$-algèbre libre
$R\langle S \rangle$ sur un ensemble $S$ est l'algèbre ayant pour $R$-base
les expressions
$$
s_1 s_2 \ldots s_n
$$
où $n \geq 0$ et où $s_1, \ldots, s_n \in S$ forment une suite
d'éléments de $S$. La multiplication est donnée par concaténation
$$
(s_1 s_2 \ldots s_n) \cdot (s'_1 s'_2 \ldots s'_m) =
s_1 \ldots s_n s'_1 \ldots s'_m
$$
Cette algèbre est caractérisée par le fait que l'application
$$
\Mor_{R\text{-alg}}(R\langle S \rangle, A) \to
\text{Map}(S, A),\quad
\varphi \longmapsto (s \mapsto \varphi(s))
$$
est une bijection pour chaque $R$-algèbre $A$.

\medskip\noindent
Dans la catégorie des $R$-algèbres graduées, notre ensemble $S$ devrait être
muni d'une graduation, que nous considérons comme une application $\deg : S \to
\mathbf{Z}$. Alors $R\langle S\rangle$ possède une graduation telle que les monômes
aient un degré
$$
\deg(s_1 s_2 \ldots s_n) = \deg(s_1) + \ldots + \deg(s_n)
$$
Dans ce contexte, l'application
$$
\Mor_{\text{}R\text{-alg graduées}}(R\langle S \rangle, A) \to
\text{Map}_{\text{ensembles gradués}}(S,
A),\quad \varphi \longmapsto (s \mapsto \varphi(s))
$$
est une bijection pour chaque $R$-algèbre graduée $A$.

\medskip\noindent
Si $A$ est une $R$-algèbre graduée et que $S$ est un ensemble gradué,
alors nous pouvons de manière similaire former $A\langle S
\rangle$. Les éléments de $A\langle S \rangle$ sont
des sommes d'éléments de la forme
$$
a_0 s_1 a_1 s_2 \ldots a_{n - 1} s_n a_n
$$
avec $a_i \in A$ modulo les relations selon lesquelles ces
expressions sont $R$-multilinéaires en $(a_0, \ldots,
a_n)$. Ainsi, pour chaque suite $s_1, \ldots, s_n$ d'éléments de $S$,
il y a une inclusion
$$
A \otimes_R \ldots \otimes_R A \subset A\langle S \rangle
$$
et l'algèbre est la somme directe de celles-ci. Avec cette définition, on
vérifie que l'application
$$
\Mor_{\text{}R\text{-alg graduées}}(A\langle S \rangle, B) \to
\Mor_{\text{}R\text{-alg graduées}}(A, B) \times
\text{Map}_{\text{ensembles gradués}}(S, B),
$$
qui envoie $\varphi$ sur $(\varphi|_A, (s \mapsto \varphi(s)))$ est
une bijection pour chaque $R$-algèbre graduée $B$. Nous
observons que si $A$ était une $R$-algèbre graduée libre,
alors $A\langle S \rangle$ l'est aussi.

\medskip\noindent
Supposons que $A$ soit une $R$-algèbre différentielle graduée
et que $S$ soit un ensemble gradué. Supposons de plus que pour
chaque $s \in S$, on nous donne un élément homogène $f_s \in A$ avec $\deg(f_s) = \deg(s)
+ 1$ et $\text{d}f_s = 0$. Alors il existe une structure unique
d'algèbre différentielle graduée sur $A\langle S \rangle$ avec
$\text{d}(s) = f_s$. Par exemple, étant donnés $a, b, c \in A$ et
$s, t \in S$, on définirait
\begin{align*}
\text{d}(asbtc)
& =
\text{d}(a)sbtc + (-1)^{\deg(a)}a f_s b t c +
(-1)^{\deg(a) + \deg(s)} as\text{d}(b)tc \\ &
+ (-1)^{\deg(a) + \deg(s) + \deg(b)} asb f_t c +
(-1)^{\deg(a) + \deg(s) + \deg(b) + \deg(t)} asbt\text{d}(c)
\end{align*}
Nous omettons les détails.

\begin{lemma}
\label{lemma-K-flat-resolution}
Soit $R$ un anneau. Soit $(B, \text{d})$ une $R$-algèbre différentielle
graduée. Il existe un quasi-isomorphisme $(A, \text{d}) \to (B, \text{d})$ de
$R$-algèbres différentielles graduées avec les propriétés suivantes
\begin{enumerate}
\item $A$ est K-plat en tant que complexe de $R$-modules,
\item $A$ est une $R$-algèbre graduée libre.
\end{enumerate}
\end{lemma}

\begin{proof}
Tout d'abord, nous affirmons que nous pouvons trouver $(A_0, \text{d}) \to
(B, \text{d})$ vérifiant (1) et (2) et induisant une surjection
en cohomologie. En effet, prenons un ensemble gradué $S$ et pour
chaque $s \in S$ un élément homogène $b_s \in \Ker(d : B \to B)$ de degré $\deg(s)$
tel que les classes $\overline{b}_s$ dans $H^*(B)$
engendrent $H^*(B)$ en tant que $R$-module.
Ensuite, nous pouvons définir $A_0 = R\langle S \rangle$ avec une différentielle nulle
et $A_0 \to B$ défini en envoyant $s$ sur $b_s$.

\medskip\noindent
Étant donné $A_0 \to B$ induisant une surjection en cohomologie, nous
construisons une suite
$$
A_0 \to A_1 \to A_2 \to \ldots B
$$
par récurrence. Étant donné $A_n \to B$, prenons un ensemble gradué
$S_n$ et, pour chaque $s \in S_n$, soit $a_s \in \Ker(\text{d} : A_n \to A_n)$
un élément homogène de degré $\deg(s) + 1$ dont l'image
est une classe $\overline{a}_s$ dans $H^*(A_n)$ qui
s'envoie sur zéro dans $H^*(B)$. Nous choisissons $S_n$ assez grand pour
que les éléments $\overline{a}_s$ engendrent $\Ker(H^*(A_n) \to H^*(B))$ comme
$R$-module. Ensuite, nous posons
$$
A_{n + 1} = A_n\langle S_n \rangle
$$
avec la différentielle donnée par $\text{d}(s) = a_s$ ; voir la discussion
ci-dessus. Alors chaque $(A_n, \text{d})$ satisfait (1) et (2), nous omettons les détails.
L'application $H^*(A_n) \to H^*(B)$ est surjective comme c'était le cas pour $n = 0$.

\medskip\noindent
Il est clair que $A = \colim A_n$ est une $R$-algèbre graduée libre.
Elle est K-plate d'après Compléments d'algèbre, lemme \ref{more-algebra-lemma-colimit-K-flat}.
L'application $H^*(A) \to H^*(B)$ est un isomorphisme car elle est surjective et
injective : tout élément de $H^*(A)$ provient d'un élément de $H^*(A_n)$
pour un certain $n$ et s'il devient nul dans $H^*(B)$, alors il devient nul
dans $H^*(A_{n + 1})$, donc dans $H^*(A)$.
\end{proof}

\noindent
Comme application, nous prouvons la version « correcte »
du lemme \ref{lemma-compose-tensor-functors-general-algebra}.

\begin{lemma}
\label{lemma-compose-tensor-functors-tor}
Soit $R$ un anneau. Soient $(A, \text{d})$, $(B, \text{d})$, et
$(C, \text{d})$ des $R$-algèbres différentielles graduées.
Supposons que $A \otimes_R C$ représente $A \otimes^\mathbf{L}_R C$ dans
$D(R)$. Soit $N$ un $(A, B)$-bimodule différentiel
gradué. Soit $N'$ un $(B, C)$-bimodule différentiel gradué.
Alors la composition
$$
\xymatrix{
D(A, \text{d}) \ar[rr]^{- \otimes_A^\mathbf{L} N} & &
D(B, \text{d}) \ar[rr]^{- \otimes_B^\mathbf{L} N'} & &
D(C, \text{d})
}
$$
est isomorphe à $- \otimes_A^\mathbf{L} N''$ pour un certain $(A, C)$-bimodule
différentiel gradué $N''$.
\end{lemma}

\begin{proof}
En utilisant le lemme
\ref{lemma-K-flat-resolution}, nous choisissons un quasi-isomorphisme $(B', \text{d}) \to (B,
\text{d})$ avec $B'$ K-plat en tant que complexe de
$R$-modules. D'après le lemme
\ref{lemma-qis-equivalence}, le foncteur $-\otimes^\mathbf{L}_{B'} B : D(B', \text{d}) \to D(B,
\text{d})$ est une équivalence dont le quasi-inverse est donné par la
restriction. Notons que la restriction est canoniquement isomorphe au
foncteur $- \otimes^\mathbf{L}_B B : D(B, \text{d}) \to D(B', \text{d})$ où
$B$ est considéré comme un bimodule sur $(B, B')$.
Il suffit donc de prouver le lemme pour les compositions
$$
D(A) \to D(B) \to D(B'),\quad
D(B') \to D(B) \to D(C),\quad
D(A) \to D(B') \to D(C).
$$
Le premier cas résulte du lemme
\ref{lemma-compose-tensor-functors} parce que $B'$ est K-plat en tant que complexe de
$R$-modules. Le second cas est vrai parce
que $B \otimes_B^\mathbf{L} N' = N' = B \otimes_B N'$ et
donc le lemme \ref{lemma-compose-tensor-functors-general} s'applique. Ainsi,
nous nous ramenons au cas où $B$ est K-plat en tant que complexe de
$R$-modules.

\medskip\noindent
Supposons que $B$ soit K-plat en tant que complexe de $R$-modules. Il
suffit de montrer que (\ref{equation-plain-versus-derived-algebras}) est
un isomorphisme,
voir le lemme \ref{lemma-compose-tensor-functors-general-algebra}.
Choisissons un quasi-isomorphisme $L \to A$ où $L$ est un $R$-module
différentiel gradué qui a la propriété (P). Alors il est clair que $P = L
\otimes_R B$ a la propriété (P) en tant que $B$-module différentiel gradué. Par
conséquent, nous devons montrer que $P \to A \otimes_R B$
induit un quasi-isomorphisme
$$
P \otimes_B (B \otimes_R C)
\longrightarrow
(A \otimes_R B) \otimes_B (B \otimes_R C)
$$
Nous pouvons réécrire ceci comme
$$
P \otimes_R B \otimes_R C \longrightarrow A \otimes_R B \otimes_R C
$$
Puisque $B$ est K-plat en tant que complexe de
$R$-modules, il découle
de Compléments d'algèbre, lemme
\ref{more-algebra-lemma-K-flat-quasi-isomorphism} qu'il suffit de
montrer que
$$
P \otimes_R C \to A \otimes_R C
$$
est un quasi-isomorphisme, ce qui est exactement notre hypothèse.
\end{proof}

\noindent
Le lemme suivant n'appartient pas vraiment à cette section, mais il ne
semble pas y avoir de place naturelle appropriée pour lui.

\begin{lemma}
\label{lemma-countable}
Soit $(A, \text{d})$ une algèbre différentielle graduée avec
$H^i(A)$ dénombrable pour chaque $i$. Soit $M$ un objet de $D(A, \text{d})$. Alors
les énoncés suivants sont équivalents
\begin{enumerate}
\item $M = \text{hocolim} E_n$ avec $E_n$ compact dans $D(A, \text{d})$, et
\item $H^i(M)$ est dénombrable pour chaque $i$.
\end{enumerate}
\end{lemma}

\begin{proof}
Supposons que (1) soit vrai. Ensuite, nous avons $H^i(M) = \colim
H^i(E_n)$ par Catégories dérivées, lemme
\ref{derived-lemma-cohomology-of-hocolim}. Ainsi, il suffit de prouver que $H^i(E_n)$ est dénombrable pour chaque $n$.
Par la proposition \ref{proposition-compact}, nous voyons que
$E_n$ est isomorphe dans $D(A, \text{d})$ à un facteur direct d'un
module différentiel gradué $P$ qui possède une filtration finie
$F_\bullet$ par des sous-modules différentiels gradués telle
que les $F_jP/F_{j - 1}P$ soient des sommes directes finies de décalages de $A$.
Par hypothèse, les groupes $H^i(F_jP/F_{j - 1}P)$ sont dénombrables.
En raisonnant par récurrence sur la longueur de la filtration et en
utilisant la suite exacte longue de cohomologie, nous concluons que (2) est
vrai. L'implication intéressante est l'autre.

\medskip\noindent
Nous affirmons qu'il existe une sous-algèbre différentielle
graduée dénombrable $A' \subset A$ telle que l'application d'inclusion
$A' \to A$ définisse un isomorphisme sur la cohomologie. Pour
construire $A'$, nous choisissons des sous-algèbres différentielles graduées
dénombrables
$$
A_1 \subset A_2 \subset A_3 \subset \ldots
$$
de telle sorte que (a) $H^i(A_1) \to H^i(A)$ soit surjectif, et
(b) pour $n > 1$ le noyau de l'application $H^i(A_{n - 1}) \to
H^i(A_n)$ soit le même que le noyau de l'application $H^i(A_{n - 1}) \to
H^i(A)$. Pour construire $A_1$, prenons une collection dénombrable de
cochaînes $S \subset A$ engendrant la cohomologie de $A$ (comme anneau ou
comme groupe abélien gradué) et soit
$A_1$ la sous-algèbre différentielle graduée de $A$ engendrée par $S$. Pour
construire $A_n$ à partir de $A_{n - 1}$ pour chaque cochaîne $a \in A_{n - 1}^i$
qui s'envoie sur zéro dans $H^i(A)$, choisissons $s_a \in A^{i
- 1}$ avec $\text{d}(s_a) = a$ et soit $A_n$ la sous-algèbre
différentielle graduée de $A$ engendrée par $A_{n - 1}$ et les éléments $s_a$. Enfin,
prenons $A' = \bigcup A_n$.

\medskip\noindent
D'après le lemme
\ref{lemma-qis-equivalence}, l'application de restriction $D(A, \text{d}) \to D(A',
\text{d})$, $M \mapsto M_{A'}$ est une équivalence. Puisque les
groupes de cohomologie de $M$ et $M_{A'}$ sont les mêmes, nous voyons
qu'il suffit de prouver l'implication (2) $\Rightarrow$ (1)
pour $(A', \text{d})$.

\medskip\noindent
Supposons que $A$ soit dénombrable. Par le même type d'argument que celui
donné ci-dessus, nous voyons que pour $M$ dans $D(A,
\text{d})$, les énoncés suivants sont équivalents : $H^i(M)$ est dénombrable pour
chaque $i$ et $M$ peut être représenté par un module différentiel gradué dénombrable.
Par conséquent, afin de prouver l'implication (2) $\Rightarrow$ (1),
nous nous ramenons à la situation décrite dans le paragraphe suivant.

\medskip\noindent
Supposons que $A$ soit dénombrable et que $M$ soit un module différentiel gradué
dénombrable sur $A$. Nous affirmons qu'il existe un homomorphisme $P \to M$ de
$A$-modules différentiels gradués vérifiant les propriétés suivantes :
\begin{enumerate}
\item $P \to M$ est un quasi-isomorphisme,
\item $P$ possède la propriété (P), et
\item $P$ est dénombrable.
\end{enumerate}
En regardant la démonstration de la construction des P-résolutions dans
le lemme \ref{lemma-resolve}, nous voyons qu'il suffit de montrer que nous
pouvons prouver le lemme \ref{lemma-good-quotient}
dans le cadre des modules différentiels gradués dénombrables. Cela
découle immédiatement de la démonstration.

\medskip\noindent
Supposons que $A$ soit dénombrable et que $M$ soit un module
différentiel gradué dénombrable avec la propriété (P). Choisissons une filtration
$$
0 = F_{-1}P \subset F_0P \subset F_1P \subset \ldots \subset P
$$
par des sous-modules différentiels gradués, telle que l'on ait
\begin{enumerate}
\item $P = \bigcup F_pP$,
\item $F_iP \to F_{i + 1}P$ est un monomorphisme admissible,
\item des isomorphismes de modules différentiels
gradués $F_iP/F_{i - 1}P \to \bigoplus_{j \in J_i} A[k_j]$ pour
certains ensembles $J_i$ et entiers $k_j$.
\end{enumerate}
Bien sûr, $J_i$ est dénombrable pour chaque $i$. Pour chaque $i$
et $j \in J_i$, choisissons $x_{i, j} \in F_iP$ de degré $k_j$ dont
l'image dans $F_iP/F_{i - 1}P$ engendre le facteur direct correspondant
à $j$.

\medskip\noindent
Affirmation : Étant donné $n$ et des sous-ensembles finis $S_i \subset J_i$, $i = 1,
\ldots, n$, il existe des sous-ensembles finis $S_i \subset T_i \subset J_i$, $i = 1,
\ldots, n$ tels que $P' = \bigoplus_{i \leq n} \bigoplus_{j \in T_i} Ax_{i, j}$ soit un
sous-module différentiel gradué de $P$. Cela a été montré dans la
démonstration du lemme \ref{lemma-factor-through-nicer} mais cela se montre aussi
facilement directement : les éléments $x_{i, j}$ engendrent librement $P$ en
tant que module à droite sur $A$. La structure de $P$ montre que
$$
\text{d}(x_{i, j}) = \sum\nolimits_{i' < i} x_{i', j'}a_{i', j'}
$$
où bien sûr la somme est finie.
Ainsi, étant donnés $S_0, \ldots, S_n$, nous pouvons d'abord
choisir $S_0 \subset S'_0, \ldots, S_{n - 1} \subset S'_{n - 1}$
avec $\text{d}(x_{n, j}) \in \bigoplus_{i' < n, j' \in S'_{i'}} x_{i', j'}A$ pour
tous $j \in S_n$. Ensuite, par récurrence sur $n$, nous pouvons
choisir $S'_0 \subset T_0, \ldots, S'_{n - 1} \subset T_{n - 1}$
pour s'assurer que $\bigoplus_{i' < n, j' \in T_{i'}} x_{i', j'}A$ soit
un $A$-sous-module différentiel gradué. En définissant $T_n = S_n$, nous trouvons
que $P' = \bigoplus_{i \leq n, j \in T_i} x_{i, j}A$ convient.

\medskip\noindent
D'après l'assertion, il est clair que $P = \bigcup P'_n$
est une union croissante dénombrable de $P'_n$ comme
indiqué ci-dessus. Par construction, chaque $P'_n$ est un module différentiel
gradué avec la propriété (P) tel que la filtration soit finie et que les
quotients successifs sont des sommes directes finies de décalages de $A$.
Ainsi, $P'_n$ définit un objet compact de $D(A, \text{d})$, voir par
exemple la proposition \ref{proposition-compact}.
Puisque $P = \text{hocolim} P'_n$ dans $D(A, \text{d})$ par
le lemme \ref{lemma-homotopy-colimit}, la
démonstration de l'implication (2) $\Rightarrow$ (1) est complète.
\end{proof}




\begin{multicols}{2}[\section{Autres chapitres}]
\noindent
Préliminaires
\begin{enumerate}
\item \hyperref[introduction-section-phantom]{Introduction}
\item \hyperref[conventions-section-phantom]{Conventions}
\item \hyperref[sets-section-phantom]{Théorie des ensembles}
\item \hyperref[categories-section-phantom]{Catégories}
\item \hyperref[topology-section-phantom]{Topologie}
\item \hyperref[sheaves-section-phantom]{Faisceaux sur les espaces}
\item \hyperref[sites-section-phantom]{Sites et faisceaux}
\item \hyperref[stacks-section-phantom]{Champs}
\item \hyperref[fields-section-phantom]{Corps}
\item \hyperref[algebra-section-phantom]{Algèbre commutative}
\item \hyperref[brauer-section-phantom]{Groupes de Brauer}
\item \hyperref[homology-section-phantom]{Algèbre homologique}
\item \hyperref[derived-section-phantom]{Catégories dérivées}
\item \hyperref[simplicial-section-phantom]{Méthodes simpliciales}
\item \hyperref[more-algebra-section-phantom]{Compléments d'algèbre}
\item \hyperref[smoothing-section-phantom]{Lissage des morphismes d'anneaux}
\item \hyperref[modules-section-phantom]{Faisceaux de modules}
\item \hyperref[sites-modules-section-phantom]{Modules sur les sites}
\item \hyperref[injectives-section-phantom]{Objets injectifs}
\item \hyperref[cohomology-section-phantom]{Cohomologie des faisceaux}
\item \hyperref[sites-cohomology-section-phantom]{Cohomologie sur les sites}
\item \hyperref[dga-section-phantom]{Algèbre différentielle graduée}
\item \hyperref[dpa-section-phantom]{Algèbre à puissances divisées}
\item \hyperref[sdga-section-phantom]{Faisceaux différentiels gradués}
\item \hyperref[hypercovering-section-phantom]{Hyperrecouvrements}
\end{enumerate}
Schémas
\begin{enumerate}
\setcounter{enumi}{25}
\item \hyperref[schemes-section-phantom]{Schémas}
\item \hyperref[constructions-section-phantom]{Constructions de schémas}
\item \hyperref[properties-section-phantom]{Propriétés des schémas}
\item \hyperref[morphisms-section-phantom]{Morphismes de schémas}
\item \hyperref[coherent-section-phantom]{Cohomologie des schémas}
\item \hyperref[divisors-section-phantom]{Diviseurs}
\item \hyperref[limits-section-phantom]{Limites de schémas}
\item \hyperref[varieties-section-phantom]{Variétés}
\item \hyperref[topologies-section-phantom]{Topologies sur les schémas}
\item \hyperref[descent-section-phantom]{Descente}
\item \hyperref[perfect-section-phantom]{Catégories dérivées de schémas}
\item \hyperref[more-morphisms-section-phantom]{Compléments sur les morphismes}
\item \hyperref[flat-section-phantom]{Compléments sur la platitude}
\item \hyperref[groupoids-section-phantom]{Schémas en groupoïdes}
\item \hyperref[more-groupoids-section-phantom]{Compléments sur les schémas en groupoïdes}
\item \hyperref[etale-section-phantom]{Morphismes étales de schémas}
\end{enumerate}
Thèmes de théorie des schémas
\begin{enumerate}
\setcounter{enumi}{41}
\item \hyperref[chow-section-phantom]{Homologie de Chow}
\item \hyperref[intersection-section-phantom]{Théorie de l'intersection}
\item \hyperref[pic-section-phantom]{Schémas de Picard des courbes}
\item \hyperref[weil-section-phantom]{Théories cohomologiques de Weil}
\item \hyperref[adequate-section-phantom]{Modules adéquats}
\item \hyperref[dualizing-section-phantom]{Complexes dualisants}
\item \hyperref[duality-section-phantom]{Dualité pour les schémas}
\item \hyperref[discriminant-section-phantom]{Discriminants et différentes}
\item \hyperref[derham-section-phantom]{Cohomologie de de Rham}
\item \hyperref[local-cohomology-section-phantom]{Cohomologie locale}
\item \hyperref[algebraization-section-phantom]{Géométrie algébrique et formelle}
\item \hyperref[curves-section-phantom]{Courbes algébriques}
\item \hyperref[resolve-section-phantom]{Résolution des surfaces}
\item \hyperref[models-section-phantom]{Réduction semi-stable}
\item \hyperref[functors-section-phantom]{Foncteurs et morphismes}
\item \hyperref[equiv-section-phantom]{Catégories dérivées de variétés}
\item \hyperref[pione-section-phantom]{Groupes fondamentaux de schémas}
\item \hyperref[etale-cohomology-section-phantom]{Cohomologie étale}
\item \hyperref[crystalline-section-phantom]{Cohomologie cristalline}
\item \hyperref[proetale-section-phantom]{Cohomologie pro-étale}
\item \hyperref[relative-cycles-section-phantom]{Cycles relatifs}
\item \hyperref[more-etale-section-phantom]{Compléments de cohomologie étale}
\item \hyperref[trace-section-phantom]{La formule des traces}
\end{enumerate}
Espaces algébriques
\begin{enumerate}
\setcounter{enumi}{64}
\item \hyperref[spaces-section-phantom]{Espaces algébriques}
\item \hyperref[spaces-properties-section-phantom]{Propriétés des espaces algébriques}
\item \hyperref[spaces-morphisms-section-phantom]{Morphismes d'espaces algébriques}
\item \hyperref[decent-spaces-section-phantom]{Espaces algébriques décents}
\item \hyperref[spaces-cohomology-section-phantom]{Cohomologie des espaces algébriques}
\item \hyperref[spaces-limits-section-phantom]{Limites d'espaces algébriques}
\item \hyperref[spaces-divisors-section-phantom]{Diviseurs sur les espaces algébriques}
\item \hyperref[spaces-over-fields-section-phantom]{Espaces algébriques sur des corps}
\item \hyperref[spaces-topologies-section-phantom]{Topologies sur les espaces algébriques}
\item \hyperref[spaces-descent-section-phantom]{Descente et espaces algébriques}
\item \hyperref[spaces-perfect-section-phantom]{Catégories dérivées d'espaces}
\item \hyperref[spaces-more-morphisms-section-phantom]{Compléments sur les morphismes d'espaces}
\item \hyperref[spaces-flat-section-phantom]{Platitude sur les espaces algébriques}
\item \hyperref[spaces-groupoids-section-phantom]{Groupoïdes dans les espaces algébriques}
\item \hyperref[spaces-more-groupoids-section-phantom]{Compléments sur les groupoïdes dans les espaces}
\item \hyperref[bootstrap-section-phantom]{Amorçage}
\item \hyperref[spaces-pushouts-section-phantom]{Sommes amalgamées d'espaces algébriques}
\end{enumerate}
Thèmes de géométrie
\begin{enumerate}
\setcounter{enumi}{81}
\item \hyperref[spaces-chow-section-phantom]{Groupes de Chow des espaces}
\item \hyperref[groupoids-quotients-section-phantom]{Quotients de groupoïdes}
\item \hyperref[spaces-more-cohomology-section-phantom]{Compléments sur la cohomologie des espaces}
\item \hyperref[spaces-simplicial-section-phantom]{Espaces simpliciaux}
\item \hyperref[spaces-duality-section-phantom]{Dualité pour les espaces}
\item \hyperref[formal-spaces-section-phantom]{Espaces algébriques formels}
\item \hyperref[restricted-section-phantom]{Algébrisation des espaces formels}
\item \hyperref[spaces-resolve-section-phantom]{Résolution des surfaces revisitée}
\end{enumerate}
Théorie des déformations
\begin{enumerate}
\setcounter{enumi}{89}
\item \hyperref[formal-defos-section-phantom]{Théorie formelle des déformations}
\item \hyperref[defos-section-phantom]{Théorie des déformations}
\item \hyperref[cotangent-section-phantom]{Le complexe cotangent}
\item \hyperref[examples-defos-section-phantom]{Problèmes de déformation}
\end{enumerate}
Champs algébriques
\begin{enumerate}
\setcounter{enumi}{93}
\item \hyperref[algebraic-section-phantom]{Champs algébriques}
\item \hyperref[examples-stacks-section-phantom]{Exemples de champs}
\item \hyperref[stacks-sheaves-section-phantom]{Faisceaux sur les champs algébriques}
\item \hyperref[criteria-section-phantom]{Critères de représentabilité}
\item \hyperref[artin-section-phantom]{Axiomes d'Artin}
\item \hyperref[quot-section-phantom]{Espaces de Quot et de Hilbert}
\item \hyperref[stacks-properties-section-phantom]{Propriétés des champs algébriques}
\item \hyperref[stacks-morphisms-section-phantom]{Morphismes de champs algébriques}
\item \hyperref[stacks-limits-section-phantom]{Limites de champs algébriques}
\item \hyperref[stacks-cohomology-section-phantom]{Cohomologie des champs algébriques}
\item \hyperref[stacks-perfect-section-phantom]{Catégories dérivées de champs}
\item \hyperref[stacks-introduction-section-phantom]{Introduction aux champs algébriques}
\item \hyperref[stacks-more-morphisms-section-phantom]{Compléments sur les morphismes de champs}
\item \hyperref[stacks-geometry-section-phantom]{La géométrie des champs}
\end{enumerate}
Thèmes de théorie des espaces de modules
\begin{enumerate}
\setcounter{enumi}{107}
\item \hyperref[moduli-section-phantom]{Champs de modules}
\item \hyperref[moduli-curves-section-phantom]{Espaces de modules de courbes}
\end{enumerate}
Divers
\begin{enumerate}
\setcounter{enumi}{109}
\item \hyperref[examples-section-phantom]{Exemples}
\item \hyperref[exercises-section-phantom]{Exercices}
\item \hyperref[guide-section-phantom]{Guide bibliographique}
\item \hyperref[desirables-section-phantom]{Desiderata}
\item \hyperref[coding-section-phantom]{Style de programmation}
\item \hyperref[obsolete-section-phantom]{Obsolète}
\item \hyperref[fdl-section-phantom]{Licence de documentation libre GNU}
\item \hyperref[index-section-phantom]{Index généré automatiquement}
\end{enumerate}
\end{multicols}


\bibliography{my}
\bibliographystyle{amsalpha}

\end{document}
